% Torus knot, closed and open: two phasors into one table, at 2 : 3 and at the golden ratio; circuit, loop, sound.
% Closed: the path is the (2, 3) torus knot (a trefoil), and every line sits on one harmonic series. Open: it fills the surface.
% Geometry and data: torus-knot-closed-and-open.py (tkco-mesh.tex, tkco-knot.tex, tkco-fill.tex, tkco-closed.dat, tkco-open.dat).
% Everything here is computed (ink). Compile: lualatex torus-knot-closed-and-open.tex ; SVG via PyMuPDF.
\documentclass[border=0pt]{standalone}
\usepackage{fontspec}
\setmainfont{JetBrains Mono}
\setmonofont{JetBrains Mono}
\usepackage{tikz}
\usepackage{pgfplots}
\pgfplotsset{compat=1.16}
\usetikzlibrary{arrows.meta, calc, backgrounds, decorations.markings, positioning}
% Drafting Table, paper ground (_ops/drafting-table/design-system/colors_and_type.css)
\definecolor{paper}{HTML}{F0EEE5}
\definecolor{ink}{HTML}{1A1916}
\definecolor{ink2}{HTML}{3A3A35}
\definecolor{ink3}{HTML}{6A6862}
\definecolor{slab}{HTML}{D8D4C6}
\definecolor{measure}{HTML}{C8261A}
\tikzset{
  >={Stealth[length=1.8mm]},
  mesh/.style={draw=ink3!45, line width=0.25pt},
  knot/.style={draw=ink, line width=1.5pt, line cap=round},
  knotback/.style={draw=ink, line width=0.7pt, densely dashed},
  fill/.style={draw=ink, line width=0.45pt},
  fillback/.style={draw=ink3!60, line width=0.3pt, densely dashed},
  loop/.style={draw=ink, line width=1.5pt},
  back/.style={draw=ink, line width=1.5pt, ->},
  wire/.style={draw=ink, line width=0.6pt},
  blk/.style={rectangle, draw=ink, fill=slab, line width=0.7pt, minimum height=7mm, minimum width=24mm, inner sep=2pt, font=\scriptsize, text=ink},
  lab/.style={font=\scriptsize, text=ink2},
  rowhead/.style={font=\bfseries\small, anchor=west, text=ink},
  rowsub/.style={font=\scriptsize, anchor=west, text=ink2},
  colhead/.style={font=\bfseries\scriptsize, anchor=west, text=ink},
  spin/.style={postaction={decorate}, decoration={markings, mark=at position 0.25 with {\arrow{>}}}},
}
\pgfplotsset{
  /pgf/number format/assume math mode=true,
  snd/.style={width=6.6cm, height=3.0cm, axis lines=left, tick align=outside, ymin=0, ytick=\empty,
    axis line style={ink, line width=0.5pt}, tick style={ink}, clip=false,
    xlabel style={font=\scriptsize, text=ink2, yshift=1mm}, ylabel style={font=\scriptsize, text=ink2},
    tick label style={font=\tiny, text=ink2}},
}
\newcommand{\rowframe}[3]{%
  \draw[ink, line width=0.5pt] (-0.3,#1) rectangle (19.9,#1-5.9);
  \draw[ink, line width=0.3pt] (6.9,#1-0.95) -- (6.9,#1-5.9) (12.4,#1-0.95) -- (12.4,#1-5.9) (-0.3,#1-0.95) -- (19.9,#1-0.95);
  \node[rowhead] at (0,#1-0.32) {#2};
  \node[rowsub] at (0,#1-0.7) {#3};}
% the circuit: two phasors into one table (#1, #2: the phasor labels)
\newcommand{\twophasors}[2]{%
  \node[blk] (p1) at (1.4,-2.5) {#1};
  \node[blk] (p2) at (1.4,-4.85) {#2};
  \draw[back] (p1.north east) to[out=40, in=140, looseness=1.6] node[lab, above] {wraps} (p1.north west);
  \draw[back] (p2.north east) to[out=40, in=140, looseness=1.6] node[lab, above] {wraps} (p2.north west);
  \node[blk, minimum width=20mm, minimum height=12mm] (w) at (4.6,-3.7) {table W(φ₁, φ₂)};
  \draw[wire,->] (p1.east) -| ([xshift=-4mm]w.north);
  \draw[wire,->] (p2.east) -| ([xshift=-4mm]w.south);
  \draw[wire,->] (w.east) -- (6.65,-3.7) node[lab, above left] {out};}
% the loop: two circles, read together, are a torus
\newcommand{\twocircles}{%
  \draw[loop, spin] (7.75,-2.05) circle (0.42); \node[lab] at (7.75,-2.05) {φ₁};
  \draw[loop, spin] (7.75,-4.2) circle (0.3);  \node[lab] at (7.75,-4.2) {φ₂};
  \node[font=\small, text=ink] at (8.35,-3.15) {×};}
% the sound: lines over the dashed series of f0 = f1 / 2
\newcommand{\grid}{\foreach \k in {1,...,11} { \edef\t{\noexpand\draw[ink3!60, densely dashed, line width=0.4pt] (axis cs:\k/2,0) -- (axis cs:\k/2,1.15);} \t }}
\begin{document}
\begin{tikzpicture}

% ---------- column heads ----------
\node[colhead] at (0,1.15) {THE CIRCUIT \quad what you build};
\node[colhead] at (7.2,1.15) {THE LOOP \quad what it is};
\node[colhead] at (12.7,1.15) {THE SOUND \quad what you hear};
\node[font=\scriptsize, anchor=west, text=ink2] at (0,0.55) {Same circuit twice; only the ratio of the two phasors changes. Heavy: the path that comes back to where it began.};

% =============== CLOSED · 2 : 3 ===============
\def\y{0}
\rowframe{\y}{CLOSED \enspace f₁ : f₂ = 2 : 3}{the path comes home: twice round one loop, three times round the other}
\begin{scope}[shift={(0,\y)}]
  \twophasors{phasor f₁ = 2f₀}{phasor f₂ = 3f₀}
  \twocircles
  \begin{scope}[shift={(10.45,-3.25)}, scale=0.75]
    % torus mesh, generated by torus-knot-closed-and-open.py
\draw[mesh] (2.428,0.028) -- (2.422,0.056) -- (2.411,0.084) -- (2.397,0.111) -- (2.378,0.138) -- (2.356,0.164) -- (2.330,0.188) -- (2.300,0.212) -- (2.267,0.234) -- (2.231,0.255) -- (2.192,0.274) -- (2.150,0.292) -- (2.105,0.307) -- (2.059,0.321) -- (2.010,0.333) -- (1.960,0.343) -- (1.909,0.350) -- (1.856,0.356) -- (1.803,0.359) -- (1.750,0.360) -- (1.697,0.359) -- (1.644,0.356) -- (1.591,0.350) -- (1.540,0.343) -- (1.490,0.333) -- (1.441,0.321) -- (1.395,0.307) -- (1.350,0.292) -- (1.308,0.274) -- (1.269,0.255) -- (1.233,0.234) -- (1.200,0.212) -- (1.170,0.188) -- (1.144,0.164) -- (1.122,0.138) -- (1.103,0.111) -- (1.089,0.084) -- (1.078,0.056) -- (1.072,0.028) -- (1.070,0.000) -- (1.072,-0.028);
\draw[mesh] (2.329,0.613) -- (2.315,0.637) -- (2.297,0.660) -- (2.276,0.681) -- (2.250,0.700) -- (2.222,0.717) -- (2.190,0.732) -- (2.155,0.744) -- (2.117,0.755) -- (2.076,0.763) -- (2.034,0.769) -- (1.989,0.773) -- (1.942,0.774) -- (1.893,0.773) -- (1.844,0.769) -- (1.793,0.763) -- (1.742,0.755) -- (1.690,0.744) -- (1.639,0.732) -- (1.588,0.717) -- (1.537,0.700) -- (1.487,0.681) -- (1.439,0.660) -- (1.392,0.637) -- (1.347,0.613) -- (1.304,0.588) -- (1.264,0.561) -- (1.226,0.533) -- (1.191,0.505) -- (1.159,0.475) -- (1.130,0.445) -- (1.105,0.415) -- (1.084,0.384) -- (1.066,0.354) -- (1.052,0.323) -- (1.042,0.293) -- (1.036,0.264) -- (1.034,0.235) -- (1.036,0.207) -- (1.042,0.180) -- (1.052,0.155);
\draw[mesh] (2.076,1.128) -- (2.060,1.146) -- (2.040,1.163) -- (2.018,1.176) -- (1.992,1.187) -- (1.963,1.195) -- (1.932,1.201) -- (1.898,1.203) -- (1.862,1.203) -- (1.823,1.200) -- (1.783,1.194) -- (1.741,1.185) -- (1.698,1.174) -- (1.653,1.160) -- (1.608,1.143) -- (1.562,1.124) -- (1.516,1.102) -- (1.469,1.079) -- (1.423,1.053) -- (1.378,1.025) -- (1.334,0.996) -- (1.290,0.965) -- (1.248,0.932) -- (1.208,0.899) -- (1.169,0.864) -- (1.133,0.829) -- (1.099,0.793) -- (1.068,0.757) -- (1.039,0.721) -- (1.013,0.684) -- (0.991,0.649) -- (0.971,0.614) -- (0.955,0.579) -- (0.943,0.546) -- (0.934,0.514) -- (0.928,0.483) -- (0.927,0.454) -- (0.928,0.426) -- (0.934,0.401) -- (0.943,0.378) -- (0.955,0.356);
\draw[mesh] (1.666,1.576) -- (1.647,1.585) -- (1.626,1.591) -- (1.603,1.594) -- (1.577,1.593) -- (1.550,1.588) -- (1.520,1.581) -- (1.489,1.570) -- (1.456,1.556) -- (1.421,1.538) -- (1.386,1.518) -- (1.350,1.495) -- (1.313,1.469) -- (1.275,1.441) -- (1.237,1.410) -- (1.200,1.377) -- (1.162,1.342) -- (1.125,1.305) -- (1.089,1.266) -- (1.053,1.226) -- (1.019,1.185) -- (0.986,1.144) -- (0.955,1.101) -- (0.925,1.059) -- (0.897,1.016) -- (0.872,0.973) -- (0.848,0.931) -- (0.827,0.890) -- (0.809,0.850) -- (0.793,0.811) -- (0.780,0.773) -- (0.770,0.737) -- (0.763,0.703) -- (0.758,0.671) -- (0.757,0.642) -- (0.758,0.615) -- (0.763,0.590) -- (0.770,0.569) -- (0.780,0.550) -- (0.793,0.535) -- (0.809,0.522);
\draw[mesh] (1.165,1.899) -- (1.150,1.901) -- (1.134,1.899) -- (1.115,1.893) -- (1.096,1.884) -- (1.075,1.870) -- (1.053,1.853) -- (1.029,1.833) -- (1.005,1.809) -- (0.980,1.782) -- (0.954,1.752) -- (0.928,1.719) -- (0.902,1.684) -- (0.875,1.646) -- (0.848,1.605) -- (0.822,1.563) -- (0.796,1.519) -- (0.770,1.474) -- (0.745,1.427) -- (0.721,1.380) -- (0.697,1.332) -- (0.675,1.283) -- (0.654,1.235) -- (0.635,1.187) -- (0.616,1.140) -- (0.600,1.093) -- (0.585,1.048) -- (0.572,1.004) -- (0.561,0.962) -- (0.552,0.922) -- (0.544,0.884) -- (0.539,0.848) -- (0.536,0.816) -- (0.535,0.786) -- (0.536,0.759) -- (0.539,0.736) -- (0.544,0.716) -- (0.552,0.699) -- (0.561,0.686) -- (0.572,0.677) -- (0.585,0.671);
\draw[mesh] (0.603,2.097) -- (0.595,2.096) -- (0.587,2.091) -- (0.577,2.082) -- (0.567,2.069) -- (0.556,2.052) -- (0.545,2.032) -- (0.533,2.007) -- (0.520,1.980) -- (0.507,1.948) -- (0.494,1.914) -- (0.480,1.877) -- (0.467,1.836) -- (0.453,1.794) -- (0.439,1.749) -- (0.425,1.702) -- (0.412,1.654) -- (0.399,1.604) -- (0.386,1.553) -- (0.373,1.502) -- (0.361,1.450) -- (0.349,1.398) -- (0.339,1.346) -- (0.328,1.294) -- (0.319,1.244) -- (0.311,1.195) -- (0.303,1.147) -- (0.296,1.101) -- (0.290,1.057) -- (0.286,1.015) -- (0.282,0.976) -- (0.279,0.940) -- (0.277,0.906) -- (0.277,0.876) -- (0.277,0.850) -- (0.279,0.827) -- (0.282,0.808) -- (0.286,0.792) -- (0.290,0.781) -- (0.296,0.774) -- (0.303,0.770);
\draw[mesh] (0.000,2.162) -- (0.000,2.157) -- (0.000,2.147) -- (0.000,2.133) -- (0.000,2.115) -- (0.000,2.093) -- (0.000,2.067) -- (0.000,2.038) -- (0.000,2.005) -- (0.000,1.969) -- (0.000,1.930) -- (0.000,1.889) -- (0.000,1.844) -- (0.000,1.798) -- (0.000,1.750) -- (0.000,1.700) -- (0.000,1.649) -- (0.000,1.596) -- (0.000,1.543) -- (0.000,1.490) -- (0.000,1.437) -- (0.000,1.384) -- (0.000,1.331) -- (0.000,1.280) -- (0.000,1.229) -- (0.000,1.181) -- (0.000,1.134) -- (0.000,1.089) -- (0.000,1.047) -- (0.000,1.007) -- (0.000,0.971) -- (0.000,0.937) -- (0.000,0.907) -- (0.000,0.881) -- (0.000,0.858) -- (0.000,0.839) -- (0.000,0.824) -- (0.000,0.813) -- (0.000,0.807) -- (0.000,0.804) -- (0.000,0.806);
\draw[mesh] (-0.603,2.097) -- (-0.595,2.096) -- (-0.587,2.091) -- (-0.577,2.082) -- (-0.567,2.069) -- (-0.556,2.052) -- (-0.545,2.032) -- (-0.533,2.007) -- (-0.520,1.980) -- (-0.507,1.948) -- (-0.494,1.914) -- (-0.480,1.877) -- (-0.467,1.836) -- (-0.453,1.794) -- (-0.439,1.749) -- (-0.425,1.702) -- (-0.412,1.654) -- (-0.399,1.604) -- (-0.386,1.553) -- (-0.373,1.502) -- (-0.361,1.450) -- (-0.349,1.398) -- (-0.339,1.346) -- (-0.328,1.294) -- (-0.319,1.244) -- (-0.311,1.195) -- (-0.303,1.147) -- (-0.296,1.101) -- (-0.290,1.057) -- (-0.286,1.015) -- (-0.282,0.976) -- (-0.279,0.940) -- (-0.277,0.906) -- (-0.277,0.876) -- (-0.277,0.850) -- (-0.279,0.827) -- (-0.282,0.808) -- (-0.286,0.792) -- (-0.290,0.781) -- (-0.296,0.774) -- (-0.303,0.770);
\draw[mesh] (-1.165,1.899) -- (-1.150,1.901) -- (-1.134,1.899) -- (-1.115,1.893) -- (-1.096,1.884) -- (-1.075,1.870) -- (-1.053,1.853) -- (-1.029,1.833) -- (-1.005,1.809) -- (-0.980,1.782) -- (-0.954,1.752) -- (-0.928,1.719) -- (-0.902,1.684) -- (-0.875,1.646) -- (-0.848,1.605) -- (-0.822,1.563) -- (-0.796,1.519) -- (-0.770,1.474) -- (-0.745,1.427) -- (-0.721,1.380) -- (-0.697,1.332) -- (-0.675,1.283) -- (-0.654,1.235) -- (-0.635,1.187) -- (-0.616,1.140) -- (-0.600,1.093) -- (-0.585,1.048) -- (-0.572,1.004) -- (-0.561,0.962) -- (-0.552,0.922) -- (-0.544,0.884) -- (-0.539,0.848) -- (-0.536,0.816) -- (-0.535,0.786) -- (-0.536,0.759) -- (-0.539,0.736) -- (-0.544,0.716) -- (-0.552,0.699) -- (-0.561,0.686) -- (-0.572,0.677) -- (-0.585,0.671);
\draw[mesh] (-1.666,1.576) -- (-1.647,1.585) -- (-1.626,1.591) -- (-1.603,1.594) -- (-1.577,1.593) -- (-1.550,1.588) -- (-1.520,1.581) -- (-1.489,1.570) -- (-1.456,1.556) -- (-1.421,1.538) -- (-1.386,1.518) -- (-1.350,1.495) -- (-1.313,1.469) -- (-1.275,1.441) -- (-1.237,1.410) -- (-1.200,1.377) -- (-1.162,1.342) -- (-1.125,1.305) -- (-1.089,1.266) -- (-1.053,1.226) -- (-1.019,1.185) -- (-0.986,1.144) -- (-0.955,1.101) -- (-0.925,1.059) -- (-0.897,1.016) -- (-0.872,0.973) -- (-0.848,0.931) -- (-0.827,0.890) -- (-0.809,0.850) -- (-0.793,0.811) -- (-0.780,0.773) -- (-0.770,0.737) -- (-0.763,0.703) -- (-0.758,0.671) -- (-0.757,0.642) -- (-0.758,0.615) -- (-0.763,0.590) -- (-0.770,0.569) -- (-0.780,0.550) -- (-0.793,0.535) -- (-0.809,0.522);
\draw[mesh] (-2.076,1.128) -- (-2.060,1.146) -- (-2.040,1.163) -- (-2.018,1.176) -- (-1.992,1.187) -- (-1.963,1.195) -- (-1.932,1.201) -- (-1.898,1.203) -- (-1.862,1.203) -- (-1.823,1.200) -- (-1.783,1.194) -- (-1.741,1.185) -- (-1.698,1.174) -- (-1.653,1.160) -- (-1.608,1.143) -- (-1.562,1.124) -- (-1.516,1.102) -- (-1.469,1.079) -- (-1.423,1.053) -- (-1.378,1.025) -- (-1.334,0.996) -- (-1.290,0.965) -- (-1.248,0.932) -- (-1.208,0.899) -- (-1.169,0.864) -- (-1.133,0.829) -- (-1.099,0.793) -- (-1.068,0.757) -- (-1.039,0.721) -- (-1.013,0.684) -- (-0.991,0.649) -- (-0.971,0.614) -- (-0.955,0.579) -- (-0.943,0.546) -- (-0.934,0.514) -- (-0.928,0.483) -- (-0.927,0.454) -- (-0.928,0.426) -- (-0.934,0.401) -- (-0.943,0.378) -- (-0.955,0.356);
\draw[mesh] (-2.329,0.613) -- (-2.315,0.637) -- (-2.297,0.660) -- (-2.276,0.681) -- (-2.250,0.700) -- (-2.222,0.717) -- (-2.190,0.732) -- (-2.155,0.744) -- (-2.117,0.755) -- (-2.076,0.763) -- (-2.034,0.769) -- (-1.989,0.773) -- (-1.942,0.774) -- (-1.893,0.773) -- (-1.844,0.769) -- (-1.793,0.763) -- (-1.742,0.755) -- (-1.690,0.744) -- (-1.639,0.732) -- (-1.588,0.717) -- (-1.537,0.700) -- (-1.487,0.681) -- (-1.439,0.660) -- (-1.392,0.637) -- (-1.347,0.613) -- (-1.304,0.588) -- (-1.264,0.561) -- (-1.226,0.533) -- (-1.191,0.505) -- (-1.159,0.475) -- (-1.130,0.445) -- (-1.105,0.415) -- (-1.084,0.384) -- (-1.066,0.354) -- (-1.052,0.323) -- (-1.042,0.293) -- (-1.036,0.264) -- (-1.034,0.235) -- (-1.036,0.207) -- (-1.042,0.180) -- (-1.052,0.155);
\draw[mesh] (-2.428,0.028) -- (-2.422,0.056) -- (-2.411,0.084) -- (-2.397,0.111) -- (-2.378,0.138) -- (-2.356,0.164) -- (-2.330,0.188) -- (-2.300,0.212) -- (-2.267,0.234) -- (-2.231,0.255) -- (-2.192,0.274) -- (-2.150,0.292) -- (-2.105,0.307) -- (-2.059,0.321) -- (-2.010,0.333) -- (-1.960,0.343) -- (-1.909,0.350) -- (-1.856,0.356) -- (-1.803,0.359) -- (-1.750,0.360) -- (-1.697,0.359) -- (-1.644,0.356) -- (-1.591,0.350) -- (-1.540,0.343) -- (-1.490,0.333) -- (-1.441,0.321) -- (-1.395,0.307) -- (-1.350,0.292) -- (-1.308,0.274) -- (-1.269,0.255) -- (-1.233,0.234) -- (-1.200,0.212) -- (-1.170,0.188) -- (-1.144,0.164) -- (-1.122,0.138) -- (-1.103,0.111) -- (-1.089,0.084) -- (-1.078,0.056) -- (-1.072,0.028) -- (-1.070,0.000) -- (-1.072,-0.028);
\draw[mesh] (-2.347,-0.533) -- (-2.345,-0.505) -- (-2.339,-0.475) -- (-2.329,-0.445) -- (-2.315,-0.415) -- (-2.297,-0.384) -- (-2.276,-0.354) -- (-2.250,-0.323) -- (-2.222,-0.293) -- (-2.190,-0.264) -- (-2.155,-0.235) -- (-2.117,-0.207) -- (-2.076,-0.180) -- (-2.034,-0.155) -- (-1.989,-0.131) -- (-1.942,-0.108) -- (-1.893,-0.088) -- (-1.844,-0.069) -- (-1.793,-0.052) -- (-1.742,-0.037) -- (-1.690,-0.024) -- (-1.639,-0.013) -- (-1.588,-0.005) -- (-1.537,0.001) -- (-1.487,0.005) -- (-1.439,0.006) -- (-1.392,0.005) -- (-1.347,0.001) -- (-1.304,-0.005) -- (-1.264,-0.013) -- (-1.226,-0.024) -- (-1.191,-0.037) -- (-1.159,-0.052) -- (-1.130,-0.069) -- (-1.105,-0.088) -- (-1.084,-0.108) -- (-1.066,-0.131) -- (-1.052,-0.155) -- (-1.042,-0.180);
\draw[mesh] (-2.339,-0.588) -- (-2.345,-0.561) -- (-2.347,-0.533);
\draw[mesh] (-2.104,-1.030) -- (-2.103,-1.001) -- (-2.097,-0.970) -- (-2.088,-0.938) -- (-2.076,-0.905) -- (-2.060,-0.871) -- (-2.040,-0.835) -- (-2.018,-0.800) -- (-1.992,-0.764) -- (-1.963,-0.727) -- (-1.932,-0.691) -- (-1.898,-0.655) -- (-1.862,-0.620) -- (-1.823,-0.585) -- (-1.783,-0.552) -- (-1.741,-0.519) -- (-1.698,-0.488) -- (-1.653,-0.459) -- (-1.608,-0.431) -- (-1.562,-0.405) -- (-1.516,-0.382) -- (-1.469,-0.360) -- (-1.423,-0.341) -- (-1.378,-0.324) -- (-1.334,-0.310) -- (-1.290,-0.299) -- (-1.248,-0.290) -- (-1.208,-0.284) -- (-1.169,-0.281) -- (-1.133,-0.281) -- (-1.099,-0.283) -- (-1.068,-0.289) -- (-1.039,-0.297) -- (-1.013,-0.308) -- (-0.991,-0.322) -- (-0.971,-0.338) -- (-0.955,-0.356) -- (-0.943,-0.378);
\draw[mesh] (-2.088,-1.107) -- (-2.097,-1.083) -- (-2.103,-1.058) -- (-2.104,-1.030);
\draw[mesh] (-1.718,-1.457) -- (-1.717,-1.428) -- (-1.712,-1.396) -- (-1.705,-1.362) -- (-1.695,-1.326) -- (-1.682,-1.288) -- (-1.666,-1.249) -- (-1.647,-1.209) -- (-1.626,-1.167) -- (-1.603,-1.125) -- (-1.577,-1.083) -- (-1.550,-1.040) -- (-1.520,-0.998) -- (-1.489,-0.955) -- (-1.456,-0.913) -- (-1.421,-0.873) -- (-1.386,-0.833) -- (-1.350,-0.794) -- (-1.313,-0.757) -- (-1.275,-0.722) -- (-1.237,-0.689) -- (-1.200,-0.658) -- (-1.162,-0.630) -- (-1.125,-0.604) -- (-1.089,-0.581) -- (-1.053,-0.560) -- (-1.019,-0.543) -- (-0.986,-0.529) -- (-0.955,-0.518) -- (-0.925,-0.511) -- (-0.897,-0.506) -- (-0.872,-0.505) -- (-0.848,-0.508) -- (-0.827,-0.513) -- (-0.809,-0.522) -- (-0.793,-0.535);
\draw[mesh] (-1.682,-1.564) -- (-1.695,-1.549) -- (-1.705,-1.530) -- (-1.712,-1.509) -- (-1.717,-1.484) -- (-1.718,-1.457);
\draw[mesh] (-1.215,-1.785) -- (-1.214,-1.755) -- (-1.211,-1.722) -- (-1.206,-1.687) -- (-1.198,-1.649) -- (-1.189,-1.609) -- (-1.178,-1.567) -- (-1.165,-1.523) -- (-1.150,-1.477) -- (-1.134,-1.431) -- (-1.115,-1.384) -- (-1.096,-1.336) -- (-1.075,-1.287) -- (-1.053,-1.239) -- (-1.029,-1.191) -- (-1.005,-1.143) -- (-0.980,-1.097) -- (-0.954,-1.051) -- (-0.928,-1.007) -- (-0.902,-0.965) -- (-0.875,-0.925) -- (-0.848,-0.887) -- (-0.822,-0.851) -- (-0.796,-0.818) -- (-0.770,-0.788) -- (-0.745,-0.761) -- (-0.721,-0.737) -- (-0.697,-0.717) -- (-0.675,-0.700) -- (-0.654,-0.687) -- (-0.635,-0.677) -- (-0.616,-0.671) -- (-0.600,-0.669) -- (-0.585,-0.671) -- (-0.572,-0.677);
\draw[mesh] (-1.178,-1.894) -- (-1.189,-1.885) -- (-1.198,-1.872) -- (-1.206,-1.855) -- (-1.211,-1.835) -- (-1.214,-1.811) -- (-1.215,-1.785);
\draw[mesh] (-0.629,-1.991) -- (-0.628,-1.961) -- (-0.627,-1.927) -- (-0.624,-1.891) -- (-0.620,-1.852) -- (-0.616,-1.810) -- (-0.610,-1.766) -- (-0.603,-1.720) -- (-0.595,-1.672) -- (-0.587,-1.623) -- (-0.577,-1.573) -- (-0.567,-1.521) -- (-0.556,-1.469) -- (-0.545,-1.417) -- (-0.533,-1.365) -- (-0.520,-1.314) -- (-0.507,-1.263) -- (-0.494,-1.213) -- (-0.480,-1.165) -- (-0.467,-1.118) -- (-0.453,-1.073) -- (-0.439,-1.031) -- (-0.425,-0.990) -- (-0.412,-0.953) -- (-0.399,-0.919) -- (-0.386,-0.887) -- (-0.373,-0.860) -- (-0.361,-0.835) -- (-0.349,-0.815) -- (-0.339,-0.798) -- (-0.328,-0.785) -- (-0.319,-0.776) -- (-0.311,-0.771) -- (-0.303,-0.770) -- (-0.296,-0.774);
\draw[mesh] (-0.610,-2.093) -- (-0.616,-2.086) -- (-0.620,-2.075) -- (-0.624,-2.059) -- (-0.627,-2.040) -- (-0.628,-2.017) -- (-0.629,-1.991);
\draw[mesh] (-0.000,-2.061) -- (-0.000,-2.031) -- (-0.000,-1.997) -- (-0.000,-1.961) -- (-0.000,-1.921) -- (-0.000,-1.879) -- (-0.000,-1.834) -- (-0.000,-1.787) -- (-0.000,-1.739) -- (-0.000,-1.689) -- (-0.000,-1.637) -- (-0.000,-1.585) -- (-0.000,-1.532) -- (-0.000,-1.478) -- (-0.000,-1.425) -- (-0.000,-1.372) -- (-0.000,-1.320) -- (-0.000,-1.268) -- (-0.000,-1.218) -- (-0.000,-1.170) -- (-0.000,-1.124) -- (-0.000,-1.080) -- (-0.000,-1.038) -- (-0.000,-0.999) -- (-0.000,-0.963) -- (-0.000,-0.930) -- (-0.000,-0.901) -- (-0.000,-0.876) -- (-0.000,-0.854) -- (-0.000,-0.836) -- (-0.000,-0.822) -- (-0.000,-0.812) -- (-0.000,-0.806) -- (-0.000,-0.804);
\draw[mesh] (-0.000,-2.164) -- (-0.000,-2.161) -- (-0.000,-2.155) -- (-0.000,-2.144) -- (-0.000,-2.129) -- (-0.000,-2.110) -- (-0.000,-2.087) -- (-0.000,-2.061);
\draw[mesh] (0.629,-1.991) -- (0.628,-1.961) -- (0.627,-1.927) -- (0.624,-1.891) -- (0.620,-1.852) -- (0.616,-1.810) -- (0.610,-1.766) -- (0.603,-1.720) -- (0.595,-1.672) -- (0.587,-1.623) -- (0.577,-1.573) -- (0.567,-1.521) -- (0.556,-1.469) -- (0.545,-1.417) -- (0.533,-1.365) -- (0.520,-1.314) -- (0.507,-1.263) -- (0.494,-1.213) -- (0.480,-1.165) -- (0.467,-1.118) -- (0.453,-1.073) -- (0.439,-1.031) -- (0.425,-0.990) -- (0.412,-0.953) -- (0.399,-0.919) -- (0.386,-0.887) -- (0.373,-0.860) -- (0.361,-0.835) -- (0.349,-0.815) -- (0.339,-0.798) -- (0.328,-0.785) -- (0.319,-0.776) -- (0.311,-0.771) -- (0.303,-0.770) -- (0.296,-0.774);
\draw[mesh] (0.610,-2.093) -- (0.616,-2.086) -- (0.620,-2.075) -- (0.624,-2.059) -- (0.627,-2.040) -- (0.628,-2.017) -- (0.629,-1.991);
\draw[mesh] (1.215,-1.785) -- (1.214,-1.755) -- (1.211,-1.722) -- (1.206,-1.687) -- (1.198,-1.649) -- (1.189,-1.609) -- (1.178,-1.567) -- (1.165,-1.523) -- (1.150,-1.477) -- (1.134,-1.431) -- (1.115,-1.384) -- (1.096,-1.336) -- (1.075,-1.287) -- (1.053,-1.239) -- (1.029,-1.191) -- (1.005,-1.143) -- (0.980,-1.097) -- (0.954,-1.051) -- (0.928,-1.007) -- (0.902,-0.965) -- (0.875,-0.925) -- (0.848,-0.887) -- (0.822,-0.851) -- (0.796,-0.818) -- (0.770,-0.788) -- (0.745,-0.761) -- (0.721,-0.737) -- (0.697,-0.717) -- (0.675,-0.700) -- (0.654,-0.687) -- (0.635,-0.677) -- (0.616,-0.671) -- (0.600,-0.669) -- (0.585,-0.671) -- (0.572,-0.677);
\draw[mesh] (1.178,-1.894) -- (1.189,-1.885) -- (1.198,-1.872) -- (1.206,-1.855) -- (1.211,-1.835) -- (1.214,-1.811) -- (1.215,-1.785);
\draw[mesh] (1.718,-1.457) -- (1.717,-1.428) -- (1.712,-1.396) -- (1.705,-1.362) -- (1.695,-1.326) -- (1.682,-1.288) -- (1.666,-1.249) -- (1.647,-1.209) -- (1.626,-1.167) -- (1.603,-1.125) -- (1.577,-1.083) -- (1.550,-1.040) -- (1.520,-0.998) -- (1.489,-0.955) -- (1.456,-0.913) -- (1.421,-0.873) -- (1.386,-0.833) -- (1.350,-0.794) -- (1.313,-0.757) -- (1.275,-0.722) -- (1.237,-0.689) -- (1.200,-0.658) -- (1.162,-0.630) -- (1.125,-0.604) -- (1.089,-0.581) -- (1.053,-0.560) -- (1.019,-0.543) -- (0.986,-0.529) -- (0.955,-0.518) -- (0.925,-0.511) -- (0.897,-0.506) -- (0.872,-0.505) -- (0.848,-0.508) -- (0.827,-0.513) -- (0.809,-0.522) -- (0.793,-0.535);
\draw[mesh] (1.682,-1.564) -- (1.695,-1.549) -- (1.705,-1.530) -- (1.712,-1.509) -- (1.717,-1.484) -- (1.718,-1.457);
\draw[mesh] (2.104,-1.030) -- (2.103,-1.001) -- (2.097,-0.970) -- (2.088,-0.938) -- (2.076,-0.905) -- (2.060,-0.871) -- (2.040,-0.835) -- (2.018,-0.800) -- (1.992,-0.764) -- (1.963,-0.727) -- (1.932,-0.691) -- (1.898,-0.655) -- (1.862,-0.620) -- (1.823,-0.585) -- (1.783,-0.552) -- (1.741,-0.519) -- (1.698,-0.488) -- (1.653,-0.459) -- (1.608,-0.431) -- (1.562,-0.405) -- (1.516,-0.382) -- (1.469,-0.360) -- (1.423,-0.341) -- (1.378,-0.324) -- (1.334,-0.310) -- (1.290,-0.299) -- (1.248,-0.290) -- (1.208,-0.284) -- (1.169,-0.281) -- (1.133,-0.281) -- (1.099,-0.283) -- (1.068,-0.289) -- (1.039,-0.297) -- (1.013,-0.308) -- (0.991,-0.322) -- (0.971,-0.338) -- (0.955,-0.356) -- (0.943,-0.378);
\draw[mesh] (2.088,-1.107) -- (2.097,-1.083) -- (2.103,-1.058) -- (2.104,-1.030);
\draw[mesh] (2.347,-0.533) -- (2.345,-0.505) -- (2.339,-0.475) -- (2.329,-0.445) -- (2.315,-0.415) -- (2.297,-0.384) -- (2.276,-0.354) -- (2.250,-0.323) -- (2.222,-0.293) -- (2.190,-0.264) -- (2.155,-0.235) -- (2.117,-0.207) -- (2.076,-0.180) -- (2.034,-0.155) -- (1.989,-0.131) -- (1.942,-0.108) -- (1.893,-0.088) -- (1.844,-0.069) -- (1.793,-0.052) -- (1.742,-0.037) -- (1.690,-0.024) -- (1.639,-0.013) -- (1.588,-0.005) -- (1.537,0.001) -- (1.487,0.005) -- (1.439,0.006) -- (1.392,0.005) -- (1.347,0.001) -- (1.304,-0.005) -- (1.264,-0.013) -- (1.226,-0.024) -- (1.191,-0.037) -- (1.159,-0.052) -- (1.130,-0.069) -- (1.105,-0.088) -- (1.084,-0.108) -- (1.066,-0.131) -- (1.052,-0.155) -- (1.042,-0.180);
\draw[mesh] (2.339,-0.588) -- (2.345,-0.561) -- (2.347,-0.533);
\draw[mesh] (-2.428,-0.081) -- (-2.423,-0.162) -- (-2.413,-0.242) -- (-2.400,-0.322) -- (-2.383,-0.402) -- (-2.363,-0.481) -- (-2.339,-0.559) -- (-2.311,-0.637) -- (-2.280,-0.713) -- (-2.245,-0.789) -- (-2.207,-0.863) -- (-2.165,-0.936) -- (-2.120,-1.007) -- (-2.072,-1.077) -- (-2.020,-1.145) -- (-1.966,-1.211) -- (-1.908,-1.276) -- (-1.848,-1.338) -- (-1.784,-1.399) -- (-1.718,-1.457) -- (-1.649,-1.513) -- (-1.578,-1.567) -- (-1.504,-1.618) -- (-1.428,-1.667) -- (-1.350,-1.713) -- (-1.270,-1.757) -- (-1.187,-1.798) -- (-1.103,-1.836) -- (-1.017,-1.871) -- (-0.930,-1.904) -- (-0.841,-1.933) -- (-0.751,-1.960) -- (-0.660,-1.983) -- (-0.567,-2.004) -- (-0.474,-2.021) -- (-0.380,-2.035) -- (-0.286,-2.046) -- (-0.191,-2.054) -- (-0.095,-2.059) -- (-0.000,-2.061) -- (0.095,-2.059) -- (0.191,-2.054) -- (0.286,-2.046) -- (0.380,-2.035) -- (0.474,-2.021) -- (0.567,-2.004) -- (0.660,-1.983) -- (0.751,-1.960) -- (0.841,-1.933) -- (0.930,-1.904) -- (1.017,-1.871) -- (1.103,-1.836) -- (1.187,-1.798) -- (1.270,-1.757) -- (1.350,-1.713) -- (1.428,-1.667) -- (1.504,-1.618) -- (1.578,-1.567) -- (1.649,-1.513) -- (1.718,-1.457) -- (1.784,-1.399) -- (1.848,-1.338) -- (1.908,-1.276) -- (1.966,-1.211) -- (2.020,-1.145) -- (2.072,-1.077) -- (2.120,-1.007) -- (2.165,-0.936) -- (2.207,-0.863) -- (2.245,-0.789) -- (2.280,-0.713) -- (2.311,-0.637) -- (2.339,-0.559) -- (2.363,-0.481) -- (2.383,-0.402) -- (2.400,-0.322) -- (2.413,-0.242) -- (2.423,-0.162) -- (2.428,-0.081) -- (2.430,-0.000);
\draw[mesh] (2.231,0.255) -- (2.229,0.329) -- (2.224,0.403) -- (2.215,0.477) -- (2.203,0.551) -- (2.188,0.624) -- (2.169,0.696) -- (2.147,0.768) -- (2.122,0.839) -- (2.093,0.910) -- (2.061,0.979) -- (2.026,1.047) -- (1.988,1.114) -- (1.946,1.179) -- (1.902,1.243) -- (1.855,1.306) -- (1.805,1.367) -- (1.752,1.426) -- (1.696,1.483) -- (1.638,1.539) -- (1.577,1.593) -- (1.514,1.644) -- (1.449,1.693) -- (1.381,1.741) -- (1.311,1.785) -- (1.239,1.828) -- (1.166,1.868) -- (1.090,1.905) -- (1.013,1.940) -- (0.934,1.973) -- (0.854,2.003) -- (0.772,2.030) -- (0.689,2.054) -- (0.606,2.076) -- (0.521,2.094) -- (0.435,2.110) -- (0.349,2.123) -- (0.262,2.134) -- (0.175,2.141) -- (0.088,2.145) -- (0.000,2.147) -- (-0.088,2.145) -- (-0.175,2.141) -- (-0.262,2.134) -- (-0.349,2.123) -- (-0.435,2.110) -- (-0.521,2.094) -- (-0.606,2.076) -- (-0.689,2.054) -- (-0.772,2.030) -- (-0.854,2.003) -- (-0.934,1.973) -- (-1.013,1.940) -- (-1.090,1.905) -- (-1.166,1.868) -- (-1.239,1.828) -- (-1.311,1.785) -- (-1.381,1.741) -- (-1.449,1.693) -- (-1.514,1.644) -- (-1.577,1.593) -- (-1.638,1.539) -- (-1.696,1.483) -- (-1.752,1.426) -- (-1.805,1.367) -- (-1.855,1.306) -- (-1.902,1.243) -- (-1.946,1.179) -- (-1.988,1.114) -- (-2.026,1.047) -- (-2.061,0.979) -- (-2.093,0.910) -- (-2.122,0.839) -- (-2.147,0.768) -- (-2.169,0.696) -- (-2.188,0.624) -- (-2.203,0.551) -- (-2.215,0.477) -- (-2.224,0.403) -- (-2.229,0.329) -- (-2.231,0.255) -- (-2.229,0.181) -- (-2.224,0.106) -- (-2.215,0.032) -- (-2.203,-0.041) -- (-2.188,-0.114) -- (-2.169,-0.187) -- (-2.147,-0.259) -- (-2.122,-0.330) -- (-2.093,-0.400) -- (-2.061,-0.469) -- (-2.026,-0.537) -- (-1.988,-0.604) -- (-1.946,-0.670) -- (-1.902,-0.734) -- (-1.855,-0.796) -- (-1.805,-0.857) -- (-1.752,-0.916) -- (-1.696,-0.974) -- (-1.638,-1.029) -- (-1.577,-1.083) -- (-1.514,-1.134) -- (-1.449,-1.184) -- (-1.381,-1.231) -- (-1.311,-1.276) -- (-1.239,-1.318) -- (-1.166,-1.358) -- (-1.090,-1.396) -- (-1.013,-1.431) -- (-0.934,-1.463) -- (-0.854,-1.493) -- (-0.772,-1.520) -- (-0.689,-1.544) -- (-0.606,-1.566) -- (-0.521,-1.585) -- (-0.435,-1.601) -- (-0.349,-1.614) -- (-0.262,-1.624) -- (-0.175,-1.631) -- (-0.088,-1.636) -- (-0.000,-1.637) -- (0.088,-1.636) -- (0.175,-1.631) -- (0.262,-1.624) -- (0.349,-1.614) -- (0.435,-1.601) -- (0.521,-1.585) -- (0.606,-1.566) -- (0.689,-1.544) -- (0.772,-1.520) -- (0.854,-1.493) -- (0.934,-1.463) -- (1.013,-1.431) -- (1.090,-1.396) -- (1.166,-1.358) -- (1.239,-1.318) -- (1.311,-1.276) -- (1.381,-1.231) -- (1.449,-1.184) -- (1.514,-1.134) -- (1.577,-1.083) -- (1.638,-1.029) -- (1.696,-0.974) -- (1.752,-0.916) -- (1.805,-0.857) -- (1.855,-0.796) -- (1.902,-0.734) -- (1.946,-0.670) -- (1.988,-0.604) -- (2.026,-0.537) -- (2.061,-0.469) -- (2.093,-0.400) -- (2.122,-0.330) -- (2.147,-0.259) -- (2.169,-0.187) -- (2.188,-0.114) -- (2.203,-0.041) -- (2.215,0.032) -- (2.224,0.106) -- (2.229,0.181) -- (2.231,0.255);
\draw[mesh] (1.750,0.360) -- (1.749,0.419) -- (1.745,0.477) -- (1.738,0.535) -- (1.728,0.593) -- (1.716,0.650) -- (1.702,0.707) -- (1.684,0.763) -- (1.664,0.819) -- (1.642,0.874) -- (1.617,0.928) -- (1.589,0.982) -- (1.559,1.034) -- (1.527,1.086) -- (1.492,1.136) -- (1.455,1.185) -- (1.416,1.233) -- (1.374,1.279) -- (1.331,1.324) -- (1.285,1.368) -- (1.237,1.410) -- (1.188,1.450) -- (1.137,1.489) -- (1.083,1.526) -- (1.029,1.561) -- (0.972,1.594) -- (0.914,1.626) -- (0.855,1.655) -- (0.794,1.683) -- (0.733,1.708) -- (0.670,1.731) -- (0.606,1.753) -- (0.541,1.772) -- (0.475,1.789) -- (0.409,1.803) -- (0.341,1.816) -- (0.274,1.826) -- (0.206,1.834) -- (0.137,1.840) -- (0.069,1.843) -- (0.000,1.844) -- (-0.069,1.843) -- (-0.137,1.840) -- (-0.206,1.834) -- (-0.274,1.826) -- (-0.341,1.816) -- (-0.409,1.803) -- (-0.475,1.789) -- (-0.541,1.772) -- (-0.606,1.753) -- (-0.670,1.731) -- (-0.733,1.708) -- (-0.794,1.683) -- (-0.855,1.655) -- (-0.914,1.626) -- (-0.972,1.594) -- (-1.029,1.561) -- (-1.083,1.526) -- (-1.137,1.489) -- (-1.188,1.450) -- (-1.237,1.410) -- (-1.285,1.368) -- (-1.331,1.324) -- (-1.374,1.279) -- (-1.416,1.233) -- (-1.455,1.185) -- (-1.492,1.136) -- (-1.527,1.086) -- (-1.559,1.034) -- (-1.589,0.982) -- (-1.617,0.928) -- (-1.642,0.874) -- (-1.664,0.819) -- (-1.684,0.763) -- (-1.702,0.707) -- (-1.716,0.650) -- (-1.728,0.593) -- (-1.738,0.535) -- (-1.745,0.477) -- (-1.749,0.419) -- (-1.750,0.360) -- (-1.749,0.302) -- (-1.745,0.244) -- (-1.738,0.186) -- (-1.728,0.128) -- (-1.716,0.071) -- (-1.702,0.014) -- (-1.684,-0.042) -- (-1.664,-0.098) -- (-1.642,-0.153) -- (-1.617,-0.208) -- (-1.589,-0.261) -- (-1.559,-0.313) -- (-1.527,-0.365) -- (-1.492,-0.415) -- (-1.455,-0.464) -- (-1.416,-0.512) -- (-1.374,-0.558) -- (-1.331,-0.603) -- (-1.285,-0.647) -- (-1.237,-0.689) -- (-1.188,-0.729) -- (-1.137,-0.768) -- (-1.083,-0.805) -- (-1.029,-0.840) -- (-0.972,-0.874) -- (-0.914,-0.905) -- (-0.855,-0.935) -- (-0.794,-0.962) -- (-0.733,-0.987) -- (-0.670,-1.011) -- (-0.606,-1.032) -- (-0.541,-1.051) -- (-0.475,-1.068) -- (-0.409,-1.083) -- (-0.341,-1.095) -- (-0.274,-1.105) -- (-0.206,-1.113) -- (-0.137,-1.119) -- (-0.069,-1.123) -- (-0.000,-1.124) -- (0.069,-1.123) -- (0.137,-1.119) -- (0.206,-1.113) -- (0.274,-1.105) -- (0.341,-1.095) -- (0.409,-1.083) -- (0.475,-1.068) -- (0.541,-1.051) -- (0.606,-1.032) -- (0.670,-1.011) -- (0.733,-0.987) -- (0.794,-0.962) -- (0.855,-0.935) -- (0.914,-0.905) -- (0.972,-0.874) -- (1.029,-0.840) -- (1.083,-0.805) -- (1.137,-0.768) -- (1.188,-0.729) -- (1.237,-0.689) -- (1.285,-0.647) -- (1.331,-0.603) -- (1.374,-0.558) -- (1.416,-0.512) -- (1.455,-0.464) -- (1.492,-0.415) -- (1.527,-0.365) -- (1.559,-0.313) -- (1.589,-0.261) -- (1.617,-0.208) -- (1.642,-0.153) -- (1.664,-0.098) -- (1.684,-0.042) -- (1.702,0.014) -- (1.716,0.071) -- (1.728,0.128) -- (1.738,0.186) -- (1.745,0.244) -- (1.749,0.302) -- (1.750,0.360);
\draw[mesh] (1.269,0.255) -- (1.268,0.297) -- (1.265,0.339) -- (1.260,0.381) -- (1.254,0.423) -- (1.245,0.465) -- (1.234,0.506) -- (1.222,0.547) -- (1.207,0.587) -- (1.191,0.627) -- (1.173,0.667) -- (1.153,0.705) -- (1.131,0.743) -- (1.107,0.781) -- (1.082,0.817) -- (1.055,0.853) -- (1.027,0.887) -- (0.997,0.921) -- (0.965,0.954) -- (0.932,0.985) -- (0.897,1.016) -- (0.862,1.045) -- (0.824,1.073) -- (0.786,1.100) -- (0.746,1.126) -- (0.705,1.150) -- (0.663,1.173) -- (0.620,1.194) -- (0.576,1.214) -- (0.531,1.232) -- (0.486,1.249) -- (0.439,1.265) -- (0.392,1.278) -- (0.345,1.291) -- (0.296,1.301) -- (0.248,1.310) -- (0.199,1.318) -- (0.149,1.324) -- (0.100,1.328) -- (0.050,1.330) -- (0.000,1.331) -- (-0.050,1.330) -- (-0.100,1.328) -- (-0.149,1.324) -- (-0.199,1.318) -- (-0.248,1.310) -- (-0.296,1.301) -- (-0.345,1.291) -- (-0.392,1.278) -- (-0.439,1.265) -- (-0.486,1.249) -- (-0.531,1.232) -- (-0.576,1.214) -- (-0.620,1.194) -- (-0.663,1.173) -- (-0.705,1.150) -- (-0.746,1.126) -- (-0.786,1.100) -- (-0.824,1.073) -- (-0.862,1.045) -- (-0.897,1.016) -- (-0.932,0.985) -- (-0.965,0.954) -- (-0.997,0.921) -- (-1.027,0.887) -- (-1.055,0.853) -- (-1.082,0.817) -- (-1.107,0.781) -- (-1.131,0.743) -- (-1.153,0.705) -- (-1.173,0.667) -- (-1.191,0.627) -- (-1.207,0.587) -- (-1.222,0.547) -- (-1.234,0.506) -- (-1.245,0.465) -- (-1.254,0.423) -- (-1.260,0.381) -- (-1.265,0.339) -- (-1.268,0.297) -- (-1.269,0.255) -- (-1.268,0.213) -- (-1.265,0.170) -- (-1.260,0.128) -- (-1.254,0.086) -- (-1.245,0.045) -- (-1.234,0.004) -- (-1.222,-0.037) -- (-1.207,-0.078) -- (-1.191,-0.118) -- (-1.173,-0.157) -- (-1.153,-0.196) -- (-1.131,-0.234) -- (-1.107,-0.271) -- (-1.082,-0.308) -- (-1.055,-0.343) -- (-1.027,-0.378) -- (-0.997,-0.412) -- (-0.965,-0.444) -- (-0.932,-0.476) -- (-0.897,-0.506) -- (-0.862,-0.536) -- (-0.824,-0.564) -- (-0.786,-0.590) -- (-0.746,-0.616) -- (-0.705,-0.640) -- (-0.663,-0.663) -- (-0.620,-0.684) -- (-0.576,-0.704) -- (-0.531,-0.723) -- (-0.486,-0.740) -- (-0.439,-0.755) -- (-0.392,-0.769) -- (-0.345,-0.781) -- (-0.296,-0.792) -- (-0.248,-0.801) -- (-0.199,-0.808) -- (-0.149,-0.814) -- (-0.100,-0.818) -- (-0.050,-0.821) -- (-0.000,-0.822) -- (0.050,-0.821) -- (0.100,-0.818) -- (0.149,-0.814) -- (0.199,-0.808) -- (0.248,-0.801) -- (0.296,-0.792) -- (0.345,-0.781) -- (0.392,-0.769) -- (0.439,-0.755) -- (0.486,-0.740) -- (0.531,-0.723) -- (0.576,-0.704) -- (0.620,-0.684) -- (0.663,-0.663) -- (0.705,-0.640) -- (0.746,-0.616) -- (0.786,-0.590) -- (0.824,-0.564) -- (0.862,-0.536) -- (0.897,-0.506) -- (0.932,-0.476) -- (0.965,-0.444) -- (0.997,-0.412) -- (1.027,-0.378) -- (1.055,-0.343) -- (1.082,-0.308) -- (1.107,-0.271) -- (1.131,-0.234) -- (1.153,-0.196) -- (1.173,-0.157) -- (1.191,-0.118) -- (1.207,-0.078) -- (1.222,-0.037) -- (1.234,0.004) -- (1.245,0.045) -- (1.254,0.086) -- (1.260,0.128) -- (1.265,0.170) -- (1.268,0.213) -- (1.269,0.255);
\draw[mesh] (1.070,0.000) -- (1.069,0.036) -- (1.067,0.071) -- (1.063,0.107) -- (1.057,0.142) -- (1.049,0.177) -- (1.040,0.212) -- (1.030,0.246) -- (1.018,0.280) -- (1.004,0.314) -- (0.989,0.347) -- (0.972,0.380) -- (0.953,0.412) -- (0.934,0.443) -- (0.912,0.474) -- (0.890,0.504) -- (0.866,0.533) -- (0.840,0.562) -- (0.814,0.589) -- (0.786,0.616) -- (0.757,0.642) -- (0.726,0.666) -- (0.695,0.690) -- (0.662,0.713) -- (0.629,0.734) -- (0.594,0.754) -- (0.559,0.774) -- (0.523,0.792) -- (0.486,0.809) -- (0.448,0.824) -- (0.409,0.838) -- (0.370,0.851) -- (0.331,0.863) -- (0.290,0.873) -- (0.250,0.882) -- (0.209,0.890) -- (0.167,0.896) -- (0.126,0.901) -- (0.084,0.905) -- (0.042,0.907) -- (0.000,0.907) -- (-0.042,0.907) -- (-0.084,0.905) -- (-0.126,0.901) -- (-0.167,0.896) -- (-0.209,0.890) -- (-0.250,0.882) -- (-0.290,0.873) -- (-0.331,0.863) -- (-0.370,0.851) -- (-0.409,0.838) -- (-0.448,0.824) -- (-0.486,0.809) -- (-0.523,0.792) -- (-0.559,0.774) -- (-0.594,0.754) -- (-0.629,0.734) -- (-0.662,0.713) -- (-0.695,0.690) -- (-0.726,0.666) -- (-0.757,0.642) -- (-0.786,0.616) -- (-0.814,0.589) -- (-0.840,0.562) -- (-0.866,0.533) -- (-0.890,0.504) -- (-0.912,0.474) -- (-0.934,0.443) -- (-0.953,0.412) -- (-0.972,0.380) -- (-0.989,0.347) -- (-1.004,0.314) -- (-1.018,0.280) -- (-1.030,0.246) -- (-1.040,0.212) -- (-1.049,0.177) -- (-1.057,0.142) -- (-1.063,0.107) -- (-1.067,0.071) -- (-1.069,0.036) -- (-1.070,0.000) -- (-1.069,-0.036);

    % the (2, 3) torus knot, generated by torus-knot-closed-and-open.py; closes after t = 2 pi
\draw[knotback] (2.430,0.000) -- (2.430,0.018);
\draw[knot] (2.430,0.018) -- (2.430,0.036) -- (2.429,0.054) -- (2.428,0.073) -- (2.428,0.091) -- (2.427,0.109) -- (2.425,0.127) -- (2.424,0.145) -- (2.422,0.163) -- (2.420,0.181) -- (2.418,0.199) -- (2.416,0.217) -- (2.414,0.235) -- (2.411,0.253) -- (2.408,0.271) -- (2.405,0.289) -- (2.402,0.307) -- (2.399,0.325) -- (2.395,0.342) -- (2.392,0.360) -- (2.388,0.378) -- (2.384,0.395) -- (2.379,0.413) -- (2.375,0.430) -- (2.370,0.448) -- (2.365,0.465) -- (2.360,0.482) -- (2.355,0.500) -- (2.350,0.517) -- (2.344,0.534) -- (2.339,0.551) -- (2.333,0.568) -- (2.327,0.585) -- (2.320,0.602) -- (2.314,0.618) -- (2.307,0.635) -- (2.301,0.652) -- (2.294,0.668) -- (2.287,0.685) -- (2.279,0.701) -- (2.272,0.717) -- (2.264,0.733) -- (2.257,0.749) -- (2.249,0.765) -- (2.241,0.781) -- (2.232,0.797) -- (2.224,0.812) -- (2.215,0.828) -- (2.207,0.843) -- (2.198,0.859) -- (2.189,0.874) -- (2.180,0.889) -- (2.170,0.904) -- (2.161,0.919) -- (2.151,0.933) -- (2.142,0.948) -- (2.132,0.962) -- (2.122,0.977) -- (2.112,0.991) -- (2.101,1.005) -- (2.091,1.019) -- (2.080,1.033) -- (2.070,1.047) -- (2.059,1.060) -- (2.048,1.074) -- (2.037,1.087) -- (2.026,1.100) -- (2.014,1.113) -- (2.003,1.126) -- (1.991,1.139) -- (1.980,1.152) -- (1.968,1.164) -- (1.956,1.176) -- (1.944,1.189) -- (1.932,1.201) -- (1.920,1.213) -- (1.907,1.224) -- (1.895,1.236) -- (1.883,1.247) -- (1.870,1.259) -- (1.857,1.270) -- (1.844,1.281) -- (1.832,1.292) -- (1.819,1.302) -- (1.806,1.313) -- (1.792,1.323) -- (1.779,1.333) -- (1.766,1.343) -- (1.753,1.353) -- (1.739,1.363) -- (1.726,1.373) -- (1.712,1.382) -- (1.698,1.391) -- (1.685,1.400) -- (1.671,1.409) -- (1.657,1.418) -- (1.643,1.427) -- (1.629,1.435) -- (1.615,1.443) -- (1.601,1.451) -- (1.587,1.459) -- (1.573,1.467) -- (1.559,1.475) -- (1.544,1.482) -- (1.530,1.489) -- (1.516,1.496) -- (1.501,1.503) -- (1.487,1.510) -- (1.472,1.517) -- (1.458,1.523) -- (1.443,1.529) -- (1.429,1.535) -- (1.414,1.541) -- (1.400,1.547) -- (1.385,1.553) -- (1.370,1.558) -- (1.356,1.563) -- (1.341,1.568) -- (1.326,1.573) -- (1.312,1.578) -- (1.297,1.583) -- (1.282,1.587) -- (1.267,1.591) -- (1.253,1.595) -- (1.238,1.599) -- (1.223,1.603) -- (1.209,1.607) -- (1.194,1.610) -- (1.179,1.613) -- (1.164,1.616) -- (1.150,1.619) -- (1.135,1.622) -- (1.120,1.625) -- (1.106,1.627) -- (1.091,1.630) -- (1.077,1.632) -- (1.062,1.634) -- (1.047,1.636) -- (1.033,1.637) -- (1.018,1.639) -- (1.004,1.640) -- (0.989,1.641) -- (0.975,1.642) -- (0.961,1.643) -- (0.946,1.644) -- (0.932,1.645) -- (0.918,1.645) -- (0.903,1.645) -- (0.889,1.646) -- (0.875,1.646) -- (0.861,1.645) -- (0.847,1.645) -- (0.833,1.645) -- (0.819,1.644) -- (0.805,1.644) -- (0.791,1.643) -- (0.777,1.642) -- (0.763,1.641) -- (0.750,1.639) -- (0.736,1.638) -- (0.722,1.637) -- (0.709,1.635) -- (0.695,1.633) -- (0.682,1.631) -- (0.669,1.629) -- (0.655,1.627) -- (0.642,1.625) -- (0.629,1.622) -- (0.616,1.620) -- (0.603,1.617) -- (0.590,1.615) -- (0.577,1.612) -- (0.564,1.609) -- (0.551,1.606) -- (0.538,1.602) -- (0.526,1.599) -- (0.513,1.596) -- (0.501,1.592) -- (0.488,1.588) -- (0.476,1.585) -- (0.464,1.581) -- (0.451,1.577) -- (0.439,1.573) -- (0.427,1.569) -- (0.415,1.564) -- (0.403,1.560) -- (0.391,1.556) -- (0.380,1.551) -- (0.368,1.546) -- (0.356,1.542) -- (0.345,1.537) -- (0.333,1.532) -- (0.322,1.527) -- (0.311,1.522) -- (0.300,1.517) -- (0.289,1.511) -- (0.277,1.506) -- (0.267,1.501) -- (0.256,1.495) -- (0.245,1.490) -- (0.234,1.484) -- (0.223,1.478) -- (0.213,1.473) -- (0.202,1.467) -- (0.192,1.461) -- (0.182,1.455) -- (0.171,1.449) -- (0.161,1.443) -- (0.151,1.437) -- (0.141,1.430) -- (0.131,1.424) -- (0.121,1.418) -- (0.112,1.411) -- (0.102,1.405) -- (0.092,1.398) -- (0.083,1.392) -- (0.073,1.385) -- (0.064,1.379) -- (0.054,1.372) -- (0.045,1.365) -- (0.036,1.359) -- (0.027,1.352) -- (0.018,1.345) -- (0.009,1.338) -- (0.000,1.331) -- (-0.009,1.324) -- (-0.018,1.317) -- (-0.026,1.310) -- (-0.035,1.303) -- (-0.043,1.296) -- (-0.052,1.289) -- (-0.060,1.282) -- (-0.069,1.275) -- (-0.077,1.267) -- (-0.085,1.260) -- (-0.093,1.253) -- (-0.101,1.246) -- (-0.110,1.238) -- (-0.117,1.231) -- (-0.125,1.224) -- (-0.133,1.216) -- (-0.141,1.209) -- (-0.149,1.202) -- (-0.157,1.194) -- (-0.164,1.187) -- (-0.172,1.180) -- (-0.179,1.172) -- (-0.187,1.165) -- (-0.194,1.157) -- (-0.202,1.150) -- (-0.209,1.142) -- (-0.216,1.135) -- (-0.224,1.128) -- (-0.231,1.120) -- (-0.238,1.113) -- (-0.245,1.105) -- (-0.252,1.098) -- (-0.259,1.090) -- (-0.266,1.083) -- (-0.273,1.075) -- (-0.280,1.068) -- (-0.287,1.061) -- (-0.294,1.053) -- (-0.301,1.046) -- (-0.307,1.038) -- (-0.314,1.031) -- (-0.321,1.023) -- (-0.328,1.016) -- (-0.334,1.009) -- (-0.341,1.001) -- (-0.348,0.994) -- (-0.354,0.986) -- (-0.361,0.979) -- (-0.367,0.972) -- (-0.374,0.964) -- (-0.380,0.957) -- (-0.387,0.950) -- (-0.393,0.942) -- (-0.400,0.935) -- (-0.406,0.928) -- (-0.413,0.921) -- (-0.419,0.913) -- (-0.426,0.906) -- (-0.432,0.899) -- (-0.439,0.892) -- (-0.445,0.885) -- (-0.451,0.877) -- (-0.458,0.870) -- (-0.464,0.863) -- (-0.471,0.856) -- (-0.477,0.849) -- (-0.484,0.842) -- (-0.490,0.835) -- (-0.496,0.828) -- (-0.503,0.821) -- (-0.509,0.814) -- (-0.516,0.807) -- (-0.522,0.800) -- (-0.529,0.793) -- (-0.535,0.786) -- (-0.541,0.779) -- (-0.548,0.772) -- (-0.554,0.765) -- (-0.561,0.758) -- (-0.568,0.751) -- (-0.574,0.744) -- (-0.581,0.738) -- (-0.587,0.731) -- (-0.594,0.724) -- (-0.600,0.717) -- (-0.607,0.710) -- (-0.614,0.704) -- (-0.620,0.697) -- (-0.627,0.690) -- (-0.634,0.683) -- (-0.641,0.677) -- (-0.647,0.670) -- (-0.654,0.663) -- (-0.661,0.657) -- (-0.668,0.650) -- (-0.675,0.643) -- (-0.682,0.637) -- (-0.689,0.630) -- (-0.696,0.624) -- (-0.703,0.617) -- (-0.710,0.610) -- (-0.717,0.604) -- (-0.724,0.597) -- (-0.731,0.591) -- (-0.738,0.584) -- (-0.745,0.577) -- (-0.753,0.571) -- (-0.760,0.564) -- (-0.767,0.558) -- (-0.775,0.551) -- (-0.782,0.545) -- (-0.789,0.538) -- (-0.797,0.532) -- (-0.804,0.525) -- (-0.812,0.518);
\draw[knotback] (-0.812,0.518) -- (-0.820,0.512) -- (-0.827,0.505) -- (-0.835,0.499) -- (-0.843,0.492) -- (-0.850,0.486) -- (-0.858,0.479) -- (-0.866,0.473) -- (-0.874,0.466) -- (-0.882,0.459) -- (-0.889,0.453) -- (-0.897,0.446) -- (-0.905,0.440) -- (-0.913,0.433) -- (-0.922,0.426) -- (-0.930,0.420) -- (-0.938,0.413) -- (-0.946,0.406) -- (-0.954,0.400) -- (-0.962,0.393) -- (-0.971,0.386) -- (-0.979,0.380) -- (-0.987,0.373) -- (-0.996,0.366) -- (-1.004,0.359) -- (-1.013,0.353) -- (-1.021,0.346) -- (-1.030,0.339) -- (-1.038,0.332) -- (-1.047,0.325) -- (-1.056,0.318) -- (-1.064,0.311) -- (-1.073,0.304) -- (-1.082,0.297) -- (-1.090,0.290) -- (-1.099,0.283) -- (-1.108,0.276) -- (-1.117,0.269) -- (-1.126,0.262) -- (-1.135,0.255) -- (-1.143,0.248) -- (-1.152,0.241) -- (-1.161,0.233) -- (-1.170,0.226) -- (-1.179,0.219) -- (-1.188,0.211) -- (-1.197,0.204) -- (-1.206,0.197) -- (-1.215,0.189) -- (-1.224,0.182) -- (-1.234,0.174) -- (-1.243,0.167) -- (-1.252,0.159) -- (-1.261,0.151) -- (-1.270,0.144) -- (-1.279,0.136) -- (-1.288,0.128) -- (-1.297,0.121) -- (-1.307,0.113) -- (-1.316,0.105) -- (-1.325,0.097) -- (-1.334,0.089) -- (-1.343,0.081) -- (-1.352,0.073) -- (-1.362,0.065) -- (-1.371,0.057) -- (-1.380,0.048) -- (-1.389,0.040) -- (-1.398,0.032) -- (-1.407,0.024) -- (-1.416,0.015) -- (-1.425,0.007) -- (-1.434,-0.002) -- (-1.444,-0.010) -- (-1.453,-0.019) -- (-1.462,-0.027) -- (-1.471,-0.036) -- (-1.480,-0.045) -- (-1.488,-0.054) -- (-1.497,-0.062) -- (-1.506,-0.071) -- (-1.515,-0.080) -- (-1.524,-0.089) -- (-1.533,-0.098) -- (-1.541,-0.107) -- (-1.550,-0.116) -- (-1.559,-0.125) -- (-1.567,-0.135) -- (-1.576,-0.144) -- (-1.584,-0.153) -- (-1.593,-0.163) -- (-1.601,-0.172) -- (-1.610,-0.182) -- (-1.618,-0.191) -- (-1.626,-0.201) -- (-1.635,-0.210) -- (-1.643,-0.220) -- (-1.651,-0.230) -- (-1.659,-0.239) -- (-1.667,-0.249) -- (-1.675,-0.259) -- (-1.683,-0.269) -- (-1.690,-0.279) -- (-1.698,-0.289) -- (-1.706,-0.299) -- (-1.713,-0.309) -- (-1.721,-0.319) -- (-1.728,-0.329) -- (-1.736,-0.340) -- (-1.743,-0.350) -- (-1.750,-0.360) -- (-1.757,-0.371) -- (-1.764,-0.381) -- (-1.771,-0.392) -- (-1.778,-0.402) -- (-1.785,-0.413) -- (-1.791,-0.423) -- (-1.798,-0.434) -- (-1.804,-0.445) -- (-1.810,-0.455) -- (-1.817,-0.466) -- (-1.823,-0.477) -- (-1.829,-0.488) -- (-1.835,-0.499) -- (-1.841,-0.510) -- (-1.846,-0.520) -- (-1.852,-0.531) -- (-1.857,-0.542) -- (-1.863,-0.554) -- (-1.868,-0.565) -- (-1.873,-0.576) -- (-1.878,-0.587) -- (-1.883,-0.598) -- (-1.888,-0.609) -- (-1.892,-0.620) -- (-1.897,-0.632) -- (-1.901,-0.643) -- (-1.905,-0.654) -- (-1.909,-0.666) -- (-1.913,-0.677) -- (-1.917,-0.688) -- (-1.921,-0.700) -- (-1.925,-0.711) -- (-1.928,-0.723) -- (-1.931,-0.734) -- (-1.934,-0.745) -- (-1.937,-0.757) -- (-1.940,-0.768) -- (-1.943,-0.780) -- (-1.946,-0.791) -- (-1.948,-0.803) -- (-1.950,-0.814) -- (-1.952,-0.826) -- (-1.954,-0.837) -- (-1.956,-0.849) -- (-1.958,-0.861) -- (-1.959,-0.872) -- (-1.961,-0.884) -- (-1.962,-0.895) -- (-1.963,-0.907) -- (-1.964,-0.918) -- (-1.965,-0.930) -- (-1.965,-0.941) -- (-1.966,-0.953) -- (-1.966,-0.964) -- (-1.966,-0.976) -- (-1.966,-0.987) -- (-1.966,-0.999) -- (-1.965,-1.010) -- (-1.965,-1.022) -- (-1.964,-1.033) -- (-1.963,-1.044) -- (-1.962,-1.056) -- (-1.961,-1.067) -- (-1.959,-1.078) -- (-1.958,-1.090) -- (-1.956,-1.101) -- (-1.954,-1.112) -- (-1.952,-1.123) -- (-1.949,-1.135) -- (-1.947,-1.146) -- (-1.944,-1.157) -- (-1.941,-1.168) -- (-1.938,-1.179) -- (-1.935,-1.190) -- (-1.932,-1.201) -- (-1.928,-1.212) -- (-1.925,-1.222) -- (-1.921,-1.233) -- (-1.917,-1.244) -- (-1.913,-1.255) -- (-1.908,-1.265) -- (-1.904,-1.276) -- (-1.899,-1.286) -- (-1.894,-1.297) -- (-1.889,-1.307) -- (-1.884,-1.318) -- (-1.878,-1.328) -- (-1.873,-1.338) -- (-1.867,-1.348) -- (-1.861,-1.358) -- (-1.855,-1.368) -- (-1.848,-1.378) -- (-1.842,-1.388) -- (-1.835,-1.398) -- (-1.828,-1.408) -- (-1.821,-1.417) -- (-1.814,-1.427) -- (-1.807,-1.436) -- (-1.799,-1.446) -- (-1.792,-1.455) -- (-1.784,-1.464) -- (-1.776,-1.474) -- (-1.768,-1.483) -- (-1.759,-1.492) -- (-1.751,-1.500) -- (-1.742,-1.509) -- (-1.733,-1.518) -- (-1.724,-1.527) -- (-1.715,-1.535) -- (-1.706,-1.543) -- (-1.696,-1.552);
\draw[knot] (-1.696,-1.552) -- (-1.687,-1.560) -- (-1.677,-1.568) -- (-1.667,-1.576) -- (-1.657,-1.584) -- (-1.646,-1.592) -- (-1.636,-1.599) -- (-1.625,-1.607) -- (-1.615,-1.614) -- (-1.604,-1.622) -- (-1.593,-1.629) -- (-1.581,-1.636) -- (-1.570,-1.643) -- (-1.559,-1.650) -- (-1.547,-1.657) -- (-1.535,-1.663) -- (-1.523,-1.670) -- (-1.511,-1.676) -- (-1.499,-1.683) -- (-1.487,-1.689) -- (-1.474,-1.695) -- (-1.462,-1.701) -- (-1.449,-1.707) -- (-1.436,-1.712) -- (-1.423,-1.718) -- (-1.410,-1.723) -- (-1.397,-1.728) -- (-1.384,-1.734) -- (-1.370,-1.739) -- (-1.357,-1.743) -- (-1.343,-1.748) -- (-1.329,-1.753) -- (-1.315,-1.757) -- (-1.301,-1.762) -- (-1.287,-1.766) -- (-1.273,-1.770) -- (-1.259,-1.774) -- (-1.244,-1.778) -- (-1.230,-1.781) -- (-1.215,-1.785) -- (-1.200,-1.788) -- (-1.185,-1.791) -- (-1.170,-1.794) -- (-1.155,-1.797) -- (-1.140,-1.800) -- (-1.125,-1.803) -- (-1.110,-1.805) -- (-1.095,-1.807) -- (-1.079,-1.810) -- (-1.064,-1.812) -- (-1.048,-1.814) -- (-1.032,-1.815) -- (-1.017,-1.817) -- (-1.001,-1.818) -- (-0.985,-1.820) -- (-0.969,-1.821) -- (-0.953,-1.822) -- (-0.937,-1.823) -- (-0.921,-1.824) -- (-0.905,-1.824) -- (-0.889,-1.825) -- (-0.872,-1.825) -- (-0.856,-1.825) -- (-0.840,-1.825) -- (-0.823,-1.825) -- (-0.807,-1.824) -- (-0.790,-1.824) -- (-0.774,-1.823) -- (-0.757,-1.823) -- (-0.741,-1.822) -- (-0.724,-1.821) -- (-0.707,-1.819) -- (-0.691,-1.818) -- (-0.674,-1.817) -- (-0.657,-1.815) -- (-0.641,-1.813) -- (-0.624,-1.811) -- (-0.607,-1.809) -- (-0.590,-1.807) -- (-0.574,-1.805) -- (-0.557,-1.802) -- (-0.540,-1.799) -- (-0.523,-1.797) -- (-0.507,-1.794) -- (-0.490,-1.791) -- (-0.473,-1.787) -- (-0.456,-1.784) -- (-0.440,-1.781) -- (-0.423,-1.777) -- (-0.406,-1.773) -- (-0.389,-1.769) -- (-0.373,-1.765) -- (-0.356,-1.761) -- (-0.340,-1.757) -- (-0.323,-1.752) -- (-0.306,-1.748) -- (-0.290,-1.743) -- (-0.273,-1.738) -- (-0.257,-1.733) -- (-0.240,-1.728) -- (-0.224,-1.723) -- (-0.208,-1.718) -- (-0.191,-1.712) -- (-0.175,-1.706) -- (-0.159,-1.701) -- (-0.143,-1.695) -- (-0.127,-1.689) -- (-0.111,-1.683) -- (-0.095,-1.677) -- (-0.079,-1.670) -- (-0.063,-1.664) -- (-0.047,-1.657) -- (-0.031,-1.651) -- (-0.016,-1.644) -- (-0.000,-1.637) -- (0.016,-1.630) -- (0.031,-1.623) -- (0.046,-1.616) -- (0.062,-1.608) -- (0.077,-1.601) -- (0.092,-1.594) -- (0.107,-1.586) -- (0.122,-1.578) -- (0.137,-1.570) -- (0.152,-1.562) -- (0.167,-1.554) -- (0.181,-1.546) -- (0.196,-1.538) -- (0.210,-1.530) -- (0.225,-1.522) -- (0.239,-1.513) -- (0.253,-1.505) -- (0.267,-1.496) -- (0.281,-1.487) -- (0.295,-1.478) -- (0.309,-1.470) -- (0.323,-1.461) -- (0.336,-1.452) -- (0.350,-1.443) -- (0.363,-1.433) -- (0.376,-1.424) -- (0.389,-1.415) -- (0.402,-1.406) -- (0.415,-1.396) -- (0.428,-1.387) -- (0.441,-1.377) -- (0.453,-1.368) -- (0.466,-1.358) -- (0.478,-1.348) -- (0.490,-1.338) -- (0.502,-1.329) -- (0.514,-1.319) -- (0.526,-1.309) -- (0.538,-1.299) -- (0.550,-1.289) -- (0.561,-1.279) -- (0.572,-1.269) -- (0.584,-1.259) -- (0.595,-1.248) -- (0.606,-1.238) -- (0.617,-1.228) -- (0.627,-1.218) -- (0.638,-1.207) -- (0.648,-1.197) -- (0.659,-1.187) -- (0.669,-1.176) -- (0.679,-1.166) -- (0.689,-1.156) -- (0.699,-1.145) -- (0.709,-1.135) -- (0.718,-1.124) -- (0.728,-1.114) -- (0.737,-1.103) -- (0.746,-1.093) -- (0.755,-1.082) -- (0.764,-1.072) -- (0.773,-1.061) -- (0.781,-1.051) -- (0.790,-1.040) -- (0.798,-1.030) -- (0.807,-1.019) -- (0.815,-1.009) -- (0.823,-0.998) -- (0.831,-0.988) -- (0.838,-0.977) -- (0.846,-0.967) -- (0.853,-0.956) -- (0.861,-0.946) -- (0.868,-0.935) -- (0.875,-0.925) -- (0.882,-0.915) -- (0.889,-0.904) -- (0.895,-0.894) -- (0.902,-0.883) -- (0.908,-0.873) -- (0.915,-0.863) -- (0.921,-0.853) -- (0.927,-0.842) -- (0.933,-0.832) -- (0.939,-0.822) -- (0.945,-0.812) -- (0.950,-0.802) -- (0.956,-0.792) -- (0.961,-0.782) -- (0.966,-0.772) -- (0.971,-0.762) -- (0.976,-0.752) -- (0.981,-0.742) -- (0.986,-0.732) -- (0.990,-0.723) -- (0.995,-0.713) -- (0.999,-0.703) -- (1.004,-0.694) -- (1.008,-0.684) -- (1.012,-0.674) -- (1.016,-0.665) -- (1.020,-0.656) -- (1.023,-0.646) -- (1.027,-0.637) -- (1.030,-0.628) -- (1.034,-0.619) -- (1.037,-0.609) -- (1.040,-0.600) -- (1.043,-0.591) -- (1.046,-0.582) -- (1.049,-0.574) -- (1.052,-0.565) -- (1.055,-0.556) -- (1.057,-0.547) -- (1.060,-0.539) -- (1.062,-0.530) -- (1.065,-0.522) -- (1.067,-0.513) -- (1.069,-0.505) -- (1.071,-0.497) -- (1.073,-0.489) -- (1.075,-0.480) -- (1.077,-0.472) -- (1.078,-0.464) -- (1.080,-0.457) -- (1.082,-0.449) -- (1.083,-0.441) -- (1.085,-0.433) -- (1.086,-0.426) -- (1.087,-0.418) -- (1.088,-0.411) -- (1.089,-0.403) -- (1.091,-0.396) -- (1.092,-0.389) -- (1.092,-0.382) -- (1.093,-0.374) -- (1.094,-0.367) -- (1.095,-0.361) -- (1.095,-0.354) -- (1.096,-0.347) -- (1.097,-0.340) -- (1.097,-0.334) -- (1.098,-0.327) -- (1.098,-0.321) -- (1.098,-0.314) -- (1.098,-0.308) -- (1.099,-0.302) -- (1.099,-0.295) -- (1.099,-0.289) -- (1.099,-0.283) -- (1.099,-0.277) -- (1.099,-0.272) -- (1.099,-0.266) -- (1.099,-0.260) -- (1.099,-0.254) -- (1.099,-0.249) -- (1.099,-0.243) -- (1.098,-0.238) -- (1.098,-0.233) -- (1.098,-0.227) -- (1.098,-0.222) -- (1.097,-0.217) -- (1.097,-0.212) -- (1.097,-0.207) -- (1.096,-0.202) -- (1.096,-0.197) -- (1.095,-0.192) -- (1.095,-0.188) -- (1.094,-0.183) -- (1.094,-0.179) -- (1.093,-0.174) -- (1.093,-0.170) -- (1.092,-0.165) -- (1.092,-0.161) -- (1.091,-0.157) -- (1.091,-0.152) -- (1.090,-0.148) -- (1.089,-0.144) -- (1.089,-0.140) -- (1.088,-0.136) -- (1.088,-0.132) -- (1.087,-0.128) -- (1.086,-0.125) -- (1.086,-0.121) -- (1.085,-0.117) -- (1.084,-0.114) -- (1.084,-0.110) -- (1.083,-0.107) -- (1.083,-0.103) -- (1.082,-0.100) -- (1.081,-0.096) -- (1.081,-0.093) -- (1.080,-0.090) -- (1.080,-0.086) -- (1.079,-0.083) -- (1.079,-0.080) -- (1.078,-0.077) -- (1.078,-0.074) -- (1.077,-0.071) -- (1.077,-0.068) -- (1.076,-0.065) -- (1.076,-0.062) -- (1.075,-0.059) -- (1.075,-0.056) -- (1.074,-0.053) -- (1.074,-0.050) -- (1.074,-0.047) -- (1.073,-0.045) -- (1.073,-0.042) -- (1.072,-0.039) -- (1.072,-0.037) -- (1.072,-0.034) -- (1.072,-0.031) -- (1.071,-0.028) -- (1.071,-0.026) -- (1.071,-0.023) -- (1.071,-0.021) -- (1.071,-0.018) -- (1.070,-0.015) -- (1.070,-0.013) -- (1.070,-0.010) -- (1.070,-0.008) -- (1.070,-0.005) -- (1.070,-0.003) -- (1.070,-0.000) -- (1.070,0.003);
\draw[knotback] (1.070,0.003) -- (1.070,0.005) -- (1.070,0.008) -- (1.070,0.010) -- (1.070,0.013) -- (1.070,0.015) -- (1.071,0.018) -- (1.071,0.021) -- (1.071,0.023) -- (1.071,0.026) -- (1.071,0.028) -- (1.072,0.031) -- (1.072,0.034) -- (1.072,0.037) -- (1.072,0.039) -- (1.073,0.042) -- (1.073,0.045) -- (1.074,0.047) -- (1.074,0.050) -- (1.074,0.053) -- (1.075,0.056) -- (1.075,0.059) -- (1.076,0.062) -- (1.076,0.065) -- (1.077,0.068) -- (1.077,0.071) -- (1.078,0.074) -- (1.078,0.077) -- (1.079,0.080) -- (1.079,0.083) -- (1.080,0.086) -- (1.080,0.090) -- (1.081,0.093) -- (1.081,0.096) -- (1.082,0.100) -- (1.083,0.103) -- (1.083,0.107) -- (1.084,0.110) -- (1.084,0.114) -- (1.085,0.117) -- (1.086,0.121) -- (1.086,0.125) -- (1.087,0.128) -- (1.088,0.132) -- (1.088,0.136) -- (1.089,0.140) -- (1.089,0.144) -- (1.090,0.148) -- (1.091,0.152) -- (1.091,0.157) -- (1.092,0.161) -- (1.092,0.165) -- (1.093,0.170) -- (1.093,0.174) -- (1.094,0.179) -- (1.094,0.183) -- (1.095,0.188) -- (1.095,0.192) -- (1.096,0.197) -- (1.096,0.202) -- (1.097,0.207) -- (1.097,0.212) -- (1.097,0.217) -- (1.098,0.222) -- (1.098,0.227) -- (1.098,0.233) -- (1.098,0.238) -- (1.099,0.243) -- (1.099,0.249) -- (1.099,0.254) -- (1.099,0.260) -- (1.099,0.266) -- (1.099,0.272) -- (1.099,0.277) -- (1.099,0.283) -- (1.099,0.289) -- (1.099,0.295) -- (1.099,0.302) -- (1.098,0.308) -- (1.098,0.314) -- (1.098,0.321) -- (1.098,0.327) -- (1.097,0.334) -- (1.097,0.340) -- (1.096,0.347) -- (1.095,0.354) -- (1.095,0.361) -- (1.094,0.367) -- (1.093,0.374) -- (1.092,0.382) -- (1.092,0.389) -- (1.091,0.396) -- (1.089,0.403) -- (1.088,0.411) -- (1.087,0.418) -- (1.086,0.426) -- (1.085,0.433) -- (1.083,0.441) -- (1.082,0.449) -- (1.080,0.457) -- (1.078,0.464) -- (1.077,0.472) -- (1.075,0.480) -- (1.073,0.489) -- (1.071,0.497) -- (1.069,0.505) -- (1.067,0.513) -- (1.065,0.522) -- (1.062,0.530) -- (1.060,0.539) -- (1.057,0.547) -- (1.055,0.556) -- (1.052,0.565) -- (1.049,0.574) -- (1.046,0.582) -- (1.043,0.591) -- (1.040,0.600) -- (1.037,0.609) -- (1.034,0.619) -- (1.030,0.628) -- (1.027,0.637) -- (1.023,0.646) -- (1.020,0.656) -- (1.016,0.665) -- (1.012,0.674) -- (1.008,0.684) -- (1.004,0.694) -- (0.999,0.703) -- (0.995,0.713) -- (0.990,0.723) -- (0.986,0.732) -- (0.981,0.742) -- (0.976,0.752) -- (0.971,0.762) -- (0.966,0.772) -- (0.961,0.782) -- (0.956,0.792) -- (0.950,0.802) -- (0.945,0.812) -- (0.939,0.822) -- (0.933,0.832) -- (0.927,0.842) -- (0.921,0.853) -- (0.915,0.863) -- (0.908,0.873) -- (0.902,0.883) -- (0.895,0.894) -- (0.889,0.904) -- (0.882,0.915) -- (0.875,0.925) -- (0.868,0.935) -- (0.861,0.946) -- (0.853,0.956) -- (0.846,0.967) -- (0.838,0.977) -- (0.831,0.988) -- (0.823,0.998) -- (0.815,1.009) -- (0.807,1.019) -- (0.798,1.030) -- (0.790,1.040) -- (0.781,1.051) -- (0.773,1.061) -- (0.764,1.072) -- (0.755,1.082) -- (0.746,1.093) -- (0.737,1.103) -- (0.728,1.114) -- (0.718,1.124) -- (0.709,1.135) -- (0.699,1.145) -- (0.689,1.156) -- (0.679,1.166) -- (0.669,1.176) -- (0.659,1.187) -- (0.648,1.197) -- (0.638,1.207) -- (0.627,1.218) -- (0.617,1.228) -- (0.606,1.238) -- (0.595,1.248) -- (0.584,1.259) -- (0.572,1.269) -- (0.561,1.279) -- (0.550,1.289) -- (0.538,1.299) -- (0.526,1.309) -- (0.514,1.319) -- (0.502,1.329) -- (0.490,1.338) -- (0.478,1.348) -- (0.466,1.358) -- (0.453,1.368) -- (0.441,1.377) -- (0.428,1.387) -- (0.415,1.396) -- (0.402,1.406) -- (0.389,1.415) -- (0.376,1.424) -- (0.363,1.433) -- (0.350,1.443) -- (0.336,1.452) -- (0.323,1.461) -- (0.309,1.470) -- (0.295,1.478) -- (0.281,1.487) -- (0.267,1.496) -- (0.253,1.505) -- (0.239,1.513) -- (0.225,1.522) -- (0.210,1.530) -- (0.196,1.538) -- (0.181,1.546) -- (0.167,1.554) -- (0.152,1.562) -- (0.137,1.570) -- (0.122,1.578) -- (0.107,1.586) -- (0.092,1.594) -- (0.077,1.601) -- (0.062,1.608) -- (0.046,1.616) -- (0.031,1.623) -- (0.016,1.630) -- (0.000,1.637) -- (-0.016,1.644) -- (-0.031,1.651) -- (-0.047,1.657) -- (-0.063,1.664) -- (-0.079,1.670) -- (-0.095,1.677) -- (-0.111,1.683) -- (-0.127,1.689) -- (-0.143,1.695) -- (-0.159,1.701) -- (-0.175,1.706) -- (-0.191,1.712) -- (-0.208,1.718) -- (-0.224,1.723) -- (-0.240,1.728) -- (-0.257,1.733) -- (-0.273,1.738) -- (-0.290,1.743) -- (-0.306,1.748) -- (-0.323,1.752) -- (-0.340,1.757) -- (-0.356,1.761) -- (-0.373,1.765) -- (-0.389,1.769) -- (-0.406,1.773) -- (-0.423,1.777) -- (-0.440,1.781) -- (-0.456,1.784) -- (-0.473,1.787) -- (-0.490,1.791) -- (-0.507,1.794) -- (-0.523,1.797) -- (-0.540,1.799) -- (-0.557,1.802) -- (-0.574,1.805) -- (-0.590,1.807) -- (-0.607,1.809) -- (-0.624,1.811) -- (-0.641,1.813) -- (-0.657,1.815) -- (-0.674,1.817) -- (-0.691,1.818) -- (-0.707,1.819) -- (-0.724,1.821) -- (-0.741,1.822) -- (-0.757,1.823) -- (-0.774,1.823) -- (-0.790,1.824) -- (-0.807,1.824) -- (-0.823,1.825) -- (-0.840,1.825) -- (-0.856,1.825) -- (-0.872,1.825) -- (-0.889,1.825) -- (-0.905,1.824) -- (-0.921,1.824) -- (-0.937,1.823) -- (-0.953,1.822) -- (-0.969,1.821) -- (-0.985,1.820) -- (-1.001,1.818) -- (-1.017,1.817) -- (-1.032,1.815) -- (-1.048,1.814) -- (-1.064,1.812) -- (-1.079,1.810) -- (-1.095,1.807) -- (-1.110,1.805) -- (-1.125,1.803) -- (-1.140,1.800) -- (-1.155,1.797) -- (-1.170,1.794) -- (-1.185,1.791) -- (-1.200,1.788) -- (-1.215,1.785) -- (-1.230,1.781) -- (-1.244,1.778) -- (-1.259,1.774) -- (-1.273,1.770) -- (-1.287,1.766) -- (-1.301,1.762) -- (-1.315,1.757) -- (-1.329,1.753) -- (-1.343,1.748) -- (-1.357,1.743) -- (-1.370,1.739) -- (-1.384,1.734) -- (-1.397,1.728) -- (-1.410,1.723) -- (-1.423,1.718) -- (-1.436,1.712) -- (-1.449,1.707) -- (-1.462,1.701) -- (-1.474,1.695) -- (-1.487,1.689) -- (-1.499,1.683) -- (-1.511,1.676) -- (-1.523,1.670) -- (-1.535,1.663) -- (-1.547,1.657) -- (-1.559,1.650) -- (-1.570,1.643) -- (-1.581,1.636) -- (-1.593,1.629) -- (-1.604,1.622) -- (-1.615,1.614) -- (-1.625,1.607) -- (-1.636,1.599) -- (-1.646,1.592) -- (-1.657,1.584) -- (-1.667,1.576) -- (-1.677,1.568) -- (-1.687,1.560) -- (-1.696,1.552) -- (-1.706,1.543);
\draw[knot] (-1.706,1.543) -- (-1.715,1.535) -- (-1.724,1.527) -- (-1.733,1.518) -- (-1.742,1.509) -- (-1.751,1.500) -- (-1.759,1.492) -- (-1.768,1.483) -- (-1.776,1.474) -- (-1.784,1.464) -- (-1.792,1.455) -- (-1.799,1.446) -- (-1.807,1.436) -- (-1.814,1.427) -- (-1.821,1.417) -- (-1.828,1.408) -- (-1.835,1.398) -- (-1.842,1.388) -- (-1.848,1.378) -- (-1.855,1.368) -- (-1.861,1.358) -- (-1.867,1.348) -- (-1.873,1.338) -- (-1.878,1.328) -- (-1.884,1.318) -- (-1.889,1.307) -- (-1.894,1.297) -- (-1.899,1.286) -- (-1.904,1.276) -- (-1.908,1.265) -- (-1.913,1.255) -- (-1.917,1.244) -- (-1.921,1.233) -- (-1.925,1.222) -- (-1.928,1.212) -- (-1.932,1.201) -- (-1.935,1.190) -- (-1.938,1.179) -- (-1.941,1.168) -- (-1.944,1.157) -- (-1.947,1.146) -- (-1.949,1.135) -- (-1.952,1.123) -- (-1.954,1.112) -- (-1.956,1.101) -- (-1.958,1.090) -- (-1.959,1.078) -- (-1.961,1.067) -- (-1.962,1.056) -- (-1.963,1.044) -- (-1.964,1.033) -- (-1.965,1.022) -- (-1.965,1.010) -- (-1.966,0.999) -- (-1.966,0.987) -- (-1.966,0.976) -- (-1.966,0.964) -- (-1.966,0.953) -- (-1.965,0.941) -- (-1.965,0.930) -- (-1.964,0.918) -- (-1.963,0.907) -- (-1.962,0.895) -- (-1.961,0.884) -- (-1.959,0.872) -- (-1.958,0.861) -- (-1.956,0.849) -- (-1.954,0.837) -- (-1.952,0.826) -- (-1.950,0.814) -- (-1.948,0.803) -- (-1.946,0.791) -- (-1.943,0.780) -- (-1.940,0.768) -- (-1.937,0.757) -- (-1.934,0.745) -- (-1.931,0.734) -- (-1.928,0.723) -- (-1.925,0.711) -- (-1.921,0.700) -- (-1.917,0.688) -- (-1.913,0.677) -- (-1.909,0.666) -- (-1.905,0.654) -- (-1.901,0.643) -- (-1.897,0.632) -- (-1.892,0.620) -- (-1.888,0.609) -- (-1.883,0.598) -- (-1.878,0.587) -- (-1.873,0.576) -- (-1.868,0.565) -- (-1.863,0.554) -- (-1.857,0.542) -- (-1.852,0.531) -- (-1.846,0.520) -- (-1.841,0.510) -- (-1.835,0.499) -- (-1.829,0.488) -- (-1.823,0.477) -- (-1.817,0.466) -- (-1.810,0.455) -- (-1.804,0.445) -- (-1.798,0.434) -- (-1.791,0.423) -- (-1.785,0.413) -- (-1.778,0.402) -- (-1.771,0.392) -- (-1.764,0.381) -- (-1.757,0.371) -- (-1.750,0.360) -- (-1.743,0.350) -- (-1.736,0.340) -- (-1.728,0.329) -- (-1.721,0.319) -- (-1.713,0.309) -- (-1.706,0.299) -- (-1.698,0.289) -- (-1.690,0.279) -- (-1.683,0.269) -- (-1.675,0.259) -- (-1.667,0.249) -- (-1.659,0.239) -- (-1.651,0.230) -- (-1.643,0.220) -- (-1.635,0.210) -- (-1.626,0.201) -- (-1.618,0.191) -- (-1.610,0.182) -- (-1.601,0.172) -- (-1.593,0.163) -- (-1.584,0.153) -- (-1.576,0.144) -- (-1.567,0.135) -- (-1.559,0.125) -- (-1.550,0.116) -- (-1.541,0.107) -- (-1.533,0.098) -- (-1.524,0.089) -- (-1.515,0.080) -- (-1.506,0.071) -- (-1.497,0.062) -- (-1.488,0.054) -- (-1.480,0.045) -- (-1.471,0.036) -- (-1.462,0.027) -- (-1.453,0.019) -- (-1.444,0.010) -- (-1.434,0.002) -- (-1.425,-0.007) -- (-1.416,-0.015) -- (-1.407,-0.024) -- (-1.398,-0.032) -- (-1.389,-0.040) -- (-1.380,-0.048) -- (-1.371,-0.057) -- (-1.362,-0.065) -- (-1.352,-0.073) -- (-1.343,-0.081) -- (-1.334,-0.089) -- (-1.325,-0.097) -- (-1.316,-0.105) -- (-1.307,-0.113) -- (-1.297,-0.121) -- (-1.288,-0.128) -- (-1.279,-0.136) -- (-1.270,-0.144) -- (-1.261,-0.151) -- (-1.252,-0.159) -- (-1.243,-0.167) -- (-1.234,-0.174) -- (-1.224,-0.182) -- (-1.215,-0.189) -- (-1.206,-0.197) -- (-1.197,-0.204) -- (-1.188,-0.211) -- (-1.179,-0.219) -- (-1.170,-0.226) -- (-1.161,-0.233) -- (-1.152,-0.241) -- (-1.143,-0.248) -- (-1.135,-0.255) -- (-1.126,-0.262) -- (-1.117,-0.269) -- (-1.108,-0.276) -- (-1.099,-0.283) -- (-1.090,-0.290) -- (-1.082,-0.297) -- (-1.073,-0.304) -- (-1.064,-0.311) -- (-1.056,-0.318) -- (-1.047,-0.325) -- (-1.038,-0.332) -- (-1.030,-0.339) -- (-1.021,-0.346) -- (-1.013,-0.353) -- (-1.004,-0.359) -- (-0.996,-0.366) -- (-0.987,-0.373) -- (-0.979,-0.380) -- (-0.971,-0.386) -- (-0.962,-0.393) -- (-0.954,-0.400) -- (-0.946,-0.406) -- (-0.938,-0.413) -- (-0.930,-0.420) -- (-0.922,-0.426) -- (-0.913,-0.433) -- (-0.905,-0.440) -- (-0.897,-0.446) -- (-0.889,-0.453) -- (-0.882,-0.459) -- (-0.874,-0.466) -- (-0.866,-0.473) -- (-0.858,-0.479) -- (-0.850,-0.486) -- (-0.843,-0.492) -- (-0.835,-0.499) -- (-0.827,-0.505) -- (-0.820,-0.512) -- (-0.812,-0.518) -- (-0.804,-0.525);
\draw[knotback] (-0.804,-0.525) -- (-0.797,-0.532) -- (-0.789,-0.538) -- (-0.782,-0.545) -- (-0.775,-0.551) -- (-0.767,-0.558) -- (-0.760,-0.564) -- (-0.753,-0.571) -- (-0.745,-0.577) -- (-0.738,-0.584) -- (-0.731,-0.591) -- (-0.724,-0.597) -- (-0.717,-0.604) -- (-0.710,-0.610) -- (-0.703,-0.617) -- (-0.696,-0.624) -- (-0.689,-0.630) -- (-0.682,-0.637) -- (-0.675,-0.643) -- (-0.668,-0.650) -- (-0.661,-0.657) -- (-0.654,-0.663) -- (-0.647,-0.670) -- (-0.641,-0.677) -- (-0.634,-0.683) -- (-0.627,-0.690) -- (-0.620,-0.697) -- (-0.614,-0.704) -- (-0.607,-0.710) -- (-0.600,-0.717) -- (-0.594,-0.724) -- (-0.587,-0.731) -- (-0.581,-0.738) -- (-0.574,-0.744) -- (-0.568,-0.751) -- (-0.561,-0.758) -- (-0.554,-0.765) -- (-0.548,-0.772) -- (-0.541,-0.779) -- (-0.535,-0.786) -- (-0.529,-0.793) -- (-0.522,-0.800) -- (-0.516,-0.807) -- (-0.509,-0.814) -- (-0.503,-0.821) -- (-0.496,-0.828) -- (-0.490,-0.835) -- (-0.484,-0.842) -- (-0.477,-0.849) -- (-0.471,-0.856) -- (-0.464,-0.863) -- (-0.458,-0.870) -- (-0.451,-0.877) -- (-0.445,-0.885) -- (-0.439,-0.892) -- (-0.432,-0.899) -- (-0.426,-0.906) -- (-0.419,-0.913) -- (-0.413,-0.921) -- (-0.406,-0.928) -- (-0.400,-0.935) -- (-0.393,-0.942) -- (-0.387,-0.950) -- (-0.380,-0.957) -- (-0.374,-0.964) -- (-0.367,-0.972) -- (-0.361,-0.979) -- (-0.354,-0.986) -- (-0.348,-0.994) -- (-0.341,-1.001) -- (-0.334,-1.009) -- (-0.328,-1.016) -- (-0.321,-1.023) -- (-0.314,-1.031) -- (-0.307,-1.038) -- (-0.301,-1.046) -- (-0.294,-1.053) -- (-0.287,-1.061) -- (-0.280,-1.068) -- (-0.273,-1.075) -- (-0.266,-1.083) -- (-0.259,-1.090) -- (-0.252,-1.098) -- (-0.245,-1.105) -- (-0.238,-1.113) -- (-0.231,-1.120) -- (-0.224,-1.128) -- (-0.216,-1.135) -- (-0.209,-1.142) -- (-0.202,-1.150) -- (-0.194,-1.157) -- (-0.187,-1.165) -- (-0.179,-1.172) -- (-0.172,-1.180) -- (-0.164,-1.187) -- (-0.157,-1.194) -- (-0.149,-1.202) -- (-0.141,-1.209) -- (-0.133,-1.216) -- (-0.125,-1.224) -- (-0.117,-1.231) -- (-0.110,-1.238) -- (-0.101,-1.246) -- (-0.093,-1.253) -- (-0.085,-1.260) -- (-0.077,-1.267) -- (-0.069,-1.275) -- (-0.060,-1.282) -- (-0.052,-1.289) -- (-0.043,-1.296) -- (-0.035,-1.303) -- (-0.026,-1.310) -- (-0.018,-1.317) -- (-0.009,-1.324) -- (-0.000,-1.331) -- (0.009,-1.338) -- (0.018,-1.345) -- (0.027,-1.352) -- (0.036,-1.359) -- (0.045,-1.365) -- (0.054,-1.372) -- (0.064,-1.379) -- (0.073,-1.385) -- (0.083,-1.392) -- (0.092,-1.398) -- (0.102,-1.405) -- (0.112,-1.411) -- (0.121,-1.418) -- (0.131,-1.424) -- (0.141,-1.430) -- (0.151,-1.437) -- (0.161,-1.443) -- (0.171,-1.449) -- (0.182,-1.455) -- (0.192,-1.461) -- (0.202,-1.467) -- (0.213,-1.473) -- (0.223,-1.478) -- (0.234,-1.484) -- (0.245,-1.490) -- (0.256,-1.495) -- (0.267,-1.501) -- (0.277,-1.506) -- (0.289,-1.511) -- (0.300,-1.517) -- (0.311,-1.522) -- (0.322,-1.527) -- (0.333,-1.532) -- (0.345,-1.537) -- (0.356,-1.542) -- (0.368,-1.546) -- (0.380,-1.551) -- (0.391,-1.556) -- (0.403,-1.560) -- (0.415,-1.564) -- (0.427,-1.569) -- (0.439,-1.573) -- (0.451,-1.577) -- (0.464,-1.581) -- (0.476,-1.585) -- (0.488,-1.588) -- (0.501,-1.592) -- (0.513,-1.596) -- (0.526,-1.599) -- (0.538,-1.602) -- (0.551,-1.606) -- (0.564,-1.609) -- (0.577,-1.612) -- (0.590,-1.615) -- (0.603,-1.617) -- (0.616,-1.620) -- (0.629,-1.622) -- (0.642,-1.625) -- (0.655,-1.627) -- (0.669,-1.629) -- (0.682,-1.631) -- (0.695,-1.633) -- (0.709,-1.635) -- (0.722,-1.637) -- (0.736,-1.638) -- (0.750,-1.639) -- (0.763,-1.641) -- (0.777,-1.642) -- (0.791,-1.643) -- (0.805,-1.644) -- (0.819,-1.644) -- (0.833,-1.645) -- (0.847,-1.645) -- (0.861,-1.645) -- (0.875,-1.646) -- (0.889,-1.646) -- (0.903,-1.645) -- (0.918,-1.645) -- (0.932,-1.645) -- (0.946,-1.644) -- (0.961,-1.643) -- (0.975,-1.642) -- (0.989,-1.641) -- (1.004,-1.640) -- (1.018,-1.639) -- (1.033,-1.637) -- (1.047,-1.636) -- (1.062,-1.634) -- (1.077,-1.632) -- (1.091,-1.630) -- (1.106,-1.627) -- (1.120,-1.625) -- (1.135,-1.622) -- (1.150,-1.619) -- (1.164,-1.616) -- (1.179,-1.613) -- (1.194,-1.610) -- (1.209,-1.607) -- (1.223,-1.603) -- (1.238,-1.599) -- (1.253,-1.595) -- (1.267,-1.591) -- (1.282,-1.587) -- (1.297,-1.583) -- (1.312,-1.578) -- (1.326,-1.573) -- (1.341,-1.568) -- (1.356,-1.563) -- (1.370,-1.558) -- (1.385,-1.553) -- (1.400,-1.547) -- (1.414,-1.541) -- (1.429,-1.535) -- (1.443,-1.529) -- (1.458,-1.523) -- (1.472,-1.517) -- (1.487,-1.510) -- (1.501,-1.503) -- (1.516,-1.496) -- (1.530,-1.489) -- (1.544,-1.482) -- (1.559,-1.475) -- (1.573,-1.467) -- (1.587,-1.459) -- (1.601,-1.451) -- (1.615,-1.443) -- (1.629,-1.435) -- (1.643,-1.427) -- (1.657,-1.418) -- (1.671,-1.409) -- (1.685,-1.400) -- (1.698,-1.391) -- (1.712,-1.382) -- (1.726,-1.373) -- (1.739,-1.363) -- (1.753,-1.353) -- (1.766,-1.343) -- (1.779,-1.333) -- (1.792,-1.323) -- (1.806,-1.313) -- (1.819,-1.302) -- (1.832,-1.292) -- (1.844,-1.281) -- (1.857,-1.270) -- (1.870,-1.259) -- (1.883,-1.247) -- (1.895,-1.236) -- (1.907,-1.224) -- (1.920,-1.213) -- (1.932,-1.201) -- (1.944,-1.189) -- (1.956,-1.176) -- (1.968,-1.164) -- (1.980,-1.152) -- (1.991,-1.139) -- (2.003,-1.126) -- (2.014,-1.113) -- (2.026,-1.100) -- (2.037,-1.087) -- (2.048,-1.074) -- (2.059,-1.060) -- (2.070,-1.047) -- (2.080,-1.033) -- (2.091,-1.019) -- (2.101,-1.005) -- (2.112,-0.991) -- (2.122,-0.977) -- (2.132,-0.962) -- (2.142,-0.948) -- (2.151,-0.933) -- (2.161,-0.919) -- (2.170,-0.904) -- (2.180,-0.889) -- (2.189,-0.874) -- (2.198,-0.859) -- (2.207,-0.843) -- (2.215,-0.828) -- (2.224,-0.812) -- (2.232,-0.797) -- (2.241,-0.781) -- (2.249,-0.765) -- (2.257,-0.749) -- (2.264,-0.733) -- (2.272,-0.717) -- (2.279,-0.701) -- (2.287,-0.685) -- (2.294,-0.668) -- (2.301,-0.652) -- (2.307,-0.635) -- (2.314,-0.618) -- (2.320,-0.602) -- (2.327,-0.585) -- (2.333,-0.568) -- (2.339,-0.551) -- (2.344,-0.534) -- (2.350,-0.517) -- (2.355,-0.500) -- (2.360,-0.482) -- (2.365,-0.465) -- (2.370,-0.448) -- (2.375,-0.430) -- (2.379,-0.413) -- (2.384,-0.395) -- (2.388,-0.378) -- (2.392,-0.360) -- (2.395,-0.342) -- (2.399,-0.325) -- (2.402,-0.307) -- (2.405,-0.289) -- (2.408,-0.271) -- (2.411,-0.253) -- (2.414,-0.235) -- (2.416,-0.217) -- (2.418,-0.199) -- (2.420,-0.181) -- (2.422,-0.163) -- (2.424,-0.145) -- (2.425,-0.127) -- (2.427,-0.109) -- (2.428,-0.091) -- (2.428,-0.073) -- (2.429,-0.054) -- (2.430,-0.036) -- (2.430,-0.018) -- (2.430,-0.000);

  \end{scope}
  \node[lab, anchor=north] at (9.65,-4.95) {the (2, 3) torus knot, a trefoil};
  \begin{axis}[snd, at={(12.85cm,-4.95cm)}, xmin=0, xmax=5.8, ymax=1.75, xtick={0,1,...,5}, xlabel={frequency ÷ f₁ \quad (f₂ = 1.5 f₁)}]
    \grid
    \addplot[ycomb, line width=1.4pt, ink] table[x=freq, y=amp] {tkco-closed.dat};
    \node[lab, anchor=north east, align=right] at (axis cs:5.8,1.78) {dashed: every f₀ = f₁ ÷ 2};
  \end{axis}
\end{scope}

% =============== OPEN · GOLDEN ===============
\def\y{-6.3}
\rowframe{\y}{OPEN \enspace f₂ ÷ f₁ = φ = 1.618…}{the path never comes home: twelve turns, and the surface starts to fill}
\begin{scope}[shift={(0,\y)}]
  \twophasors{phasor f₁}{phasor f₂ = φf₁}
  \twocircles
  \begin{scope}[shift={(10.45,-3.25)}, scale=0.75]
    % torus mesh, generated by torus-knot-closed-and-open.py
\draw[mesh] (2.428,0.028) -- (2.422,0.056) -- (2.411,0.084) -- (2.397,0.111) -- (2.378,0.138) -- (2.356,0.164) -- (2.330,0.188) -- (2.300,0.212) -- (2.267,0.234) -- (2.231,0.255) -- (2.192,0.274) -- (2.150,0.292) -- (2.105,0.307) -- (2.059,0.321) -- (2.010,0.333) -- (1.960,0.343) -- (1.909,0.350) -- (1.856,0.356) -- (1.803,0.359) -- (1.750,0.360) -- (1.697,0.359) -- (1.644,0.356) -- (1.591,0.350) -- (1.540,0.343) -- (1.490,0.333) -- (1.441,0.321) -- (1.395,0.307) -- (1.350,0.292) -- (1.308,0.274) -- (1.269,0.255) -- (1.233,0.234) -- (1.200,0.212) -- (1.170,0.188) -- (1.144,0.164) -- (1.122,0.138) -- (1.103,0.111) -- (1.089,0.084) -- (1.078,0.056) -- (1.072,0.028) -- (1.070,0.000) -- (1.072,-0.028);
\draw[mesh] (2.329,0.613) -- (2.315,0.637) -- (2.297,0.660) -- (2.276,0.681) -- (2.250,0.700) -- (2.222,0.717) -- (2.190,0.732) -- (2.155,0.744) -- (2.117,0.755) -- (2.076,0.763) -- (2.034,0.769) -- (1.989,0.773) -- (1.942,0.774) -- (1.893,0.773) -- (1.844,0.769) -- (1.793,0.763) -- (1.742,0.755) -- (1.690,0.744) -- (1.639,0.732) -- (1.588,0.717) -- (1.537,0.700) -- (1.487,0.681) -- (1.439,0.660) -- (1.392,0.637) -- (1.347,0.613) -- (1.304,0.588) -- (1.264,0.561) -- (1.226,0.533) -- (1.191,0.505) -- (1.159,0.475) -- (1.130,0.445) -- (1.105,0.415) -- (1.084,0.384) -- (1.066,0.354) -- (1.052,0.323) -- (1.042,0.293) -- (1.036,0.264) -- (1.034,0.235) -- (1.036,0.207) -- (1.042,0.180) -- (1.052,0.155);
\draw[mesh] (2.076,1.128) -- (2.060,1.146) -- (2.040,1.163) -- (2.018,1.176) -- (1.992,1.187) -- (1.963,1.195) -- (1.932,1.201) -- (1.898,1.203) -- (1.862,1.203) -- (1.823,1.200) -- (1.783,1.194) -- (1.741,1.185) -- (1.698,1.174) -- (1.653,1.160) -- (1.608,1.143) -- (1.562,1.124) -- (1.516,1.102) -- (1.469,1.079) -- (1.423,1.053) -- (1.378,1.025) -- (1.334,0.996) -- (1.290,0.965) -- (1.248,0.932) -- (1.208,0.899) -- (1.169,0.864) -- (1.133,0.829) -- (1.099,0.793) -- (1.068,0.757) -- (1.039,0.721) -- (1.013,0.684) -- (0.991,0.649) -- (0.971,0.614) -- (0.955,0.579) -- (0.943,0.546) -- (0.934,0.514) -- (0.928,0.483) -- (0.927,0.454) -- (0.928,0.426) -- (0.934,0.401) -- (0.943,0.378) -- (0.955,0.356);
\draw[mesh] (1.666,1.576) -- (1.647,1.585) -- (1.626,1.591) -- (1.603,1.594) -- (1.577,1.593) -- (1.550,1.588) -- (1.520,1.581) -- (1.489,1.570) -- (1.456,1.556) -- (1.421,1.538) -- (1.386,1.518) -- (1.350,1.495) -- (1.313,1.469) -- (1.275,1.441) -- (1.237,1.410) -- (1.200,1.377) -- (1.162,1.342) -- (1.125,1.305) -- (1.089,1.266) -- (1.053,1.226) -- (1.019,1.185) -- (0.986,1.144) -- (0.955,1.101) -- (0.925,1.059) -- (0.897,1.016) -- (0.872,0.973) -- (0.848,0.931) -- (0.827,0.890) -- (0.809,0.850) -- (0.793,0.811) -- (0.780,0.773) -- (0.770,0.737) -- (0.763,0.703) -- (0.758,0.671) -- (0.757,0.642) -- (0.758,0.615) -- (0.763,0.590) -- (0.770,0.569) -- (0.780,0.550) -- (0.793,0.535) -- (0.809,0.522);
\draw[mesh] (1.165,1.899) -- (1.150,1.901) -- (1.134,1.899) -- (1.115,1.893) -- (1.096,1.884) -- (1.075,1.870) -- (1.053,1.853) -- (1.029,1.833) -- (1.005,1.809) -- (0.980,1.782) -- (0.954,1.752) -- (0.928,1.719) -- (0.902,1.684) -- (0.875,1.646) -- (0.848,1.605) -- (0.822,1.563) -- (0.796,1.519) -- (0.770,1.474) -- (0.745,1.427) -- (0.721,1.380) -- (0.697,1.332) -- (0.675,1.283) -- (0.654,1.235) -- (0.635,1.187) -- (0.616,1.140) -- (0.600,1.093) -- (0.585,1.048) -- (0.572,1.004) -- (0.561,0.962) -- (0.552,0.922) -- (0.544,0.884) -- (0.539,0.848) -- (0.536,0.816) -- (0.535,0.786) -- (0.536,0.759) -- (0.539,0.736) -- (0.544,0.716) -- (0.552,0.699) -- (0.561,0.686) -- (0.572,0.677) -- (0.585,0.671);
\draw[mesh] (0.603,2.097) -- (0.595,2.096) -- (0.587,2.091) -- (0.577,2.082) -- (0.567,2.069) -- (0.556,2.052) -- (0.545,2.032) -- (0.533,2.007) -- (0.520,1.980) -- (0.507,1.948) -- (0.494,1.914) -- (0.480,1.877) -- (0.467,1.836) -- (0.453,1.794) -- (0.439,1.749) -- (0.425,1.702) -- (0.412,1.654) -- (0.399,1.604) -- (0.386,1.553) -- (0.373,1.502) -- (0.361,1.450) -- (0.349,1.398) -- (0.339,1.346) -- (0.328,1.294) -- (0.319,1.244) -- (0.311,1.195) -- (0.303,1.147) -- (0.296,1.101) -- (0.290,1.057) -- (0.286,1.015) -- (0.282,0.976) -- (0.279,0.940) -- (0.277,0.906) -- (0.277,0.876) -- (0.277,0.850) -- (0.279,0.827) -- (0.282,0.808) -- (0.286,0.792) -- (0.290,0.781) -- (0.296,0.774) -- (0.303,0.770);
\draw[mesh] (0.000,2.162) -- (0.000,2.157) -- (0.000,2.147) -- (0.000,2.133) -- (0.000,2.115) -- (0.000,2.093) -- (0.000,2.067) -- (0.000,2.038) -- (0.000,2.005) -- (0.000,1.969) -- (0.000,1.930) -- (0.000,1.889) -- (0.000,1.844) -- (0.000,1.798) -- (0.000,1.750) -- (0.000,1.700) -- (0.000,1.649) -- (0.000,1.596) -- (0.000,1.543) -- (0.000,1.490) -- (0.000,1.437) -- (0.000,1.384) -- (0.000,1.331) -- (0.000,1.280) -- (0.000,1.229) -- (0.000,1.181) -- (0.000,1.134) -- (0.000,1.089) -- (0.000,1.047) -- (0.000,1.007) -- (0.000,0.971) -- (0.000,0.937) -- (0.000,0.907) -- (0.000,0.881) -- (0.000,0.858) -- (0.000,0.839) -- (0.000,0.824) -- (0.000,0.813) -- (0.000,0.807) -- (0.000,0.804) -- (0.000,0.806);
\draw[mesh] (-0.603,2.097) -- (-0.595,2.096) -- (-0.587,2.091) -- (-0.577,2.082) -- (-0.567,2.069) -- (-0.556,2.052) -- (-0.545,2.032) -- (-0.533,2.007) -- (-0.520,1.980) -- (-0.507,1.948) -- (-0.494,1.914) -- (-0.480,1.877) -- (-0.467,1.836) -- (-0.453,1.794) -- (-0.439,1.749) -- (-0.425,1.702) -- (-0.412,1.654) -- (-0.399,1.604) -- (-0.386,1.553) -- (-0.373,1.502) -- (-0.361,1.450) -- (-0.349,1.398) -- (-0.339,1.346) -- (-0.328,1.294) -- (-0.319,1.244) -- (-0.311,1.195) -- (-0.303,1.147) -- (-0.296,1.101) -- (-0.290,1.057) -- (-0.286,1.015) -- (-0.282,0.976) -- (-0.279,0.940) -- (-0.277,0.906) -- (-0.277,0.876) -- (-0.277,0.850) -- (-0.279,0.827) -- (-0.282,0.808) -- (-0.286,0.792) -- (-0.290,0.781) -- (-0.296,0.774) -- (-0.303,0.770);
\draw[mesh] (-1.165,1.899) -- (-1.150,1.901) -- (-1.134,1.899) -- (-1.115,1.893) -- (-1.096,1.884) -- (-1.075,1.870) -- (-1.053,1.853) -- (-1.029,1.833) -- (-1.005,1.809) -- (-0.980,1.782) -- (-0.954,1.752) -- (-0.928,1.719) -- (-0.902,1.684) -- (-0.875,1.646) -- (-0.848,1.605) -- (-0.822,1.563) -- (-0.796,1.519) -- (-0.770,1.474) -- (-0.745,1.427) -- (-0.721,1.380) -- (-0.697,1.332) -- (-0.675,1.283) -- (-0.654,1.235) -- (-0.635,1.187) -- (-0.616,1.140) -- (-0.600,1.093) -- (-0.585,1.048) -- (-0.572,1.004) -- (-0.561,0.962) -- (-0.552,0.922) -- (-0.544,0.884) -- (-0.539,0.848) -- (-0.536,0.816) -- (-0.535,0.786) -- (-0.536,0.759) -- (-0.539,0.736) -- (-0.544,0.716) -- (-0.552,0.699) -- (-0.561,0.686) -- (-0.572,0.677) -- (-0.585,0.671);
\draw[mesh] (-1.666,1.576) -- (-1.647,1.585) -- (-1.626,1.591) -- (-1.603,1.594) -- (-1.577,1.593) -- (-1.550,1.588) -- (-1.520,1.581) -- (-1.489,1.570) -- (-1.456,1.556) -- (-1.421,1.538) -- (-1.386,1.518) -- (-1.350,1.495) -- (-1.313,1.469) -- (-1.275,1.441) -- (-1.237,1.410) -- (-1.200,1.377) -- (-1.162,1.342) -- (-1.125,1.305) -- (-1.089,1.266) -- (-1.053,1.226) -- (-1.019,1.185) -- (-0.986,1.144) -- (-0.955,1.101) -- (-0.925,1.059) -- (-0.897,1.016) -- (-0.872,0.973) -- (-0.848,0.931) -- (-0.827,0.890) -- (-0.809,0.850) -- (-0.793,0.811) -- (-0.780,0.773) -- (-0.770,0.737) -- (-0.763,0.703) -- (-0.758,0.671) -- (-0.757,0.642) -- (-0.758,0.615) -- (-0.763,0.590) -- (-0.770,0.569) -- (-0.780,0.550) -- (-0.793,0.535) -- (-0.809,0.522);
\draw[mesh] (-2.076,1.128) -- (-2.060,1.146) -- (-2.040,1.163) -- (-2.018,1.176) -- (-1.992,1.187) -- (-1.963,1.195) -- (-1.932,1.201) -- (-1.898,1.203) -- (-1.862,1.203) -- (-1.823,1.200) -- (-1.783,1.194) -- (-1.741,1.185) -- (-1.698,1.174) -- (-1.653,1.160) -- (-1.608,1.143) -- (-1.562,1.124) -- (-1.516,1.102) -- (-1.469,1.079) -- (-1.423,1.053) -- (-1.378,1.025) -- (-1.334,0.996) -- (-1.290,0.965) -- (-1.248,0.932) -- (-1.208,0.899) -- (-1.169,0.864) -- (-1.133,0.829) -- (-1.099,0.793) -- (-1.068,0.757) -- (-1.039,0.721) -- (-1.013,0.684) -- (-0.991,0.649) -- (-0.971,0.614) -- (-0.955,0.579) -- (-0.943,0.546) -- (-0.934,0.514) -- (-0.928,0.483) -- (-0.927,0.454) -- (-0.928,0.426) -- (-0.934,0.401) -- (-0.943,0.378) -- (-0.955,0.356);
\draw[mesh] (-2.329,0.613) -- (-2.315,0.637) -- (-2.297,0.660) -- (-2.276,0.681) -- (-2.250,0.700) -- (-2.222,0.717) -- (-2.190,0.732) -- (-2.155,0.744) -- (-2.117,0.755) -- (-2.076,0.763) -- (-2.034,0.769) -- (-1.989,0.773) -- (-1.942,0.774) -- (-1.893,0.773) -- (-1.844,0.769) -- (-1.793,0.763) -- (-1.742,0.755) -- (-1.690,0.744) -- (-1.639,0.732) -- (-1.588,0.717) -- (-1.537,0.700) -- (-1.487,0.681) -- (-1.439,0.660) -- (-1.392,0.637) -- (-1.347,0.613) -- (-1.304,0.588) -- (-1.264,0.561) -- (-1.226,0.533) -- (-1.191,0.505) -- (-1.159,0.475) -- (-1.130,0.445) -- (-1.105,0.415) -- (-1.084,0.384) -- (-1.066,0.354) -- (-1.052,0.323) -- (-1.042,0.293) -- (-1.036,0.264) -- (-1.034,0.235) -- (-1.036,0.207) -- (-1.042,0.180) -- (-1.052,0.155);
\draw[mesh] (-2.428,0.028) -- (-2.422,0.056) -- (-2.411,0.084) -- (-2.397,0.111) -- (-2.378,0.138) -- (-2.356,0.164) -- (-2.330,0.188) -- (-2.300,0.212) -- (-2.267,0.234) -- (-2.231,0.255) -- (-2.192,0.274) -- (-2.150,0.292) -- (-2.105,0.307) -- (-2.059,0.321) -- (-2.010,0.333) -- (-1.960,0.343) -- (-1.909,0.350) -- (-1.856,0.356) -- (-1.803,0.359) -- (-1.750,0.360) -- (-1.697,0.359) -- (-1.644,0.356) -- (-1.591,0.350) -- (-1.540,0.343) -- (-1.490,0.333) -- (-1.441,0.321) -- (-1.395,0.307) -- (-1.350,0.292) -- (-1.308,0.274) -- (-1.269,0.255) -- (-1.233,0.234) -- (-1.200,0.212) -- (-1.170,0.188) -- (-1.144,0.164) -- (-1.122,0.138) -- (-1.103,0.111) -- (-1.089,0.084) -- (-1.078,0.056) -- (-1.072,0.028) -- (-1.070,0.000) -- (-1.072,-0.028);
\draw[mesh] (-2.347,-0.533) -- (-2.345,-0.505) -- (-2.339,-0.475) -- (-2.329,-0.445) -- (-2.315,-0.415) -- (-2.297,-0.384) -- (-2.276,-0.354) -- (-2.250,-0.323) -- (-2.222,-0.293) -- (-2.190,-0.264) -- (-2.155,-0.235) -- (-2.117,-0.207) -- (-2.076,-0.180) -- (-2.034,-0.155) -- (-1.989,-0.131) -- (-1.942,-0.108) -- (-1.893,-0.088) -- (-1.844,-0.069) -- (-1.793,-0.052) -- (-1.742,-0.037) -- (-1.690,-0.024) -- (-1.639,-0.013) -- (-1.588,-0.005) -- (-1.537,0.001) -- (-1.487,0.005) -- (-1.439,0.006) -- (-1.392,0.005) -- (-1.347,0.001) -- (-1.304,-0.005) -- (-1.264,-0.013) -- (-1.226,-0.024) -- (-1.191,-0.037) -- (-1.159,-0.052) -- (-1.130,-0.069) -- (-1.105,-0.088) -- (-1.084,-0.108) -- (-1.066,-0.131) -- (-1.052,-0.155) -- (-1.042,-0.180);
\draw[mesh] (-2.339,-0.588) -- (-2.345,-0.561) -- (-2.347,-0.533);
\draw[mesh] (-2.104,-1.030) -- (-2.103,-1.001) -- (-2.097,-0.970) -- (-2.088,-0.938) -- (-2.076,-0.905) -- (-2.060,-0.871) -- (-2.040,-0.835) -- (-2.018,-0.800) -- (-1.992,-0.764) -- (-1.963,-0.727) -- (-1.932,-0.691) -- (-1.898,-0.655) -- (-1.862,-0.620) -- (-1.823,-0.585) -- (-1.783,-0.552) -- (-1.741,-0.519) -- (-1.698,-0.488) -- (-1.653,-0.459) -- (-1.608,-0.431) -- (-1.562,-0.405) -- (-1.516,-0.382) -- (-1.469,-0.360) -- (-1.423,-0.341) -- (-1.378,-0.324) -- (-1.334,-0.310) -- (-1.290,-0.299) -- (-1.248,-0.290) -- (-1.208,-0.284) -- (-1.169,-0.281) -- (-1.133,-0.281) -- (-1.099,-0.283) -- (-1.068,-0.289) -- (-1.039,-0.297) -- (-1.013,-0.308) -- (-0.991,-0.322) -- (-0.971,-0.338) -- (-0.955,-0.356) -- (-0.943,-0.378);
\draw[mesh] (-2.088,-1.107) -- (-2.097,-1.083) -- (-2.103,-1.058) -- (-2.104,-1.030);
\draw[mesh] (-1.718,-1.457) -- (-1.717,-1.428) -- (-1.712,-1.396) -- (-1.705,-1.362) -- (-1.695,-1.326) -- (-1.682,-1.288) -- (-1.666,-1.249) -- (-1.647,-1.209) -- (-1.626,-1.167) -- (-1.603,-1.125) -- (-1.577,-1.083) -- (-1.550,-1.040) -- (-1.520,-0.998) -- (-1.489,-0.955) -- (-1.456,-0.913) -- (-1.421,-0.873) -- (-1.386,-0.833) -- (-1.350,-0.794) -- (-1.313,-0.757) -- (-1.275,-0.722) -- (-1.237,-0.689) -- (-1.200,-0.658) -- (-1.162,-0.630) -- (-1.125,-0.604) -- (-1.089,-0.581) -- (-1.053,-0.560) -- (-1.019,-0.543) -- (-0.986,-0.529) -- (-0.955,-0.518) -- (-0.925,-0.511) -- (-0.897,-0.506) -- (-0.872,-0.505) -- (-0.848,-0.508) -- (-0.827,-0.513) -- (-0.809,-0.522) -- (-0.793,-0.535);
\draw[mesh] (-1.682,-1.564) -- (-1.695,-1.549) -- (-1.705,-1.530) -- (-1.712,-1.509) -- (-1.717,-1.484) -- (-1.718,-1.457);
\draw[mesh] (-1.215,-1.785) -- (-1.214,-1.755) -- (-1.211,-1.722) -- (-1.206,-1.687) -- (-1.198,-1.649) -- (-1.189,-1.609) -- (-1.178,-1.567) -- (-1.165,-1.523) -- (-1.150,-1.477) -- (-1.134,-1.431) -- (-1.115,-1.384) -- (-1.096,-1.336) -- (-1.075,-1.287) -- (-1.053,-1.239) -- (-1.029,-1.191) -- (-1.005,-1.143) -- (-0.980,-1.097) -- (-0.954,-1.051) -- (-0.928,-1.007) -- (-0.902,-0.965) -- (-0.875,-0.925) -- (-0.848,-0.887) -- (-0.822,-0.851) -- (-0.796,-0.818) -- (-0.770,-0.788) -- (-0.745,-0.761) -- (-0.721,-0.737) -- (-0.697,-0.717) -- (-0.675,-0.700) -- (-0.654,-0.687) -- (-0.635,-0.677) -- (-0.616,-0.671) -- (-0.600,-0.669) -- (-0.585,-0.671) -- (-0.572,-0.677);
\draw[mesh] (-1.178,-1.894) -- (-1.189,-1.885) -- (-1.198,-1.872) -- (-1.206,-1.855) -- (-1.211,-1.835) -- (-1.214,-1.811) -- (-1.215,-1.785);
\draw[mesh] (-0.629,-1.991) -- (-0.628,-1.961) -- (-0.627,-1.927) -- (-0.624,-1.891) -- (-0.620,-1.852) -- (-0.616,-1.810) -- (-0.610,-1.766) -- (-0.603,-1.720) -- (-0.595,-1.672) -- (-0.587,-1.623) -- (-0.577,-1.573) -- (-0.567,-1.521) -- (-0.556,-1.469) -- (-0.545,-1.417) -- (-0.533,-1.365) -- (-0.520,-1.314) -- (-0.507,-1.263) -- (-0.494,-1.213) -- (-0.480,-1.165) -- (-0.467,-1.118) -- (-0.453,-1.073) -- (-0.439,-1.031) -- (-0.425,-0.990) -- (-0.412,-0.953) -- (-0.399,-0.919) -- (-0.386,-0.887) -- (-0.373,-0.860) -- (-0.361,-0.835) -- (-0.349,-0.815) -- (-0.339,-0.798) -- (-0.328,-0.785) -- (-0.319,-0.776) -- (-0.311,-0.771) -- (-0.303,-0.770) -- (-0.296,-0.774);
\draw[mesh] (-0.610,-2.093) -- (-0.616,-2.086) -- (-0.620,-2.075) -- (-0.624,-2.059) -- (-0.627,-2.040) -- (-0.628,-2.017) -- (-0.629,-1.991);
\draw[mesh] (-0.000,-2.061) -- (-0.000,-2.031) -- (-0.000,-1.997) -- (-0.000,-1.961) -- (-0.000,-1.921) -- (-0.000,-1.879) -- (-0.000,-1.834) -- (-0.000,-1.787) -- (-0.000,-1.739) -- (-0.000,-1.689) -- (-0.000,-1.637) -- (-0.000,-1.585) -- (-0.000,-1.532) -- (-0.000,-1.478) -- (-0.000,-1.425) -- (-0.000,-1.372) -- (-0.000,-1.320) -- (-0.000,-1.268) -- (-0.000,-1.218) -- (-0.000,-1.170) -- (-0.000,-1.124) -- (-0.000,-1.080) -- (-0.000,-1.038) -- (-0.000,-0.999) -- (-0.000,-0.963) -- (-0.000,-0.930) -- (-0.000,-0.901) -- (-0.000,-0.876) -- (-0.000,-0.854) -- (-0.000,-0.836) -- (-0.000,-0.822) -- (-0.000,-0.812) -- (-0.000,-0.806) -- (-0.000,-0.804);
\draw[mesh] (-0.000,-2.164) -- (-0.000,-2.161) -- (-0.000,-2.155) -- (-0.000,-2.144) -- (-0.000,-2.129) -- (-0.000,-2.110) -- (-0.000,-2.087) -- (-0.000,-2.061);
\draw[mesh] (0.629,-1.991) -- (0.628,-1.961) -- (0.627,-1.927) -- (0.624,-1.891) -- (0.620,-1.852) -- (0.616,-1.810) -- (0.610,-1.766) -- (0.603,-1.720) -- (0.595,-1.672) -- (0.587,-1.623) -- (0.577,-1.573) -- (0.567,-1.521) -- (0.556,-1.469) -- (0.545,-1.417) -- (0.533,-1.365) -- (0.520,-1.314) -- (0.507,-1.263) -- (0.494,-1.213) -- (0.480,-1.165) -- (0.467,-1.118) -- (0.453,-1.073) -- (0.439,-1.031) -- (0.425,-0.990) -- (0.412,-0.953) -- (0.399,-0.919) -- (0.386,-0.887) -- (0.373,-0.860) -- (0.361,-0.835) -- (0.349,-0.815) -- (0.339,-0.798) -- (0.328,-0.785) -- (0.319,-0.776) -- (0.311,-0.771) -- (0.303,-0.770) -- (0.296,-0.774);
\draw[mesh] (0.610,-2.093) -- (0.616,-2.086) -- (0.620,-2.075) -- (0.624,-2.059) -- (0.627,-2.040) -- (0.628,-2.017) -- (0.629,-1.991);
\draw[mesh] (1.215,-1.785) -- (1.214,-1.755) -- (1.211,-1.722) -- (1.206,-1.687) -- (1.198,-1.649) -- (1.189,-1.609) -- (1.178,-1.567) -- (1.165,-1.523) -- (1.150,-1.477) -- (1.134,-1.431) -- (1.115,-1.384) -- (1.096,-1.336) -- (1.075,-1.287) -- (1.053,-1.239) -- (1.029,-1.191) -- (1.005,-1.143) -- (0.980,-1.097) -- (0.954,-1.051) -- (0.928,-1.007) -- (0.902,-0.965) -- (0.875,-0.925) -- (0.848,-0.887) -- (0.822,-0.851) -- (0.796,-0.818) -- (0.770,-0.788) -- (0.745,-0.761) -- (0.721,-0.737) -- (0.697,-0.717) -- (0.675,-0.700) -- (0.654,-0.687) -- (0.635,-0.677) -- (0.616,-0.671) -- (0.600,-0.669) -- (0.585,-0.671) -- (0.572,-0.677);
\draw[mesh] (1.178,-1.894) -- (1.189,-1.885) -- (1.198,-1.872) -- (1.206,-1.855) -- (1.211,-1.835) -- (1.214,-1.811) -- (1.215,-1.785);
\draw[mesh] (1.718,-1.457) -- (1.717,-1.428) -- (1.712,-1.396) -- (1.705,-1.362) -- (1.695,-1.326) -- (1.682,-1.288) -- (1.666,-1.249) -- (1.647,-1.209) -- (1.626,-1.167) -- (1.603,-1.125) -- (1.577,-1.083) -- (1.550,-1.040) -- (1.520,-0.998) -- (1.489,-0.955) -- (1.456,-0.913) -- (1.421,-0.873) -- (1.386,-0.833) -- (1.350,-0.794) -- (1.313,-0.757) -- (1.275,-0.722) -- (1.237,-0.689) -- (1.200,-0.658) -- (1.162,-0.630) -- (1.125,-0.604) -- (1.089,-0.581) -- (1.053,-0.560) -- (1.019,-0.543) -- (0.986,-0.529) -- (0.955,-0.518) -- (0.925,-0.511) -- (0.897,-0.506) -- (0.872,-0.505) -- (0.848,-0.508) -- (0.827,-0.513) -- (0.809,-0.522) -- (0.793,-0.535);
\draw[mesh] (1.682,-1.564) -- (1.695,-1.549) -- (1.705,-1.530) -- (1.712,-1.509) -- (1.717,-1.484) -- (1.718,-1.457);
\draw[mesh] (2.104,-1.030) -- (2.103,-1.001) -- (2.097,-0.970) -- (2.088,-0.938) -- (2.076,-0.905) -- (2.060,-0.871) -- (2.040,-0.835) -- (2.018,-0.800) -- (1.992,-0.764) -- (1.963,-0.727) -- (1.932,-0.691) -- (1.898,-0.655) -- (1.862,-0.620) -- (1.823,-0.585) -- (1.783,-0.552) -- (1.741,-0.519) -- (1.698,-0.488) -- (1.653,-0.459) -- (1.608,-0.431) -- (1.562,-0.405) -- (1.516,-0.382) -- (1.469,-0.360) -- (1.423,-0.341) -- (1.378,-0.324) -- (1.334,-0.310) -- (1.290,-0.299) -- (1.248,-0.290) -- (1.208,-0.284) -- (1.169,-0.281) -- (1.133,-0.281) -- (1.099,-0.283) -- (1.068,-0.289) -- (1.039,-0.297) -- (1.013,-0.308) -- (0.991,-0.322) -- (0.971,-0.338) -- (0.955,-0.356) -- (0.943,-0.378);
\draw[mesh] (2.088,-1.107) -- (2.097,-1.083) -- (2.103,-1.058) -- (2.104,-1.030);
\draw[mesh] (2.347,-0.533) -- (2.345,-0.505) -- (2.339,-0.475) -- (2.329,-0.445) -- (2.315,-0.415) -- (2.297,-0.384) -- (2.276,-0.354) -- (2.250,-0.323) -- (2.222,-0.293) -- (2.190,-0.264) -- (2.155,-0.235) -- (2.117,-0.207) -- (2.076,-0.180) -- (2.034,-0.155) -- (1.989,-0.131) -- (1.942,-0.108) -- (1.893,-0.088) -- (1.844,-0.069) -- (1.793,-0.052) -- (1.742,-0.037) -- (1.690,-0.024) -- (1.639,-0.013) -- (1.588,-0.005) -- (1.537,0.001) -- (1.487,0.005) -- (1.439,0.006) -- (1.392,0.005) -- (1.347,0.001) -- (1.304,-0.005) -- (1.264,-0.013) -- (1.226,-0.024) -- (1.191,-0.037) -- (1.159,-0.052) -- (1.130,-0.069) -- (1.105,-0.088) -- (1.084,-0.108) -- (1.066,-0.131) -- (1.052,-0.155) -- (1.042,-0.180);
\draw[mesh] (2.339,-0.588) -- (2.345,-0.561) -- (2.347,-0.533);
\draw[mesh] (-2.428,-0.081) -- (-2.423,-0.162) -- (-2.413,-0.242) -- (-2.400,-0.322) -- (-2.383,-0.402) -- (-2.363,-0.481) -- (-2.339,-0.559) -- (-2.311,-0.637) -- (-2.280,-0.713) -- (-2.245,-0.789) -- (-2.207,-0.863) -- (-2.165,-0.936) -- (-2.120,-1.007) -- (-2.072,-1.077) -- (-2.020,-1.145) -- (-1.966,-1.211) -- (-1.908,-1.276) -- (-1.848,-1.338) -- (-1.784,-1.399) -- (-1.718,-1.457) -- (-1.649,-1.513) -- (-1.578,-1.567) -- (-1.504,-1.618) -- (-1.428,-1.667) -- (-1.350,-1.713) -- (-1.270,-1.757) -- (-1.187,-1.798) -- (-1.103,-1.836) -- (-1.017,-1.871) -- (-0.930,-1.904) -- (-0.841,-1.933) -- (-0.751,-1.960) -- (-0.660,-1.983) -- (-0.567,-2.004) -- (-0.474,-2.021) -- (-0.380,-2.035) -- (-0.286,-2.046) -- (-0.191,-2.054) -- (-0.095,-2.059) -- (-0.000,-2.061) -- (0.095,-2.059) -- (0.191,-2.054) -- (0.286,-2.046) -- (0.380,-2.035) -- (0.474,-2.021) -- (0.567,-2.004) -- (0.660,-1.983) -- (0.751,-1.960) -- (0.841,-1.933) -- (0.930,-1.904) -- (1.017,-1.871) -- (1.103,-1.836) -- (1.187,-1.798) -- (1.270,-1.757) -- (1.350,-1.713) -- (1.428,-1.667) -- (1.504,-1.618) -- (1.578,-1.567) -- (1.649,-1.513) -- (1.718,-1.457) -- (1.784,-1.399) -- (1.848,-1.338) -- (1.908,-1.276) -- (1.966,-1.211) -- (2.020,-1.145) -- (2.072,-1.077) -- (2.120,-1.007) -- (2.165,-0.936) -- (2.207,-0.863) -- (2.245,-0.789) -- (2.280,-0.713) -- (2.311,-0.637) -- (2.339,-0.559) -- (2.363,-0.481) -- (2.383,-0.402) -- (2.400,-0.322) -- (2.413,-0.242) -- (2.423,-0.162) -- (2.428,-0.081) -- (2.430,-0.000);
\draw[mesh] (2.231,0.255) -- (2.229,0.329) -- (2.224,0.403) -- (2.215,0.477) -- (2.203,0.551) -- (2.188,0.624) -- (2.169,0.696) -- (2.147,0.768) -- (2.122,0.839) -- (2.093,0.910) -- (2.061,0.979) -- (2.026,1.047) -- (1.988,1.114) -- (1.946,1.179) -- (1.902,1.243) -- (1.855,1.306) -- (1.805,1.367) -- (1.752,1.426) -- (1.696,1.483) -- (1.638,1.539) -- (1.577,1.593) -- (1.514,1.644) -- (1.449,1.693) -- (1.381,1.741) -- (1.311,1.785) -- (1.239,1.828) -- (1.166,1.868) -- (1.090,1.905) -- (1.013,1.940) -- (0.934,1.973) -- (0.854,2.003) -- (0.772,2.030) -- (0.689,2.054) -- (0.606,2.076) -- (0.521,2.094) -- (0.435,2.110) -- (0.349,2.123) -- (0.262,2.134) -- (0.175,2.141) -- (0.088,2.145) -- (0.000,2.147) -- (-0.088,2.145) -- (-0.175,2.141) -- (-0.262,2.134) -- (-0.349,2.123) -- (-0.435,2.110) -- (-0.521,2.094) -- (-0.606,2.076) -- (-0.689,2.054) -- (-0.772,2.030) -- (-0.854,2.003) -- (-0.934,1.973) -- (-1.013,1.940) -- (-1.090,1.905) -- (-1.166,1.868) -- (-1.239,1.828) -- (-1.311,1.785) -- (-1.381,1.741) -- (-1.449,1.693) -- (-1.514,1.644) -- (-1.577,1.593) -- (-1.638,1.539) -- (-1.696,1.483) -- (-1.752,1.426) -- (-1.805,1.367) -- (-1.855,1.306) -- (-1.902,1.243) -- (-1.946,1.179) -- (-1.988,1.114) -- (-2.026,1.047) -- (-2.061,0.979) -- (-2.093,0.910) -- (-2.122,0.839) -- (-2.147,0.768) -- (-2.169,0.696) -- (-2.188,0.624) -- (-2.203,0.551) -- (-2.215,0.477) -- (-2.224,0.403) -- (-2.229,0.329) -- (-2.231,0.255) -- (-2.229,0.181) -- (-2.224,0.106) -- (-2.215,0.032) -- (-2.203,-0.041) -- (-2.188,-0.114) -- (-2.169,-0.187) -- (-2.147,-0.259) -- (-2.122,-0.330) -- (-2.093,-0.400) -- (-2.061,-0.469) -- (-2.026,-0.537) -- (-1.988,-0.604) -- (-1.946,-0.670) -- (-1.902,-0.734) -- (-1.855,-0.796) -- (-1.805,-0.857) -- (-1.752,-0.916) -- (-1.696,-0.974) -- (-1.638,-1.029) -- (-1.577,-1.083) -- (-1.514,-1.134) -- (-1.449,-1.184) -- (-1.381,-1.231) -- (-1.311,-1.276) -- (-1.239,-1.318) -- (-1.166,-1.358) -- (-1.090,-1.396) -- (-1.013,-1.431) -- (-0.934,-1.463) -- (-0.854,-1.493) -- (-0.772,-1.520) -- (-0.689,-1.544) -- (-0.606,-1.566) -- (-0.521,-1.585) -- (-0.435,-1.601) -- (-0.349,-1.614) -- (-0.262,-1.624) -- (-0.175,-1.631) -- (-0.088,-1.636) -- (-0.000,-1.637) -- (0.088,-1.636) -- (0.175,-1.631) -- (0.262,-1.624) -- (0.349,-1.614) -- (0.435,-1.601) -- (0.521,-1.585) -- (0.606,-1.566) -- (0.689,-1.544) -- (0.772,-1.520) -- (0.854,-1.493) -- (0.934,-1.463) -- (1.013,-1.431) -- (1.090,-1.396) -- (1.166,-1.358) -- (1.239,-1.318) -- (1.311,-1.276) -- (1.381,-1.231) -- (1.449,-1.184) -- (1.514,-1.134) -- (1.577,-1.083) -- (1.638,-1.029) -- (1.696,-0.974) -- (1.752,-0.916) -- (1.805,-0.857) -- (1.855,-0.796) -- (1.902,-0.734) -- (1.946,-0.670) -- (1.988,-0.604) -- (2.026,-0.537) -- (2.061,-0.469) -- (2.093,-0.400) -- (2.122,-0.330) -- (2.147,-0.259) -- (2.169,-0.187) -- (2.188,-0.114) -- (2.203,-0.041) -- (2.215,0.032) -- (2.224,0.106) -- (2.229,0.181) -- (2.231,0.255);
\draw[mesh] (1.750,0.360) -- (1.749,0.419) -- (1.745,0.477) -- (1.738,0.535) -- (1.728,0.593) -- (1.716,0.650) -- (1.702,0.707) -- (1.684,0.763) -- (1.664,0.819) -- (1.642,0.874) -- (1.617,0.928) -- (1.589,0.982) -- (1.559,1.034) -- (1.527,1.086) -- (1.492,1.136) -- (1.455,1.185) -- (1.416,1.233) -- (1.374,1.279) -- (1.331,1.324) -- (1.285,1.368) -- (1.237,1.410) -- (1.188,1.450) -- (1.137,1.489) -- (1.083,1.526) -- (1.029,1.561) -- (0.972,1.594) -- (0.914,1.626) -- (0.855,1.655) -- (0.794,1.683) -- (0.733,1.708) -- (0.670,1.731) -- (0.606,1.753) -- (0.541,1.772) -- (0.475,1.789) -- (0.409,1.803) -- (0.341,1.816) -- (0.274,1.826) -- (0.206,1.834) -- (0.137,1.840) -- (0.069,1.843) -- (0.000,1.844) -- (-0.069,1.843) -- (-0.137,1.840) -- (-0.206,1.834) -- (-0.274,1.826) -- (-0.341,1.816) -- (-0.409,1.803) -- (-0.475,1.789) -- (-0.541,1.772) -- (-0.606,1.753) -- (-0.670,1.731) -- (-0.733,1.708) -- (-0.794,1.683) -- (-0.855,1.655) -- (-0.914,1.626) -- (-0.972,1.594) -- (-1.029,1.561) -- (-1.083,1.526) -- (-1.137,1.489) -- (-1.188,1.450) -- (-1.237,1.410) -- (-1.285,1.368) -- (-1.331,1.324) -- (-1.374,1.279) -- (-1.416,1.233) -- (-1.455,1.185) -- (-1.492,1.136) -- (-1.527,1.086) -- (-1.559,1.034) -- (-1.589,0.982) -- (-1.617,0.928) -- (-1.642,0.874) -- (-1.664,0.819) -- (-1.684,0.763) -- (-1.702,0.707) -- (-1.716,0.650) -- (-1.728,0.593) -- (-1.738,0.535) -- (-1.745,0.477) -- (-1.749,0.419) -- (-1.750,0.360) -- (-1.749,0.302) -- (-1.745,0.244) -- (-1.738,0.186) -- (-1.728,0.128) -- (-1.716,0.071) -- (-1.702,0.014) -- (-1.684,-0.042) -- (-1.664,-0.098) -- (-1.642,-0.153) -- (-1.617,-0.208) -- (-1.589,-0.261) -- (-1.559,-0.313) -- (-1.527,-0.365) -- (-1.492,-0.415) -- (-1.455,-0.464) -- (-1.416,-0.512) -- (-1.374,-0.558) -- (-1.331,-0.603) -- (-1.285,-0.647) -- (-1.237,-0.689) -- (-1.188,-0.729) -- (-1.137,-0.768) -- (-1.083,-0.805) -- (-1.029,-0.840) -- (-0.972,-0.874) -- (-0.914,-0.905) -- (-0.855,-0.935) -- (-0.794,-0.962) -- (-0.733,-0.987) -- (-0.670,-1.011) -- (-0.606,-1.032) -- (-0.541,-1.051) -- (-0.475,-1.068) -- (-0.409,-1.083) -- (-0.341,-1.095) -- (-0.274,-1.105) -- (-0.206,-1.113) -- (-0.137,-1.119) -- (-0.069,-1.123) -- (-0.000,-1.124) -- (0.069,-1.123) -- (0.137,-1.119) -- (0.206,-1.113) -- (0.274,-1.105) -- (0.341,-1.095) -- (0.409,-1.083) -- (0.475,-1.068) -- (0.541,-1.051) -- (0.606,-1.032) -- (0.670,-1.011) -- (0.733,-0.987) -- (0.794,-0.962) -- (0.855,-0.935) -- (0.914,-0.905) -- (0.972,-0.874) -- (1.029,-0.840) -- (1.083,-0.805) -- (1.137,-0.768) -- (1.188,-0.729) -- (1.237,-0.689) -- (1.285,-0.647) -- (1.331,-0.603) -- (1.374,-0.558) -- (1.416,-0.512) -- (1.455,-0.464) -- (1.492,-0.415) -- (1.527,-0.365) -- (1.559,-0.313) -- (1.589,-0.261) -- (1.617,-0.208) -- (1.642,-0.153) -- (1.664,-0.098) -- (1.684,-0.042) -- (1.702,0.014) -- (1.716,0.071) -- (1.728,0.128) -- (1.738,0.186) -- (1.745,0.244) -- (1.749,0.302) -- (1.750,0.360);
\draw[mesh] (1.269,0.255) -- (1.268,0.297) -- (1.265,0.339) -- (1.260,0.381) -- (1.254,0.423) -- (1.245,0.465) -- (1.234,0.506) -- (1.222,0.547) -- (1.207,0.587) -- (1.191,0.627) -- (1.173,0.667) -- (1.153,0.705) -- (1.131,0.743) -- (1.107,0.781) -- (1.082,0.817) -- (1.055,0.853) -- (1.027,0.887) -- (0.997,0.921) -- (0.965,0.954) -- (0.932,0.985) -- (0.897,1.016) -- (0.862,1.045) -- (0.824,1.073) -- (0.786,1.100) -- (0.746,1.126) -- (0.705,1.150) -- (0.663,1.173) -- (0.620,1.194) -- (0.576,1.214) -- (0.531,1.232) -- (0.486,1.249) -- (0.439,1.265) -- (0.392,1.278) -- (0.345,1.291) -- (0.296,1.301) -- (0.248,1.310) -- (0.199,1.318) -- (0.149,1.324) -- (0.100,1.328) -- (0.050,1.330) -- (0.000,1.331) -- (-0.050,1.330) -- (-0.100,1.328) -- (-0.149,1.324) -- (-0.199,1.318) -- (-0.248,1.310) -- (-0.296,1.301) -- (-0.345,1.291) -- (-0.392,1.278) -- (-0.439,1.265) -- (-0.486,1.249) -- (-0.531,1.232) -- (-0.576,1.214) -- (-0.620,1.194) -- (-0.663,1.173) -- (-0.705,1.150) -- (-0.746,1.126) -- (-0.786,1.100) -- (-0.824,1.073) -- (-0.862,1.045) -- (-0.897,1.016) -- (-0.932,0.985) -- (-0.965,0.954) -- (-0.997,0.921) -- (-1.027,0.887) -- (-1.055,0.853) -- (-1.082,0.817) -- (-1.107,0.781) -- (-1.131,0.743) -- (-1.153,0.705) -- (-1.173,0.667) -- (-1.191,0.627) -- (-1.207,0.587) -- (-1.222,0.547) -- (-1.234,0.506) -- (-1.245,0.465) -- (-1.254,0.423) -- (-1.260,0.381) -- (-1.265,0.339) -- (-1.268,0.297) -- (-1.269,0.255) -- (-1.268,0.213) -- (-1.265,0.170) -- (-1.260,0.128) -- (-1.254,0.086) -- (-1.245,0.045) -- (-1.234,0.004) -- (-1.222,-0.037) -- (-1.207,-0.078) -- (-1.191,-0.118) -- (-1.173,-0.157) -- (-1.153,-0.196) -- (-1.131,-0.234) -- (-1.107,-0.271) -- (-1.082,-0.308) -- (-1.055,-0.343) -- (-1.027,-0.378) -- (-0.997,-0.412) -- (-0.965,-0.444) -- (-0.932,-0.476) -- (-0.897,-0.506) -- (-0.862,-0.536) -- (-0.824,-0.564) -- (-0.786,-0.590) -- (-0.746,-0.616) -- (-0.705,-0.640) -- (-0.663,-0.663) -- (-0.620,-0.684) -- (-0.576,-0.704) -- (-0.531,-0.723) -- (-0.486,-0.740) -- (-0.439,-0.755) -- (-0.392,-0.769) -- (-0.345,-0.781) -- (-0.296,-0.792) -- (-0.248,-0.801) -- (-0.199,-0.808) -- (-0.149,-0.814) -- (-0.100,-0.818) -- (-0.050,-0.821) -- (-0.000,-0.822) -- (0.050,-0.821) -- (0.100,-0.818) -- (0.149,-0.814) -- (0.199,-0.808) -- (0.248,-0.801) -- (0.296,-0.792) -- (0.345,-0.781) -- (0.392,-0.769) -- (0.439,-0.755) -- (0.486,-0.740) -- (0.531,-0.723) -- (0.576,-0.704) -- (0.620,-0.684) -- (0.663,-0.663) -- (0.705,-0.640) -- (0.746,-0.616) -- (0.786,-0.590) -- (0.824,-0.564) -- (0.862,-0.536) -- (0.897,-0.506) -- (0.932,-0.476) -- (0.965,-0.444) -- (0.997,-0.412) -- (1.027,-0.378) -- (1.055,-0.343) -- (1.082,-0.308) -- (1.107,-0.271) -- (1.131,-0.234) -- (1.153,-0.196) -- (1.173,-0.157) -- (1.191,-0.118) -- (1.207,-0.078) -- (1.222,-0.037) -- (1.234,0.004) -- (1.245,0.045) -- (1.254,0.086) -- (1.260,0.128) -- (1.265,0.170) -- (1.268,0.213) -- (1.269,0.255);
\draw[mesh] (1.070,0.000) -- (1.069,0.036) -- (1.067,0.071) -- (1.063,0.107) -- (1.057,0.142) -- (1.049,0.177) -- (1.040,0.212) -- (1.030,0.246) -- (1.018,0.280) -- (1.004,0.314) -- (0.989,0.347) -- (0.972,0.380) -- (0.953,0.412) -- (0.934,0.443) -- (0.912,0.474) -- (0.890,0.504) -- (0.866,0.533) -- (0.840,0.562) -- (0.814,0.589) -- (0.786,0.616) -- (0.757,0.642) -- (0.726,0.666) -- (0.695,0.690) -- (0.662,0.713) -- (0.629,0.734) -- (0.594,0.754) -- (0.559,0.774) -- (0.523,0.792) -- (0.486,0.809) -- (0.448,0.824) -- (0.409,0.838) -- (0.370,0.851) -- (0.331,0.863) -- (0.290,0.873) -- (0.250,0.882) -- (0.209,0.890) -- (0.167,0.896) -- (0.126,0.901) -- (0.084,0.905) -- (0.042,0.907) -- (0.000,0.907) -- (-0.042,0.907) -- (-0.084,0.905) -- (-0.126,0.901) -- (-0.167,0.896) -- (-0.209,0.890) -- (-0.250,0.882) -- (-0.290,0.873) -- (-0.331,0.863) -- (-0.370,0.851) -- (-0.409,0.838) -- (-0.448,0.824) -- (-0.486,0.809) -- (-0.523,0.792) -- (-0.559,0.774) -- (-0.594,0.754) -- (-0.629,0.734) -- (-0.662,0.713) -- (-0.695,0.690) -- (-0.726,0.666) -- (-0.757,0.642) -- (-0.786,0.616) -- (-0.814,0.589) -- (-0.840,0.562) -- (-0.866,0.533) -- (-0.890,0.504) -- (-0.912,0.474) -- (-0.934,0.443) -- (-0.953,0.412) -- (-0.972,0.380) -- (-0.989,0.347) -- (-1.004,0.314) -- (-1.018,0.280) -- (-1.030,0.246) -- (-1.040,0.212) -- (-1.049,0.177) -- (-1.057,0.142) -- (-1.063,0.107) -- (-1.067,0.071) -- (-1.069,0.036) -- (-1.070,0.000) -- (-1.069,-0.036);

    % u = t, v = phi t for 12 turns of u, generated by torus-knot-closed-and-open.py
\draw[fillback] (2.430,0.000) -- (2.430,0.033);
\draw[fill] (2.430,0.033) -- (2.429,0.066) -- (2.427,0.100) -- (2.425,0.133) -- (2.422,0.166) -- (2.418,0.199) -- (2.414,0.232) -- (2.409,0.264) -- (2.403,0.297) -- (2.397,0.330) -- (2.390,0.362) -- (2.383,0.394) -- (2.374,0.426) -- (2.366,0.458) -- (2.356,0.489) -- (2.346,0.521) -- (2.335,0.552) -- (2.324,0.583) -- (2.312,0.613) -- (2.300,0.643) -- (2.287,0.673) -- (2.273,0.703) -- (2.259,0.732) -- (2.245,0.761) -- (2.229,0.790) -- (2.214,0.818) -- (2.197,0.846) -- (2.180,0.873) -- (2.163,0.900) -- (2.145,0.927) -- (2.127,0.953) -- (2.108,0.979) -- (2.089,1.004) -- (2.069,1.029) -- (2.049,1.053) -- (2.029,1.077) -- (2.008,1.100) -- (1.987,1.123) -- (1.965,1.146) -- (1.943,1.168) -- (1.921,1.189) -- (1.898,1.210) -- (1.875,1.230) -- (1.851,1.250) -- (1.828,1.269) -- (1.804,1.288) -- (1.779,1.306) -- (1.755,1.324) -- (1.730,1.341) -- (1.705,1.357) -- (1.680,1.373) -- (1.655,1.388) -- (1.629,1.403) -- (1.603,1.417) -- (1.577,1.431) -- (1.551,1.444) -- (1.525,1.456) -- (1.499,1.468) -- (1.473,1.479) -- (1.446,1.490) -- (1.419,1.500) -- (1.393,1.510) -- (1.366,1.519) -- (1.339,1.527) -- (1.313,1.535) -- (1.286,1.542) -- (1.259,1.549) -- (1.233,1.555) -- (1.206,1.560) -- (1.179,1.565) -- (1.153,1.570) -- (1.126,1.574) -- (1.100,1.577) -- (1.073,1.580) -- (1.047,1.582) -- (1.021,1.584) -- (0.995,1.585) -- (0.969,1.586) -- (0.943,1.586) -- (0.917,1.586) -- (0.892,1.585) -- (0.867,1.584) -- (0.841,1.582) -- (0.816,1.580) -- (0.792,1.577) -- (0.767,1.574) -- (0.743,1.571) -- (0.719,1.567) -- (0.695,1.562) -- (0.671,1.557) -- (0.647,1.552) -- (0.624,1.547) -- (0.601,1.541) -- (0.578,1.534) -- (0.556,1.528) -- (0.534,1.521) -- (0.512,1.513) -- (0.490,1.506) -- (0.468,1.498) -- (0.447,1.489) -- (0.426,1.481) -- (0.405,1.472) -- (0.385,1.462) -- (0.365,1.453) -- (0.345,1.443) -- (0.325,1.433) -- (0.306,1.423) -- (0.287,1.413) -- (0.268,1.402) -- (0.249,1.391) -- (0.231,1.380) -- (0.213,1.369) -- (0.195,1.358) -- (0.177,1.346) -- (0.160,1.334) -- (0.143,1.323) -- (0.126,1.311) -- (0.110,1.299) -- (0.093,1.286) -- (0.077,1.274) -- (0.061,1.262) -- (0.046,1.249) -- (0.030,1.237) -- (0.015,1.224) -- (0.000,1.212) -- (-0.015,1.199) -- (-0.029,1.186) -- (-0.044,1.173) -- (-0.058,1.161) -- (-0.072,1.148) -- (-0.086,1.135) -- (-0.100,1.122) -- (-0.114,1.109) -- (-0.127,1.097) -- (-0.141,1.084) -- (-0.154,1.071) -- (-0.167,1.058) -- (-0.180,1.046) -- (-0.193,1.033) -- (-0.206,1.021) -- (-0.219,1.008) -- (-0.231,0.996) -- (-0.244,0.983) -- (-0.257,0.971) -- (-0.269,0.959) -- (-0.282,0.946) -- (-0.294,0.934) -- (-0.307,0.922) -- (-0.319,0.910) -- (-0.332,0.898) -- (-0.344,0.886) -- (-0.356,0.875) -- (-0.369,0.863) -- (-0.381,0.852) -- (-0.394,0.840) -- (-0.406,0.829) -- (-0.419,0.817) -- (-0.432,0.806) -- (-0.445,0.795) -- (-0.457,0.784) -- (-0.470,0.773) -- (-0.483,0.762) -- (-0.496,0.752) -- (-0.510,0.741) -- (-0.523,0.730) -- (-0.536,0.720) -- (-0.550,0.709) -- (-0.564,0.699) -- (-0.577,0.689) -- (-0.591,0.678) -- (-0.605,0.668) -- (-0.620,0.658) -- (-0.634,0.648) -- (-0.648,0.638) -- (-0.663,0.628) -- (-0.678,0.618) -- (-0.693,0.608) -- (-0.708,0.599);
\draw[fillback] (-0.708,0.599) -- (-0.723,0.589) -- (-0.739,0.579) -- (-0.755,0.569) -- (-0.771,0.560) -- (-0.787,0.550) -- (-0.803,0.540) -- (-0.819,0.530) -- (-0.836,0.521) -- (-0.852,0.511) -- (-0.869,0.501) -- (-0.886,0.491) -- (-0.904,0.482) -- (-0.921,0.472) -- (-0.939,0.462) -- (-0.956,0.452) -- (-0.974,0.442) -- (-0.992,0.432) -- (-1.010,0.422) -- (-1.029,0.412) -- (-1.047,0.401) -- (-1.066,0.391) -- (-1.085,0.381) -- (-1.104,0.370) -- (-1.123,0.359) -- (-1.142,0.349) -- (-1.161,0.338) -- (-1.180,0.327) -- (-1.200,0.316) -- (-1.219,0.305) -- (-1.239,0.293) -- (-1.259,0.282) -- (-1.278,0.270) -- (-1.298,0.259) -- (-1.318,0.247) -- (-1.338,0.235) -- (-1.358,0.223) -- (-1.378,0.210) -- (-1.398,0.198) -- (-1.418,0.185) -- (-1.438,0.172) -- (-1.457,0.159) -- (-1.477,0.146) -- (-1.497,0.133) -- (-1.517,0.119) -- (-1.537,0.105) -- (-1.556,0.091) -- (-1.576,0.077) -- (-1.595,0.063) -- (-1.614,0.048) -- (-1.634,0.034) -- (-1.653,0.019) -- (-1.671,0.004) -- (-1.690,-0.011) -- (-1.709,-0.027) -- (-1.727,-0.043) -- (-1.745,-0.058) -- (-1.763,-0.074) -- (-1.781,-0.091) -- (-1.798,-0.107) -- (-1.815,-0.124) -- (-1.832,-0.140) -- (-1.849,-0.157) -- (-1.865,-0.174) -- (-1.881,-0.192) -- (-1.897,-0.209) -- (-1.912,-0.227) -- (-1.927,-0.245) -- (-1.942,-0.263) -- (-1.956,-0.281) -- (-1.970,-0.299) -- (-1.983,-0.317) -- (-1.996,-0.336) -- (-2.009,-0.355) -- (-2.021,-0.373) -- (-2.033,-0.392) -- (-2.045,-0.411) -- (-2.055,-0.431) -- (-2.066,-0.450) -- (-2.076,-0.469) -- (-2.085,-0.489) -- (-2.094,-0.508) -- (-2.103,-0.528) -- (-2.111,-0.548) -- (-2.118,-0.567) -- (-2.125,-0.587) -- (-2.131,-0.607) -- (-2.137,-0.627) -- (-2.142,-0.647) -- (-2.147,-0.667) -- (-2.151,-0.687) -- (-2.154,-0.707) -- (-2.157,-0.726) -- (-2.159,-0.746) -- (-2.161,-0.766) -- (-2.162,-0.786) -- (-2.163,-0.806) -- (-2.162,-0.826) -- (-2.162,-0.845) -- (-2.160,-0.865) -- (-2.158,-0.884) -- (-2.156,-0.904) -- (-2.152,-0.923) -- (-2.149,-0.942) -- (-2.144,-0.961) -- (-2.139,-0.980) -- (-2.133,-0.999) -- (-2.127,-1.018) -- (-2.120,-1.036) -- (-2.112,-1.055) -- (-2.104,-1.073) -- (-2.095,-1.090) -- (-2.085,-1.108) -- (-2.075,-1.126) -- (-2.065,-1.143) -- (-2.053,-1.160) -- (-2.041,-1.177);
\draw[fill] (-2.041,-1.177) -- (-2.029,-1.193) -- (-2.015,-1.209) -- (-2.002,-1.225) -- (-1.987,-1.241) -- (-1.972,-1.256) -- (-1.957,-1.271) -- (-1.941,-1.286) -- (-1.924,-1.300) -- (-1.907,-1.315) -- (-1.889,-1.328) -- (-1.871,-1.342) -- (-1.852,-1.355) -- (-1.832,-1.368) -- (-1.813,-1.380) -- (-1.792,-1.392) -- (-1.771,-1.403) -- (-1.750,-1.415) -- (-1.728,-1.425) -- (-1.706,-1.436) -- (-1.683,-1.446) -- (-1.660,-1.455) -- (-1.636,-1.465) -- (-1.612,-1.473) -- (-1.588,-1.482) -- (-1.563,-1.490) -- (-1.538,-1.497) -- (-1.513,-1.504) -- (-1.487,-1.511) -- (-1.461,-1.517) -- (-1.434,-1.522) -- (-1.407,-1.528) -- (-1.380,-1.533) -- (-1.353,-1.537) -- (-1.325,-1.541) -- (-1.297,-1.544) -- (-1.269,-1.547) -- (-1.241,-1.550) -- (-1.213,-1.552) -- (-1.184,-1.553) -- (-1.155,-1.555) -- (-1.126,-1.555) -- (-1.097,-1.556) -- (-1.068,-1.555) -- (-1.038,-1.555) -- (-1.009,-1.554) -- (-0.979,-1.552) -- (-0.949,-1.550) -- (-0.920,-1.548) -- (-0.890,-1.545) -- (-0.860,-1.542) -- (-0.830,-1.538) -- (-0.801,-1.534) -- (-0.771,-1.529) -- (-0.741,-1.525) -- (-0.712,-1.519) -- (-0.682,-1.514) -- (-0.653,-1.507) -- (-0.623,-1.501) -- (-0.594,-1.494) -- (-0.565,-1.487) -- (-0.536,-1.479) -- (-0.507,-1.471) -- (-0.478,-1.463) -- (-0.449,-1.454) -- (-0.421,-1.445) -- (-0.393,-1.436) -- (-0.365,-1.426) -- (-0.337,-1.417) -- (-0.309,-1.406) -- (-0.282,-1.396) -- (-0.255,-1.385) -- (-0.228,-1.374) -- (-0.202,-1.363) -- (-0.175,-1.351) -- (-0.149,-1.339) -- (-0.124,-1.327) -- (-0.098,-1.315) -- (-0.073,-1.303) -- (-0.049,-1.290) -- (-0.024,-1.277) -- (-0.000,-1.264) -- (0.024,-1.251) -- (0.047,-1.238) -- (0.070,-1.224) -- (0.093,-1.211) -- (0.115,-1.197) -- (0.137,-1.183) -- (0.159,-1.169) -- (0.180,-1.155) -- (0.201,-1.141) -- (0.222,-1.127) -- (0.242,-1.113) -- (0.261,-1.098) -- (0.281,-1.084) -- (0.300,-1.070) -- (0.318,-1.056) -- (0.336,-1.041) -- (0.354,-1.027) -- (0.372,-1.013) -- (0.389,-0.998) -- (0.405,-0.984) -- (0.422,-0.970) -- (0.437,-0.956) -- (0.453,-0.942) -- (0.468,-0.928) -- (0.483,-0.914) -- (0.497,-0.900) -- (0.511,-0.886) -- (0.525,-0.873) -- (0.538,-0.859) -- (0.551,-0.846) -- (0.564,-0.833) -- (0.576,-0.820) -- (0.588,-0.807) -- (0.600,-0.794) -- (0.611,-0.782) -- (0.622,-0.769) -- (0.632,-0.757) -- (0.643,-0.745) -- (0.653,-0.733) -- (0.662,-0.721) -- (0.672,-0.710) -- (0.681,-0.698) -- (0.690,-0.687) -- (0.699,-0.676) -- (0.707,-0.666) -- (0.715,-0.655) -- (0.723,-0.645) -- (0.731,-0.635) -- (0.738,-0.625) -- (0.745,-0.616) -- (0.752,-0.606) -- (0.759,-0.597) -- (0.766,-0.588) -- (0.773,-0.580) -- (0.779,-0.571) -- (0.785,-0.563) -- (0.791,-0.555) -- (0.797,-0.548) -- (0.803,-0.540) -- (0.809,-0.533) -- (0.814,-0.526) -- (0.820,-0.519) -- (0.825,-0.512) -- (0.830,-0.506) -- (0.836,-0.500) -- (0.841,-0.494) -- (0.846,-0.488) -- (0.851,-0.482) -- (0.856,-0.477) -- (0.861,-0.472);
\draw[fillback] (0.861,-0.472) -- (0.866,-0.467) -- (0.871,-0.462) -- (0.876,-0.457) -- (0.881,-0.453) -- (0.886,-0.448) -- (0.891,-0.444) -- (0.896,-0.440) -- (0.902,-0.436) -- (0.907,-0.432) -- (0.912,-0.429) -- (0.917,-0.425) -- (0.923,-0.422) -- (0.928,-0.418) -- (0.934,-0.415) -- (0.939,-0.412) -- (0.945,-0.409) -- (0.951,-0.406) -- (0.956,-0.403) -- (0.962,-0.400) -- (0.968,-0.397) -- (0.975,-0.394) -- (0.981,-0.392) -- (0.987,-0.389) -- (0.994,-0.386) -- (1.000,-0.383) -- (1.007,-0.380) -- (1.014,-0.377) -- (1.021,-0.375) -- (1.028,-0.372) -- (1.035,-0.369) -- (1.042,-0.365) -- (1.050,-0.362) -- (1.057,-0.359) -- (1.065,-0.356) -- (1.073,-0.352) -- (1.081,-0.348) -- (1.089,-0.345) -- (1.097,-0.341) -- (1.105,-0.337) -- (1.114,-0.333) -- (1.122,-0.328) -- (1.131,-0.324) -- (1.139,-0.319) -- (1.148,-0.314) -- (1.157,-0.309) -- (1.166,-0.303) -- (1.175,-0.298) -- (1.184,-0.292) -- (1.193,-0.286) -- (1.202,-0.279) -- (1.211,-0.273) -- (1.221,-0.266) -- (1.230,-0.259) -- (1.239,-0.251) -- (1.249,-0.243) -- (1.258,-0.235) -- (1.267,-0.227) -- (1.277,-0.218) -- (1.286,-0.209) -- (1.295,-0.200) -- (1.304,-0.190) -- (1.314,-0.180) -- (1.323,-0.170) -- (1.332,-0.159) -- (1.341,-0.148) -- (1.350,-0.137) -- (1.359,-0.125) -- (1.368,-0.113) -- (1.376,-0.101) -- (1.385,-0.088) -- (1.393,-0.075) -- (1.401,-0.061) -- (1.409,-0.047) -- (1.417,-0.033) -- (1.425,-0.018) -- (1.432,-0.003) -- (1.440,0.012) -- (1.447,0.028) -- (1.454,0.044) -- (1.460,0.060) -- (1.467,0.077) -- (1.473,0.095) -- (1.479,0.112) -- (1.484,0.130) -- (1.490,0.148) -- (1.495,0.167) -- (1.499,0.186) -- (1.504,0.206) -- (1.508,0.225) -- (1.511,0.245) -- (1.515,0.266) -- (1.518,0.286) -- (1.520,0.307) -- (1.522,0.329) -- (1.524,0.350) -- (1.525,0.372) -- (1.526,0.394) -- (1.527,0.417) -- (1.527,0.440) -- (1.527,0.463) -- (1.526,0.486) -- (1.525,0.509) -- (1.523,0.533) -- (1.521,0.557) -- (1.518,0.581) -- (1.515,0.606) -- (1.512,0.630) -- (1.507,0.655) -- (1.503,0.680) -- (1.498,0.705) -- (1.492,0.731) -- (1.486,0.756) -- (1.479,0.782) -- (1.472,0.807) -- (1.465,0.833) -- (1.456,0.859) -- (1.448,0.885) -- (1.438,0.911) -- (1.428,0.937) -- (1.418,0.963) -- (1.407,0.989) -- (1.396,1.015) -- (1.384,1.041) -- (1.371,1.068) -- (1.358,1.094) -- (1.345,1.120) -- (1.330,1.146) -- (1.316,1.172) -- (1.301,1.197) -- (1.285,1.223) -- (1.269,1.249) -- (1.252,1.274) -- (1.235,1.300) -- (1.217,1.325) -- (1.198,1.350) -- (1.180,1.375) -- (1.160,1.399) -- (1.141,1.424) -- (1.120,1.448) -- (1.100,1.472) -- (1.078,1.496) -- (1.057,1.519) -- (1.034,1.542) -- (1.012,1.565) -- (0.989,1.588) -- (0.965,1.610) -- (0.942,1.632) -- (0.917,1.653) -- (0.893,1.675) -- (0.868,1.696) -- (0.842,1.716) -- (0.816,1.736) -- (0.790,1.756) -- (0.764,1.775) -- (0.737,1.794) -- (0.710,1.813) -- (0.682,1.831) -- (0.655,1.848) -- (0.627,1.865) -- (0.598,1.882) -- (0.570,1.898) -- (0.541,1.914) -- (0.512,1.929) -- (0.483,1.944) -- (0.454,1.958) -- (0.424,1.972) -- (0.394,1.985) -- (0.365,1.998) -- (0.335,2.010) -- (0.304,2.021) -- (0.274,2.032) -- (0.244,2.043) -- (0.213,2.053) -- (0.183,2.062) -- (0.153,2.071) -- (0.122,2.079) -- (0.092,2.086) -- (0.061,2.093) -- (0.030,2.100) -- (0.000,2.106) -- (-0.030,2.111) -- (-0.061,2.116) -- (-0.091,2.120) -- (-0.121,2.124) -- (-0.151,2.127) -- (-0.181,2.129) -- (-0.211,2.131) -- (-0.241,2.132) -- (-0.270,2.133) -- (-0.300,2.133) -- (-0.329,2.132) -- (-0.358,2.131) -- (-0.386,2.129) -- (-0.415,2.127) -- (-0.443,2.124) -- (-0.471,2.121) -- (-0.498,2.117) -- (-0.526,2.113) -- (-0.553,2.108) -- (-0.580,2.102);
\draw[fill] (-0.580,2.102) -- (-0.606,2.096) -- (-0.632,2.089) -- (-0.658,2.082) -- (-0.683,2.074) -- (-0.708,2.066) -- (-0.733,2.057) -- (-0.757,2.048) -- (-0.781,2.039) -- (-0.804,2.028) -- (-0.827,2.018) -- (-0.850,2.007) -- (-0.872,1.995) -- (-0.893,1.983) -- (-0.915,1.970) -- (-0.935,1.957) -- (-0.956,1.944) -- (-0.976,1.930) -- (-0.995,1.916) -- (-1.014,1.902) -- (-1.032,1.887) -- (-1.050,1.871) -- (-1.068,1.856) -- (-1.085,1.839) -- (-1.101,1.823) -- (-1.117,1.806) -- (-1.132,1.789) -- (-1.147,1.772) -- (-1.162,1.754) -- (-1.175,1.736) -- (-1.189,1.718) -- (-1.202,1.699) -- (-1.214,1.680) -- (-1.226,1.661) -- (-1.237,1.642) -- (-1.248,1.622) -- (-1.258,1.603) -- (-1.268,1.583) -- (-1.278,1.563) -- (-1.286,1.542) -- (-1.295,1.522) -- (-1.302,1.501) -- (-1.310,1.480) -- (-1.317,1.459) -- (-1.323,1.438) -- (-1.329,1.417) -- (-1.334,1.395) -- (-1.339,1.374) -- (-1.344,1.352) -- (-1.348,1.330) -- (-1.351,1.308) -- (-1.355,1.287) -- (-1.357,1.265) -- (-1.360,1.243) -- (-1.362,1.221) -- (-1.363,1.199) -- (-1.364,1.176) -- (-1.365,1.154) -- (-1.365,1.132) -- (-1.365,1.110) -- (-1.365,1.088) -- (-1.364,1.066) -- (-1.363,1.044) -- (-1.362,1.022) -- (-1.360,0.999) -- (-1.358,0.977) -- (-1.356,0.955) -- (-1.353,0.933) -- (-1.350,0.912) -- (-1.347,0.890) -- (-1.343,0.868) -- (-1.340,0.846) -- (-1.336,0.824) -- (-1.332,0.803) -- (-1.327,0.781) -- (-1.323,0.760) -- (-1.318,0.739) -- (-1.313,0.717) -- (-1.308,0.696) -- (-1.302,0.675) -- (-1.297,0.654) -- (-1.291,0.633) -- (-1.286,0.612) -- (-1.280,0.591) -- (-1.274,0.571) -- (-1.268,0.550) -- (-1.262,0.530) -- (-1.255,0.510) -- (-1.249,0.489) -- (-1.243,0.469) -- (-1.236,0.449) -- (-1.230,0.429) -- (-1.223,0.409) -- (-1.217,0.390) -- (-1.210,0.370) -- (-1.204,0.350) -- (-1.197,0.331) -- (-1.191,0.312) -- (-1.184,0.292) -- (-1.178,0.273) -- (-1.172,0.254) -- (-1.165,0.235) -- (-1.159,0.216) -- (-1.153,0.197) -- (-1.146,0.178) -- (-1.140,0.159) -- (-1.134,0.141) -- (-1.128,0.122) -- (-1.122,0.104) -- (-1.116,0.085) -- (-1.111,0.067) -- (-1.105,0.048) -- (-1.099,0.030) -- (-1.094,0.011) -- (-1.089,-0.007) -- (-1.083,-0.025) -- (-1.078,-0.043) -- (-1.073,-0.062) -- (-1.068,-0.080) -- (-1.064,-0.098) -- (-1.059,-0.116) -- (-1.054,-0.134) -- (-1.050,-0.153);
\draw[fillback] (-1.050,-0.153) -- (-1.046,-0.171) -- (-1.041,-0.189) -- (-1.037,-0.207) -- (-1.033,-0.225) -- (-1.029,-0.244) -- (-1.026,-0.262) -- (-1.022,-0.280) -- (-1.018,-0.299) -- (-1.015,-0.317) -- (-1.012,-0.335) -- (-1.008,-0.354) -- (-1.005,-0.372) -- (-1.002,-0.391) -- (-0.999,-0.410) -- (-0.996,-0.428) -- (-0.993,-0.447) -- (-0.991,-0.466) -- (-0.988,-0.484) -- (-0.985,-0.503) -- (-0.982,-0.522) -- (-0.980,-0.541) -- (-0.977,-0.561) -- (-0.975,-0.580) -- (-0.972,-0.599) -- (-0.970,-0.618) -- (-0.967,-0.638) -- (-0.965,-0.657) -- (-0.962,-0.677) -- (-0.960,-0.696) -- (-0.957,-0.716) -- (-0.955,-0.736) -- (-0.952,-0.756) -- (-0.949,-0.776) -- (-0.946,-0.796) -- (-0.944,-0.816) -- (-0.941,-0.836) -- (-0.938,-0.856) -- (-0.934,-0.877) -- (-0.931,-0.897) -- (-0.928,-0.917) -- (-0.924,-0.938) -- (-0.921,-0.958) -- (-0.917,-0.979) -- (-0.913,-1.000) -- (-0.909,-1.020) -- (-0.905,-1.041) -- (-0.900,-1.062) -- (-0.895,-1.082) -- (-0.891,-1.103) -- (-0.885,-1.124) -- (-0.880,-1.145) -- (-0.874,-1.166) -- (-0.869,-1.187) -- (-0.863,-1.208) -- (-0.856,-1.228) -- (-0.850,-1.249) -- (-0.843,-1.270) -- (-0.836,-1.291) -- (-0.828,-1.312) -- (-0.820,-1.332) -- (-0.812,-1.353) -- (-0.804,-1.373) -- (-0.795,-1.394) -- (-0.786,-1.414) -- (-0.776,-1.435) -- (-0.767,-1.455) -- (-0.756,-1.475) -- (-0.746,-1.495) -- (-0.735,-1.515) -- (-0.724,-1.535) -- (-0.712,-1.554) -- (-0.700,-1.574) -- (-0.688,-1.593) -- (-0.675,-1.612) -- (-0.662,-1.631) -- (-0.648,-1.650) -- (-0.634,-1.668) -- (-0.620,-1.686) -- (-0.605,-1.704) -- (-0.590,-1.722) -- (-0.574,-1.740) -- (-0.558,-1.757) -- (-0.542,-1.774) -- (-0.525,-1.791) -- (-0.507,-1.807) -- (-0.490,-1.823) -- (-0.472,-1.839) -- (-0.453,-1.855) -- (-0.434,-1.870) -- (-0.415,-1.885) -- (-0.395,-1.899) -- (-0.375,-1.913) -- (-0.354,-1.927) -- (-0.333,-1.940) -- (-0.312,-1.953) -- (-0.290,-1.966) -- (-0.268,-1.978) -- (-0.245,-1.990) -- (-0.222,-2.001) -- (-0.199,-2.012) -- (-0.175,-2.022) -- (-0.151,-2.032) -- (-0.127,-2.041) -- (-0.102,-2.050) -- (-0.077,-2.059) -- (-0.052,-2.067) -- (-0.026,-2.074) -- (-0.000,-2.081) -- (0.026,-2.087) -- (0.053,-2.093) -- (0.080,-2.098) -- (0.107,-2.103) -- (0.134,-2.107) -- (0.162,-2.111) -- (0.190,-2.114) -- (0.218,-2.116) -- (0.247,-2.118) -- (0.275,-2.120) -- (0.304,-2.120) -- (0.333,-2.121) -- (0.362,-2.120) -- (0.391,-2.119) -- (0.421,-2.117) -- (0.450,-2.115) -- (0.480,-2.112) -- (0.510,-2.109) -- (0.540,-2.105) -- (0.570,-2.100) -- (0.600,-2.095) -- (0.630,-2.089) -- (0.660,-2.082) -- (0.690,-2.075) -- (0.720,-2.067) -- (0.750,-2.059);
\draw[fill] (0.750,-2.059) -- (0.780,-2.050) -- (0.810,-2.040) -- (0.840,-2.030) -- (0.870,-2.019) -- (0.900,-2.008) -- (0.929,-1.996) -- (0.959,-1.983) -- (0.988,-1.970) -- (1.017,-1.956) -- (1.047,-1.942) -- (1.075,-1.927) -- (1.104,-1.911) -- (1.133,-1.895) -- (1.161,-1.878) -- (1.189,-1.861) -- (1.217,-1.843) -- (1.244,-1.825) -- (1.271,-1.806) -- (1.298,-1.786) -- (1.325,-1.767) -- (1.351,-1.746) -- (1.377,-1.725) -- (1.403,-1.704) -- (1.428,-1.682) -- (1.453,-1.659) -- (1.477,-1.637) -- (1.501,-1.613) -- (1.525,-1.590) -- (1.548,-1.565) -- (1.571,-1.541) -- (1.593,-1.516) -- (1.615,-1.491) -- (1.636,-1.465) -- (1.657,-1.439) -- (1.678,-1.412) -- (1.698,-1.386) -- (1.717,-1.358) -- (1.736,-1.331) -- (1.754,-1.303) -- (1.772,-1.275) -- (1.789,-1.247) -- (1.806,-1.219) -- (1.822,-1.190) -- (1.838,-1.161) -- (1.853,-1.132) -- (1.867,-1.102) -- (1.881,-1.073) -- (1.894,-1.043) -- (1.907,-1.013) -- (1.919,-0.983) -- (1.931,-0.953) -- (1.942,-0.922) -- (1.952,-0.892) -- (1.961,-0.862) -- (1.971,-0.831) -- (1.979,-0.801) -- (1.987,-0.770) -- (1.994,-0.740) -- (2.001,-0.709) -- (2.007,-0.678) -- (2.012,-0.648) -- (2.017,-0.617) -- (2.021,-0.587) -- (2.024,-0.557) -- (2.027,-0.526) -- (2.030,-0.496) -- (2.032,-0.466) -- (2.033,-0.436) -- (2.033,-0.407) -- (2.033,-0.377) -- (2.033,-0.347) -- (2.031,-0.318) -- (2.030,-0.289) -- (2.027,-0.260) -- (2.024,-0.232) -- (2.021,-0.203) -- (2.017,-0.175) -- (2.012,-0.147) -- (2.007,-0.120) -- (2.002,-0.093) -- (1.996,-0.066) -- (1.989,-0.039) -- (1.982,-0.013) -- (1.975,0.014) -- (1.966,0.039) -- (1.958,0.065) -- (1.949,0.090) -- (1.940,0.114) -- (1.930,0.138) -- (1.919,0.162) -- (1.909,0.186) -- (1.898,0.209) -- (1.886,0.232) -- (1.874,0.254) -- (1.862,0.276) -- (1.849,0.297) -- (1.836,0.318) -- (1.823,0.339) -- (1.809,0.359) -- (1.796,0.379) -- (1.781,0.398) -- (1.767,0.417) -- (1.752,0.435) -- (1.737,0.453) -- (1.722,0.470) -- (1.706,0.487) -- (1.690,0.504) -- (1.674,0.520) -- (1.658,0.536) -- (1.642,0.551) -- (1.625,0.565) -- (1.609,0.580) -- (1.592,0.593) -- (1.575,0.607) -- (1.558,0.619) -- (1.541,0.632) -- (1.523,0.644) -- (1.506,0.655) -- (1.488,0.666) -- (1.471,0.677) -- (1.453,0.687) -- (1.436,0.697) -- (1.418,0.706) -- (1.400,0.715) -- (1.382,0.723) -- (1.365,0.731) -- (1.347,0.739) -- (1.329,0.746) -- (1.312,0.753) -- (1.294,0.759) -- (1.276,0.765) -- (1.259,0.771) -- (1.241,0.776) -- (1.224,0.781) -- (1.207,0.785) -- (1.189,0.790) -- (1.172,0.793) -- (1.155,0.797) -- (1.138,0.800) -- (1.121,0.803) -- (1.104,0.806) -- (1.088,0.808) -- (1.071,0.810) -- (1.055,0.812) -- (1.039,0.813) -- (1.023,0.814) -- (1.007,0.815) -- (0.991,0.816) -- (0.975,0.816) -- (0.960,0.817) -- (0.944,0.817) -- (0.929,0.817) -- (0.914,0.816) -- (0.899,0.816) -- (0.884,0.815) -- (0.870,0.815) -- (0.855,0.814) -- (0.841,0.813) -- (0.827,0.811) -- (0.813,0.810) -- (0.799,0.809) -- (0.785,0.807) -- (0.772,0.806) -- (0.759,0.804) -- (0.745,0.802) -- (0.732,0.801) -- (0.719,0.799) -- (0.707,0.797) -- (0.694,0.795) -- (0.681,0.793) -- (0.669,0.792) -- (0.657,0.790) -- (0.645,0.788) -- (0.632,0.786) -- (0.621,0.784) -- (0.609,0.783) -- (0.597,0.781) -- (0.585,0.779) -- (0.574,0.778) -- (0.562,0.776) -- (0.550,0.775) -- (0.539,0.774) -- (0.528,0.772) -- (0.516,0.771) -- (0.505,0.770) -- (0.494,0.769) -- (0.482,0.768) -- (0.471,0.768) -- (0.460,0.767) -- (0.449,0.767) -- (0.437,0.766) -- (0.426,0.766) -- (0.415,0.766) -- (0.403,0.766) -- (0.392,0.767) -- (0.380,0.767) -- (0.369,0.768) -- (0.357,0.768) -- (0.345,0.769) -- (0.333,0.770) -- (0.321,0.771) -- (0.309,0.773) -- (0.297,0.774) -- (0.285,0.776) -- (0.273,0.778) -- (0.260,0.780) -- (0.247,0.782) -- (0.234,0.784);
\draw[fillback] (0.234,0.784) -- (0.221,0.786) -- (0.208,0.789) -- (0.195,0.792) -- (0.181,0.795) -- (0.167,0.798) -- (0.153,0.801) -- (0.139,0.804) -- (0.124,0.808) -- (0.110,0.811) -- (0.095,0.815) -- (0.080,0.819) -- (0.064,0.823) -- (0.049,0.827) -- (0.033,0.831) -- (0.016,0.835) -- (0.000,0.840) -- (-0.017,0.844) -- (-0.034,0.849) -- (-0.051,0.853) -- (-0.068,0.858) -- (-0.086,0.863) -- (-0.104,0.867) -- (-0.123,0.872) -- (-0.141,0.877) -- (-0.160,0.882) -- (-0.180,0.887) -- (-0.199,0.892) -- (-0.219,0.897) -- (-0.239,0.902) -- (-0.260,0.907) -- (-0.280,0.912) -- (-0.301,0.917) -- (-0.323,0.922) -- (-0.344,0.927) -- (-0.366,0.932) -- (-0.388,0.936) -- (-0.411,0.941) -- (-0.434,0.946) -- (-0.457,0.950) -- (-0.480,0.955) -- (-0.503,0.959) -- (-0.527,0.963) -- (-0.551,0.967) -- (-0.575,0.971) -- (-0.600,0.975) -- (-0.625,0.979) -- (-0.650,0.983) -- (-0.675,0.986) -- (-0.701,0.989) -- (-0.726,0.992) -- (-0.752,0.995) -- (-0.778,0.998) -- (-0.804,1.000) -- (-0.831,1.003) -- (-0.857,1.005) -- (-0.884,1.006) -- (-0.911,1.008) -- (-0.938,1.009) -- (-0.965,1.010) -- (-0.992,1.011) -- (-1.020,1.012) -- (-1.047,1.012) -- (-1.075,1.012) -- (-1.102,1.012) -- (-1.130,1.011) -- (-1.158,1.010) -- (-1.186,1.009) -- (-1.213,1.008) -- (-1.241,1.006) -- (-1.269,1.004) -- (-1.296,1.002) -- (-1.324,0.999) -- (-1.352,0.996) -- (-1.379,0.993) -- (-1.407,0.989) -- (-1.434,0.985) -- (-1.462,0.980) -- (-1.489,0.976) -- (-1.516,0.971) -- (-1.543,0.965) -- (-1.569,0.960) -- (-1.596,0.953) -- (-1.622,0.947) -- (-1.648,0.940) -- (-1.674,0.933) -- (-1.700,0.926) -- (-1.725,0.918) -- (-1.751,0.910) -- (-1.776,0.901) -- (-1.800,0.892) -- (-1.824,0.883) -- (-1.848,0.873) -- (-1.872,0.863) -- (-1.895,0.853) -- (-1.918,0.843) -- (-1.941,0.832) -- (-1.963,0.821) -- (-1.985,0.809) -- (-2.006,0.797) -- (-2.027,0.785) -- (-2.048,0.772) -- (-2.068,0.760) -- (-2.087,0.747) -- (-2.106,0.733) -- (-2.125,0.720) -- (-2.143,0.706) -- (-2.161,0.691) -- (-2.178,0.677) -- (-2.194,0.662) -- (-2.210,0.647) -- (-2.226,0.632) -- (-2.241,0.616) -- (-2.255,0.600) -- (-2.269,0.585) -- (-2.282,0.568) -- (-2.295,0.552) -- (-2.307,0.535) -- (-2.318,0.519) -- (-2.329,0.502) -- (-2.339,0.484) -- (-2.349,0.467) -- (-2.358,0.450) -- (-2.366,0.432) -- (-2.373,0.414) -- (-2.380,0.396) -- (-2.387,0.378) -- (-2.392,0.360) -- (-2.398,0.342) -- (-2.402,0.324) -- (-2.406,0.305) -- (-2.409,0.287);
\draw[fill] (-2.409,0.287) -- (-2.411,0.268) -- (-2.413,0.250) -- (-2.414,0.231) -- (-2.414,0.213) -- (-2.414,0.194) -- (-2.413,0.175) -- (-2.412,0.157) -- (-2.409,0.138) -- (-2.406,0.119) -- (-2.403,0.101) -- (-2.399,0.082) -- (-2.394,0.064) -- (-2.388,0.045) -- (-2.382,0.027) -- (-2.376,0.009) -- (-2.368,-0.010) -- (-2.360,-0.028) -- (-2.352,-0.046) -- (-2.342,-0.064) -- (-2.333,-0.081) -- (-2.322,-0.099) -- (-2.311,-0.116) -- (-2.300,-0.134) -- (-2.287,-0.151) -- (-2.275,-0.168) -- (-2.261,-0.185) -- (-2.248,-0.202) -- (-2.233,-0.218) -- (-2.218,-0.234) -- (-2.203,-0.250) -- (-2.187,-0.266) -- (-2.171,-0.282) -- (-2.154,-0.297) -- (-2.137,-0.312) -- (-2.119,-0.327) -- (-2.101,-0.342) -- (-2.082,-0.357) -- (-2.063,-0.371) -- (-2.043,-0.385) -- (-2.023,-0.399) -- (-2.003,-0.412) -- (-1.983,-0.425) -- (-1.962,-0.438) -- (-1.940,-0.451) -- (-1.919,-0.463) -- (-1.897,-0.476) -- (-1.874,-0.487) -- (-1.852,-0.499) -- (-1.829,-0.510) -- (-1.806,-0.521) -- (-1.782,-0.532) -- (-1.759,-0.542) -- (-1.735,-0.553) -- (-1.711,-0.562) -- (-1.687,-0.572) -- (-1.662,-0.581) -- (-1.638,-0.590) -- (-1.613,-0.599) -- (-1.588,-0.608) -- (-1.563,-0.616) -- (-1.538,-0.624) -- (-1.513,-0.631) -- (-1.487,-0.639) -- (-1.462,-0.646) -- (-1.437,-0.653) -- (-1.411,-0.659) -- (-1.386,-0.666) -- (-1.360,-0.672) -- (-1.335,-0.677) -- (-1.309,-0.683) -- (-1.284,-0.688) -- (-1.258,-0.693) -- (-1.233,-0.698) -- (-1.207,-0.703) -- (-1.182,-0.707) -- (-1.157,-0.712) -- (-1.132,-0.716) -- (-1.107,-0.719) -- (-1.082,-0.723) -- (-1.057,-0.726) -- (-1.032,-0.730) -- (-1.008,-0.733) -- (-0.983,-0.735) -- (-0.959,-0.738) -- (-0.935,-0.741) -- (-0.911,-0.743) -- (-0.887,-0.745) -- (-0.863,-0.747) -- (-0.840,-0.749) -- (-0.816,-0.751) -- (-0.793,-0.753) -- (-0.771,-0.755) -- (-0.748,-0.756) -- (-0.726,-0.758) -- (-0.703,-0.759) -- (-0.681,-0.760) -- (-0.660,-0.762) -- (-0.638,-0.763) -- (-0.617,-0.764) -- (-0.596,-0.765) -- (-0.575,-0.766) -- (-0.554,-0.767) -- (-0.534,-0.768) -- (-0.514,-0.769) -- (-0.494,-0.770) -- (-0.474,-0.771) -- (-0.455,-0.772) -- (-0.436,-0.773) -- (-0.417,-0.774) -- (-0.398,-0.775) -- (-0.380,-0.776) -- (-0.362,-0.777) -- (-0.344,-0.778) -- (-0.326,-0.779) -- (-0.308,-0.780) -- (-0.291,-0.782) -- (-0.274,-0.783) -- (-0.257,-0.784) -- (-0.241,-0.786) -- (-0.224,-0.787) -- (-0.208,-0.789) -- (-0.192,-0.791) -- (-0.176,-0.793) -- (-0.161,-0.795);
\draw[fillback] (-0.161,-0.795) -- (-0.145,-0.797) -- (-0.130,-0.799) -- (-0.115,-0.801) -- (-0.100,-0.804) -- (-0.085,-0.806) -- (-0.071,-0.809) -- (-0.056,-0.812) -- (-0.042,-0.814) -- (-0.028,-0.817) -- (-0.014,-0.821) -- (0.000,-0.824) -- (0.014,-0.827) -- (0.028,-0.831) -- (0.041,-0.834) -- (0.055,-0.838) -- (0.068,-0.842) -- (0.082,-0.846) -- (0.095,-0.850) -- (0.108,-0.854) -- (0.121,-0.859) -- (0.135,-0.863) -- (0.148,-0.868) -- (0.161,-0.872) -- (0.174,-0.877) -- (0.187,-0.882) -- (0.201,-0.887) -- (0.214,-0.892) -- (0.227,-0.897) -- (0.240,-0.903) -- (0.253,-0.908) -- (0.267,-0.914) -- (0.280,-0.919) -- (0.294,-0.925) -- (0.307,-0.930) -- (0.321,-0.936) -- (0.334,-0.942) -- (0.348,-0.947) -- (0.362,-0.953) -- (0.376,-0.959) -- (0.390,-0.965) -- (0.405,-0.971) -- (0.419,-0.977) -- (0.433,-0.982) -- (0.448,-0.988) -- (0.463,-0.994) -- (0.478,-1.000) -- (0.493,-1.006) -- (0.508,-1.011) -- (0.524,-1.017) -- (0.539,-1.022) -- (0.555,-1.028) -- (0.571,-1.033) -- (0.587,-1.039) -- (0.603,-1.044) -- (0.620,-1.049) -- (0.637,-1.054) -- (0.653,-1.058) -- (0.671,-1.063) -- (0.688,-1.068) -- (0.705,-1.072) -- (0.723,-1.076) -- (0.741,-1.080) -- (0.759,-1.084) -- (0.777,-1.087) -- (0.795,-1.090) -- (0.814,-1.093) -- (0.833,-1.096) -- (0.852,-1.099) -- (0.871,-1.101) -- (0.890,-1.103) -- (0.910,-1.105) -- (0.929,-1.106) -- (0.949,-1.107) -- (0.969,-1.108) -- (0.990,-1.108) -- (1.010,-1.108) -- (1.030,-1.108) -- (1.051,-1.108) -- (1.072,-1.107) -- (1.093,-1.105) -- (1.114,-1.104) -- (1.135,-1.102) -- (1.156,-1.099) -- (1.178,-1.096) -- (1.199,-1.093) -- (1.221,-1.089) -- (1.242,-1.085) -- (1.264,-1.080) -- (1.286,-1.075) -- (1.308,-1.070) -- (1.329,-1.064) -- (1.351,-1.057) -- (1.373,-1.050) -- (1.395,-1.043) -- (1.417,-1.035) -- (1.439,-1.027) -- (1.461,-1.018) -- (1.483,-1.009) -- (1.505,-0.999) -- (1.526,-0.988) -- (1.548,-0.978) -- (1.570,-0.966) -- (1.591,-0.954) -- (1.613,-0.942) -- (1.634,-0.929) -- (1.655,-0.916) -- (1.676,-0.902) -- (1.697,-0.887) -- (1.718,-0.872) -- (1.738,-0.857) -- (1.758,-0.841) -- (1.778,-0.825) -- (1.798,-0.808) -- (1.818,-0.790) -- (1.837,-0.772) -- (1.856,-0.754) -- (1.875,-0.735) -- (1.894,-0.715) -- (1.912,-0.695) -- (1.930,-0.675) -- (1.947,-0.654) -- (1.965,-0.632) -- (1.982,-0.610) -- (1.998,-0.588) -- (2.014,-0.565) -- (2.030,-0.542) -- (2.045,-0.518) -- (2.060,-0.494) -- (2.075,-0.469) -- (2.089,-0.444) -- (2.103,-0.419) -- (2.116,-0.393) -- (2.129,-0.367) -- (2.141,-0.340) -- (2.153,-0.313) -- (2.164,-0.286) -- (2.174,-0.258) -- (2.185,-0.230) -- (2.194,-0.202) -- (2.203,-0.173) -- (2.212,-0.144) -- (2.220,-0.115) -- (2.227,-0.085) -- (2.234,-0.056) -- (2.241,-0.025) -- (2.246,0.005) -- (2.251,0.035) -- (2.256,0.066) -- (2.260,0.097) -- (2.263,0.128) -- (2.266,0.160) -- (2.268,0.191) -- (2.269,0.223) -- (2.270,0.254) -- (2.270,0.286) -- (2.270,0.318) -- (2.269,0.350) -- (2.267,0.382) -- (2.265,0.415) -- (2.262,0.447) -- (2.258,0.479) -- (2.254,0.511) -- (2.249,0.543) -- (2.243,0.576) -- (2.237,0.608) -- (2.230,0.640) -- (2.222,0.672) -- (2.214,0.704) -- (2.205,0.736) -- (2.196,0.767) -- (2.186,0.799) -- (2.175,0.830) -- (2.164,0.862) -- (2.152,0.893) -- (2.139,0.923) -- (2.126,0.954) -- (2.113,0.985) -- (2.098,1.015) -- (2.083,1.045) -- (2.068,1.074) -- (2.052,1.104) -- (2.035,1.133) -- (2.018,1.161) -- (2.000,1.190) -- (1.982,1.218) -- (1.963,1.246) -- (1.944,1.273) -- (1.924,1.300) -- (1.904,1.326) -- (1.883,1.353) -- (1.862,1.378) -- (1.840,1.404) -- (1.818,1.428) -- (1.795,1.453) -- (1.772,1.477) -- (1.749,1.500) -- (1.725,1.523) -- (1.701,1.546) -- (1.676,1.568) -- (1.651,1.589) -- (1.626,1.610) -- (1.600,1.631) -- (1.574,1.651);
\draw[fill] (1.574,1.651) -- (1.548,1.670) -- (1.521,1.689) -- (1.495,1.707) -- (1.467,1.725) -- (1.440,1.742) -- (1.412,1.758) -- (1.385,1.774) -- (1.357,1.789) -- (1.328,1.804) -- (1.300,1.818) -- (1.272,1.832) -- (1.243,1.845) -- (1.214,1.857) -- (1.185,1.869) -- (1.156,1.880) -- (1.127,1.891) -- (1.098,1.900) -- (1.069,1.910) -- (1.039,1.918) -- (1.010,1.926) -- (0.981,1.934) -- (0.951,1.941) -- (0.922,1.947) -- (0.893,1.953) -- (0.864,1.958) -- (0.834,1.962) -- (0.805,1.966) -- (0.776,1.969) -- (0.747,1.971) -- (0.719,1.973) -- (0.690,1.975) -- (0.661,1.976) -- (0.633,1.976) -- (0.605,1.976) -- (0.577,1.975) -- (0.549,1.973) -- (0.521,1.971) -- (0.494,1.969) -- (0.466,1.965) -- (0.439,1.962) -- (0.412,1.957) -- (0.386,1.953) -- (0.360,1.947) -- (0.333,1.942) -- (0.308,1.935) -- (0.282,1.929) -- (0.257,1.921) -- (0.232,1.914) -- (0.207,1.906) -- (0.183,1.897) -- (0.159,1.888) -- (0.135,1.878) -- (0.112,1.868) -- (0.089,1.858) -- (0.066,1.847) -- (0.044,1.836) -- (0.022,1.825) -- (-0.000,1.813) -- (-0.021,1.800) -- (-0.042,1.788) -- (-0.063,1.775) -- (-0.083,1.761) -- (-0.103,1.748) -- (-0.123,1.734) -- (-0.142,1.720) -- (-0.161,1.705) -- (-0.179,1.690) -- (-0.198,1.675) -- (-0.215,1.660) -- (-0.233,1.644) -- (-0.250,1.628) -- (-0.267,1.612) -- (-0.283,1.596) -- (-0.299,1.580) -- (-0.315,1.563) -- (-0.330,1.546) -- (-0.345,1.529) -- (-0.360,1.512) -- (-0.374,1.495) -- (-0.388,1.477) -- (-0.402,1.459) -- (-0.416,1.442) -- (-0.429,1.424) -- (-0.442,1.406) -- (-0.454,1.388) -- (-0.466,1.370) -- (-0.478,1.351) -- (-0.490,1.333) -- (-0.501,1.315) -- (-0.512,1.296) -- (-0.523,1.278) -- (-0.534,1.259) -- (-0.544,1.241) -- (-0.555,1.222) -- (-0.565,1.204) -- (-0.574,1.185) -- (-0.584,1.167) -- (-0.593,1.148) -- (-0.602,1.130) -- (-0.611,1.111) -- (-0.620,1.093) -- (-0.629,1.074) -- (-0.637,1.056) -- (-0.646,1.037) -- (-0.654,1.019) -- (-0.662,1.001) -- (-0.670,0.982) -- (-0.678,0.964) -- (-0.686,0.946) -- (-0.693,0.928) -- (-0.701,0.910) -- (-0.709,0.892) -- (-0.716,0.874) -- (-0.723,0.856) -- (-0.731,0.838) -- (-0.738,0.820) -- (-0.745,0.803) -- (-0.753,0.785) -- (-0.760,0.768) -- (-0.767,0.750) -- (-0.775,0.733) -- (-0.782,0.716) -- (-0.789,0.698) -- (-0.797,0.681) -- (-0.804,0.664) -- (-0.812,0.647) -- (-0.819,0.630) -- (-0.827,0.614) -- (-0.834,0.597) -- (-0.842,0.580) -- (-0.850,0.563) -- (-0.858,0.547) -- (-0.866,0.530) -- (-0.874,0.514) -- (-0.882,0.498) -- (-0.890,0.481) -- (-0.898,0.465) -- (-0.907,0.449) -- (-0.915,0.432) -- (-0.924,0.416) -- (-0.933,0.400) -- (-0.942,0.384) -- (-0.951,0.368) -- (-0.960,0.352) -- (-0.969,0.335) -- (-0.979,0.319) -- (-0.988,0.303);
\draw[fillback] (-0.988,0.303) -- (-0.998,0.287) -- (-1.008,0.271) -- (-1.018,0.255) -- (-1.028,0.239) -- (-1.038,0.223) -- (-1.049,0.207) -- (-1.059,0.191) -- (-1.070,0.174) -- (-1.080,0.158) -- (-1.091,0.142) -- (-1.102,0.126) -- (-1.113,0.109) -- (-1.124,0.093) -- (-1.135,0.076) -- (-1.147,0.060) -- (-1.158,0.043) -- (-1.170,0.027) -- (-1.181,0.010) -- (-1.193,-0.007) -- (-1.205,-0.024) -- (-1.217,-0.041) -- (-1.229,-0.058) -- (-1.241,-0.075) -- (-1.253,-0.092) -- (-1.265,-0.110) -- (-1.277,-0.127) -- (-1.289,-0.145) -- (-1.301,-0.162) -- (-1.313,-0.180) -- (-1.325,-0.198) -- (-1.337,-0.216) -- (-1.349,-0.234) -- (-1.361,-0.252) -- (-1.373,-0.271) -- (-1.385,-0.289) -- (-1.397,-0.308) -- (-1.409,-0.327) -- (-1.420,-0.346) -- (-1.432,-0.365) -- (-1.444,-0.384) -- (-1.455,-0.403) -- (-1.466,-0.422) -- (-1.477,-0.442) -- (-1.488,-0.462) -- (-1.499,-0.481) -- (-1.510,-0.501) -- (-1.520,-0.521) -- (-1.530,-0.541) -- (-1.540,-0.561) -- (-1.550,-0.582) -- (-1.560,-0.602) -- (-1.569,-0.623) -- (-1.578,-0.643) -- (-1.587,-0.664) -- (-1.595,-0.685) -- (-1.604,-0.706) -- (-1.612,-0.727) -- (-1.619,-0.748) -- (-1.626,-0.769) -- (-1.633,-0.791) -- (-1.640,-0.812) -- (-1.646,-0.833) -- (-1.652,-0.855) -- (-1.658,-0.876) -- (-1.663,-0.898) -- (-1.668,-0.920) -- (-1.672,-0.941) -- (-1.676,-0.963) -- (-1.679,-0.984) -- (-1.682,-1.006) -- (-1.685,-1.028) -- (-1.687,-1.049) -- (-1.689,-1.071) -- (-1.690,-1.093) -- (-1.691,-1.114) -- (-1.691,-1.136) -- (-1.691,-1.157) -- (-1.690,-1.179) -- (-1.689,-1.200) -- (-1.687,-1.222) -- (-1.685,-1.243) -- (-1.682,-1.264) -- (-1.679,-1.285) -- (-1.675,-1.306) -- (-1.670,-1.327) -- (-1.665,-1.347) -- (-1.660,-1.368) -- (-1.654,-1.388) -- (-1.647,-1.408) -- (-1.640,-1.428) -- (-1.632,-1.448) -- (-1.624,-1.468) -- (-1.615,-1.487) -- (-1.606,-1.506) -- (-1.596,-1.525) -- (-1.585,-1.544) -- (-1.574,-1.562) -- (-1.562,-1.581) -- (-1.550,-1.598) -- (-1.537,-1.616) -- (-1.524,-1.633) -- (-1.510,-1.651) -- (-1.496,-1.667) -- (-1.481,-1.684) -- (-1.465,-1.700) -- (-1.449,-1.716) -- (-1.432,-1.731) -- (-1.415,-1.746) -- (-1.397,-1.761) -- (-1.379,-1.775) -- (-1.360,-1.789) -- (-1.341,-1.802) -- (-1.321,-1.815) -- (-1.301,-1.828) -- (-1.280,-1.840) -- (-1.259,-1.852) -- (-1.237,-1.864);
\draw[fill] (-1.237,-1.864) -- (-1.215,-1.875) -- (-1.193,-1.885) -- (-1.170,-1.895) -- (-1.146,-1.905) -- (-1.122,-1.914) -- (-1.098,-1.922) -- (-1.073,-1.930) -- (-1.048,-1.938) -- (-1.023,-1.945) -- (-0.997,-1.952) -- (-0.971,-1.958) -- (-0.944,-1.964) -- (-0.917,-1.969) -- (-0.890,-1.973) -- (-0.863,-1.977) -- (-0.835,-1.981) -- (-0.807,-1.984) -- (-0.779,-1.986) -- (-0.750,-1.988) -- (-0.721,-1.989) -- (-0.692,-1.990) -- (-0.663,-1.990) -- (-0.634,-1.990) -- (-0.604,-1.989) -- (-0.575,-1.988) -- (-0.545,-1.986) -- (-0.515,-1.984) -- (-0.485,-1.981) -- (-0.454,-1.977) -- (-0.424,-1.973) -- (-0.394,-1.968) -- (-0.363,-1.963) -- (-0.333,-1.958) -- (-0.302,-1.951) -- (-0.272,-1.945) -- (-0.242,-1.937) -- (-0.211,-1.930) -- (-0.181,-1.921) -- (-0.150,-1.912) -- (-0.120,-1.903) -- (-0.090,-1.893) -- (-0.060,-1.883) -- (-0.030,-1.872) -- (0.000,-1.861) -- (0.030,-1.849) -- (0.059,-1.837) -- (0.089,-1.824) -- (0.118,-1.811) -- (0.147,-1.798) -- (0.176,-1.784) -- (0.204,-1.769) -- (0.232,-1.754) -- (0.260,-1.739) -- (0.288,-1.723) -- (0.316,-1.707) -- (0.343,-1.691) -- (0.370,-1.674) -- (0.397,-1.657) -- (0.423,-1.640) -- (0.449,-1.622) -- (0.474,-1.604) -- (0.500,-1.585) -- (0.525,-1.567) -- (0.549,-1.547) -- (0.573,-1.528) -- (0.597,-1.509) -- (0.620,-1.489) -- (0.643,-1.469) -- (0.666,-1.448) -- (0.688,-1.428) -- (0.710,-1.407) -- (0.731,-1.386) -- (0.752,-1.365) -- (0.772,-1.344) -- (0.792,-1.322) -- (0.811,-1.301) -- (0.830,-1.279) -- (0.849,-1.257) -- (0.867,-1.235) -- (0.884,-1.213) -- (0.902,-1.191) -- (0.918,-1.169) -- (0.934,-1.147) -- (0.950,-1.125) -- (0.965,-1.103) -- (0.980,-1.080) -- (0.994,-1.058) -- (1.008,-1.036) -- (1.021,-1.014) -- (1.034,-0.992) -- (1.046,-0.969) -- (1.058,-0.947) -- (1.069,-0.925) -- (1.080,-0.903) -- (1.091,-0.882) -- (1.100,-0.860) -- (1.110,-0.838) -- (1.119,-0.817) -- (1.127,-0.796) -- (1.136,-0.774) -- (1.143,-0.753) -- (1.150,-0.733) -- (1.157,-0.712) -- (1.164,-0.691) -- (1.169,-0.671) -- (1.175,-0.651) -- (1.180,-0.631) -- (1.185,-0.612) -- (1.189,-0.592) -- (1.193,-0.573) -- (1.196,-0.554) -- (1.200,-0.535) -- (1.202,-0.517) -- (1.205,-0.499) -- (1.207,-0.481) -- (1.209,-0.463) -- (1.210,-0.446) -- (1.211,-0.429) -- (1.212,-0.412) -- (1.213,-0.396) -- (1.213,-0.379) -- (1.213,-0.364) -- (1.213,-0.348) -- (1.212,-0.333) -- (1.211,-0.318) -- (1.210,-0.303) -- (1.209,-0.289) -- (1.207,-0.275) -- (1.206,-0.261) -- (1.204,-0.248) -- (1.202,-0.235) -- (1.200,-0.222) -- (1.197,-0.210) -- (1.195,-0.197) -- (1.192,-0.186) -- (1.189,-0.174) -- (1.186,-0.163) -- (1.183,-0.152) -- (1.180,-0.141) -- (1.177,-0.131) -- (1.173,-0.121) -- (1.170,-0.111) -- (1.166,-0.102) -- (1.163,-0.093) -- (1.159,-0.084) -- (1.156,-0.075) -- (1.152,-0.067) -- (1.148,-0.059) -- (1.145,-0.051) -- (1.141,-0.044) -- (1.137,-0.036) -- (1.134,-0.029) -- (1.130,-0.022) -- (1.127,-0.016) -- (1.123,-0.009) -- (1.119,-0.003) -- (1.116,0.003) -- (1.113,0.009) -- (1.109,0.014) -- (1.106,0.020) -- (1.103,0.025) -- (1.100,0.030) -- (1.097,0.035) -- (1.094,0.040) -- (1.091,0.045) -- (1.088,0.049) -- (1.085,0.054) -- (1.083,0.058) -- (1.080,0.063) -- (1.078,0.067) -- (1.076,0.071) -- (1.074,0.075) -- (1.072,0.079) -- (1.070,0.084) -- (1.068,0.088) -- (1.066,0.092) -- (1.065,0.096) -- (1.063,0.100) -- (1.062,0.104) -- (1.061,0.108) -- (1.059,0.112) -- (1.058,0.116) -- (1.057,0.120);
\draw[fillback] (1.057,0.120) -- (1.057,0.124) -- (1.056,0.129) -- (1.055,0.133) -- (1.055,0.137) -- (1.054,0.142) -- (1.054,0.147) -- (1.054,0.152) -- (1.053,0.157) -- (1.053,0.162) -- (1.053,0.167) -- (1.053,0.172) -- (1.053,0.178) -- (1.053,0.184) -- (1.053,0.190) -- (1.054,0.196) -- (1.054,0.202) -- (1.054,0.208) -- (1.054,0.215) -- (1.054,0.222) -- (1.054,0.229) -- (1.055,0.237) -- (1.055,0.244) -- (1.055,0.252) -- (1.055,0.260) -- (1.055,0.269) -- (1.055,0.277) -- (1.055,0.286) -- (1.055,0.295) -- (1.054,0.304) -- (1.054,0.314) -- (1.054,0.324) -- (1.053,0.334) -- (1.052,0.345) -- (1.051,0.355) -- (1.050,0.366) -- (1.049,0.378) -- (1.048,0.389) -- (1.046,0.401) -- (1.045,0.413) -- (1.043,0.426) -- (1.041,0.438) -- (1.038,0.451) -- (1.036,0.465) -- (1.033,0.478) -- (1.030,0.492) -- (1.027,0.506) -- (1.023,0.520) -- (1.020,0.535) -- (1.015,0.550) -- (1.011,0.565) -- (1.006,0.580) -- (1.001,0.596) -- (0.996,0.612) -- (0.990,0.628) -- (0.984,0.644) -- (0.978,0.661) -- (0.971,0.678) -- (0.964,0.695) -- (0.957,0.712) -- (0.949,0.729) -- (0.941,0.747) -- (0.933,0.765) -- (0.924,0.783) -- (0.914,0.801) -- (0.904,0.819) -- (0.894,0.837) -- (0.884,0.856) -- (0.873,0.875) -- (0.861,0.893) -- (0.850,0.912) -- (0.837,0.931) -- (0.825,0.950) -- (0.811,0.970) -- (0.798,0.989) -- (0.784,1.008) -- (0.769,1.028) -- (0.755,1.047) -- (0.739,1.066) -- (0.723,1.086) -- (0.707,1.105) -- (0.690,1.125) -- (0.673,1.144) -- (0.656,1.164) -- (0.638,1.183) -- (0.619,1.202) -- (0.600,1.221) -- (0.581,1.240) -- (0.561,1.260) -- (0.541,1.278) -- (0.520,1.297) -- (0.499,1.316) -- (0.478,1.334) -- (0.456,1.353) -- (0.434,1.371) -- (0.411,1.389) -- (0.388,1.407) -- (0.365,1.425) -- (0.341,1.442) -- (0.316,1.459) -- (0.292,1.476) -- (0.267,1.493) -- (0.242,1.509) -- (0.216,1.526) -- (0.190,1.541) -- (0.164,1.557) -- (0.137,1.572) -- (0.110,1.587) -- (0.083,1.602) -- (0.056,1.616) -- (0.028,1.630) -- (-0.000,1.644) -- (-0.028,1.657) -- (-0.057,1.670) -- (-0.085,1.683) -- (-0.114,1.695) -- (-0.143,1.706) -- (-0.173,1.718) -- (-0.202,1.728) -- (-0.231,1.739) -- (-0.261,1.749) -- (-0.291,1.759) -- (-0.321,1.768) -- (-0.351,1.776) -- (-0.381,1.785) -- (-0.411,1.792) -- (-0.442,1.800) -- (-0.472,1.806) -- (-0.502,1.813) -- (-0.532,1.819) -- (-0.563,1.824) -- (-0.593,1.829) -- (-0.623,1.833) -- (-0.653,1.837) -- (-0.684,1.840) -- (-0.714,1.843) -- (-0.743,1.846) -- (-0.773,1.847) -- (-0.803,1.849) -- (-0.833,1.849) -- (-0.862,1.850) -- (-0.891,1.850) -- (-0.920,1.849) -- (-0.949,1.848) -- (-0.978,1.846) -- (-1.006,1.844) -- (-1.034,1.841) -- (-1.062,1.837) -- (-1.089,1.834) -- (-1.117,1.829) -- (-1.144,1.825) -- (-1.170,1.819) -- (-1.197,1.813) -- (-1.223,1.807) -- (-1.248,1.800) -- (-1.274,1.793) -- (-1.298,1.785) -- (-1.323,1.777) -- (-1.347,1.769) -- (-1.371,1.759) -- (-1.394,1.750) -- (-1.417,1.740) -- (-1.439,1.729) -- (-1.461,1.718) -- (-1.482,1.707) -- (-1.503,1.695) -- (-1.524,1.683) -- (-1.544,1.671) -- (-1.563,1.658) -- (-1.582,1.644) -- (-1.600,1.631) -- (-1.618,1.616);
\draw[fill] (-1.618,1.616) -- (-1.635,1.602) -- (-1.652,1.587) -- (-1.668,1.572) -- (-1.684,1.556) -- (-1.699,1.540) -- (-1.714,1.524) -- (-1.728,1.508) -- (-1.741,1.491) -- (-1.754,1.474) -- (-1.766,1.457) -- (-1.778,1.439) -- (-1.789,1.421) -- (-1.799,1.403) -- (-1.809,1.385) -- (-1.818,1.366) -- (-1.827,1.347) -- (-1.835,1.328) -- (-1.843,1.309) -- (-1.850,1.289) -- (-1.856,1.270) -- (-1.862,1.250) -- (-1.867,1.230) -- (-1.871,1.210) -- (-1.875,1.190) -- (-1.879,1.170) -- (-1.881,1.149) -- (-1.884,1.129) -- (-1.885,1.108) -- (-1.886,1.087) -- (-1.887,1.066) -- (-1.887,1.046) -- (-1.886,1.025) -- (-1.885,1.004) -- (-1.884,0.983) -- (-1.882,0.962) -- (-1.879,0.941) -- (-1.876,0.920) -- (-1.872,0.898) -- (-1.868,0.877) -- (-1.863,0.856) -- (-1.858,0.835) -- (-1.852,0.814) -- (-1.846,0.793) -- (-1.840,0.772) -- (-1.833,0.751) -- (-1.826,0.731) -- (-1.818,0.710) -- (-1.809,0.689) -- (-1.801,0.669) -- (-1.792,0.648) -- (-1.782,0.628) -- (-1.773,0.607) -- (-1.763,0.587) -- (-1.752,0.567) -- (-1.741,0.547) -- (-1.730,0.527) -- (-1.719,0.507) -- (-1.707,0.488) -- (-1.695,0.468) -- (-1.683,0.449) -- (-1.670,0.429) -- (-1.657,0.410) -- (-1.644,0.391) -- (-1.631,0.372) -- (-1.618,0.353) -- (-1.604,0.335) -- (-1.590,0.316) -- (-1.576,0.298) -- (-1.562,0.280) -- (-1.548,0.262) -- (-1.533,0.244) -- (-1.519,0.226) -- (-1.504,0.209) -- (-1.489,0.191) -- (-1.474,0.174) -- (-1.459,0.157) -- (-1.444,0.140) -- (-1.429,0.123) -- (-1.414,0.106) -- (-1.398,0.090) -- (-1.383,0.073) -- (-1.368,0.057) -- (-1.352,0.041) -- (-1.337,0.025) -- (-1.322,0.009) -- (-1.307,-0.007) -- (-1.291,-0.022) -- (-1.276,-0.038) -- (-1.261,-0.053) -- (-1.246,-0.069) -- (-1.231,-0.084) -- (-1.216,-0.099) -- (-1.201,-0.114) -- (-1.186,-0.129) -- (-1.171,-0.144) -- (-1.156,-0.158) -- (-1.142,-0.173) -- (-1.127,-0.188) -- (-1.113,-0.202) -- (-1.099,-0.217) -- (-1.085,-0.231) -- (-1.071,-0.245) -- (-1.057,-0.259) -- (-1.043,-0.274) -- (-1.029,-0.288) -- (-1.016,-0.302) -- (-1.002,-0.316) -- (-0.989,-0.330) -- (-0.976,-0.344) -- (-0.963,-0.358) -- (-0.950,-0.372) -- (-0.938,-0.386) -- (-0.925,-0.400) -- (-0.913,-0.414) -- (-0.901,-0.428) -- (-0.889,-0.443);
\draw[fillback] (-0.889,-0.443) -- (-0.877,-0.457) -- (-0.865,-0.471) -- (-0.853,-0.485) -- (-0.842,-0.499) -- (-0.830,-0.513) -- (-0.819,-0.527) -- (-0.808,-0.542) -- (-0.797,-0.556) -- (-0.786,-0.570) -- (-0.775,-0.585) -- (-0.765,-0.599) -- (-0.754,-0.614) -- (-0.743,-0.628) -- (-0.733,-0.643) -- (-0.723,-0.658) -- (-0.713,-0.672) -- (-0.703,-0.687) -- (-0.693,-0.702) -- (-0.683,-0.717) -- (-0.673,-0.732) -- (-0.663,-0.747) -- (-0.653,-0.762) -- (-0.643,-0.778) -- (-0.633,-0.793) -- (-0.624,-0.809) -- (-0.614,-0.824) -- (-0.604,-0.840) -- (-0.595,-0.855) -- (-0.585,-0.871) -- (-0.575,-0.887) -- (-0.565,-0.903) -- (-0.555,-0.919) -- (-0.546,-0.935) -- (-0.536,-0.951) -- (-0.526,-0.967) -- (-0.516,-0.983) -- (-0.506,-1.000) -- (-0.495,-1.016) -- (-0.485,-1.032) -- (-0.475,-1.049) -- (-0.464,-1.065) -- (-0.453,-1.082) -- (-0.443,-1.098) -- (-0.432,-1.115) -- (-0.421,-1.132) -- (-0.409,-1.148) -- (-0.398,-1.165) -- (-0.386,-1.181) -- (-0.375,-1.198) -- (-0.363,-1.215) -- (-0.350,-1.231) -- (-0.338,-1.248) -- (-0.325,-1.264) -- (-0.313,-1.281) -- (-0.300,-1.297) -- (-0.286,-1.314) -- (-0.273,-1.330) -- (-0.259,-1.346) -- (-0.245,-1.362) -- (-0.231,-1.378) -- (-0.216,-1.394) -- (-0.201,-1.410) -- (-0.186,-1.426) -- (-0.171,-1.442) -- (-0.155,-1.457) -- (-0.139,-1.472) -- (-0.123,-1.488) -- (-0.106,-1.503) -- (-0.089,-1.517) -- (-0.072,-1.532) -- (-0.054,-1.547) -- (-0.037,-1.561) -- (-0.018,-1.575) -- (-0.000,-1.589) -- (0.019,-1.602) -- (0.038,-1.615) -- (0.057,-1.628) -- (0.077,-1.641) -- (0.097,-1.654) -- (0.118,-1.666) -- (0.138,-1.678) -- (0.159,-1.689) -- (0.181,-1.700) -- (0.202,-1.711) -- (0.224,-1.722) -- (0.247,-1.732) -- (0.269,-1.742) -- (0.292,-1.751) -- (0.315,-1.760) -- (0.339,-1.769) -- (0.363,-1.777) -- (0.387,-1.785) -- (0.411,-1.792) -- (0.436,-1.799) -- (0.461,-1.806) -- (0.486,-1.812) -- (0.511,-1.817) -- (0.537,-1.822) -- (0.563,-1.827) -- (0.589,-1.831) -- (0.615,-1.834) -- (0.642,-1.838) -- (0.668,-1.840) -- (0.695,-1.842) -- (0.722,-1.844) -- (0.750,-1.845) -- (0.777,-1.845) -- (0.805,-1.845) -- (0.833,-1.844) -- (0.861,-1.843) -- (0.889,-1.841) -- (0.917,-1.839) -- (0.945,-1.836) -- (0.973,-1.832) -- (1.002,-1.828) -- (1.030,-1.823) -- (1.059,-1.818) -- (1.087,-1.812) -- (1.116,-1.805) -- (1.144,-1.798) -- (1.173,-1.791) -- (1.201,-1.782) -- (1.230,-1.773) -- (1.258,-1.764) -- (1.287,-1.754) -- (1.315,-1.743) -- (1.343,-1.731) -- (1.371,-1.719) -- (1.399,-1.707) -- (1.427,-1.694) -- (1.455,-1.680) -- (1.482,-1.665) -- (1.510,-1.651) -- (1.537,-1.635) -- (1.564,-1.619) -- (1.591,-1.602) -- (1.617,-1.585) -- (1.643,-1.567) -- (1.669,-1.548) -- (1.695,-1.529) -- (1.720,-1.510) -- (1.745,-1.490) -- (1.770,-1.469) -- (1.795,-1.448) -- (1.819,-1.426) -- (1.842,-1.404) -- (1.866,-1.381) -- (1.889,-1.358) -- (1.911,-1.334) -- (1.933,-1.310) -- (1.955,-1.285) -- (1.976,-1.260) -- (1.997,-1.235) -- (2.017,-1.209) -- (2.037,-1.182) -- (2.056,-1.155);
\draw[fill] (2.056,-1.155) -- (2.075,-1.128) -- (2.094,-1.100) -- (2.111,-1.072) -- (2.129,-1.044) -- (2.145,-1.015) -- (2.162,-0.986) -- (2.177,-0.957) -- (2.192,-0.927) -- (2.207,-0.897) -- (2.221,-0.867) -- (2.234,-0.836) -- (2.247,-0.805) -- (2.259,-0.774) -- (2.270,-0.743) -- (2.281,-0.711) -- (2.291,-0.680) -- (2.301,-0.648) -- (2.310,-0.616) -- (2.318,-0.583) -- (2.326,-0.551) -- (2.333,-0.519) -- (2.339,-0.486) -- (2.345,-0.453) -- (2.350,-0.421) -- (2.355,-0.388) -- (2.359,-0.355) -- (2.362,-0.322) -- (2.364,-0.289) -- (2.366,-0.256) -- (2.367,-0.224) -- (2.368,-0.191) -- (2.368,-0.158) -- (2.367,-0.126) -- (2.365,-0.093) -- (2.363,-0.061) -- (2.361,-0.028) -- (2.357,0.004) -- (2.353,0.036) -- (2.349,0.068) -- (2.343,0.100) -- (2.337,0.131) -- (2.331,0.162) -- (2.324,0.193) -- (2.316,0.224) -- (2.308,0.255) -- (2.299,0.285) -- (2.289,0.315) -- (2.279,0.345) -- (2.269,0.374) -- (2.257,0.403) -- (2.246,0.432) -- (2.233,0.461) -- (2.221,0.489) -- (2.207,0.516) -- (2.193,0.544) -- (2.179,0.571) -- (2.164,0.597) -- (2.149,0.623) -- (2.133,0.649) -- (2.117,0.674) -- (2.100,0.699) -- (2.083,0.724) -- (2.066,0.748) -- (2.048,0.771) -- (2.029,0.794) -- (2.011,0.817) -- (1.992,0.839) -- (1.972,0.860) -- (1.952,0.881) -- (1.932,0.902) -- (1.912,0.922) -- (1.891,0.942) -- (1.870,0.961) -- (1.849,0.979) -- (1.827,0.997) -- (1.805,1.014) -- (1.783,1.031) -- (1.761,1.048) -- (1.738,1.064) -- (1.716,1.079) -- (1.693,1.094) -- (1.670,1.108) -- (1.647,1.121) -- (1.623,1.135) -- (1.600,1.147) -- (1.576,1.159) -- (1.553,1.171) -- (1.529,1.182) -- (1.505,1.192) -- (1.481,1.202) -- (1.457,1.211) -- (1.433,1.220) -- (1.409,1.228) -- (1.385,1.236) -- (1.361,1.243) -- (1.337,1.250) -- (1.313,1.256) -- (1.289,1.262) -- (1.265,1.267) -- (1.241,1.272) -- (1.217,1.276) -- (1.193,1.280) -- (1.170,1.283) -- (1.146,1.286) -- (1.122,1.289) -- (1.099,1.291) -- (1.076,1.292) -- (1.053,1.293) -- (1.030,1.294) -- (1.007,1.294) -- (0.984,1.294) -- (0.962,1.293) -- (0.939,1.292) -- (0.917,1.290) -- (0.895,1.289) -- (0.873,1.287) -- (0.851,1.284) -- (0.830,1.281) -- (0.808,1.278) -- (0.787,1.274) -- (0.766,1.271) -- (0.746,1.266) -- (0.725,1.262) -- (0.705,1.257) -- (0.685,1.252) -- (0.665,1.247) -- (0.645,1.241) -- (0.626,1.236) -- (0.607,1.230) -- (0.588,1.224) -- (0.569,1.217) -- (0.551,1.210) -- (0.532,1.204) -- (0.514,1.197) -- (0.496,1.189) -- (0.479,1.182) -- (0.461,1.175) -- (0.444,1.167) -- (0.427,1.159) -- (0.410,1.151) -- (0.393,1.144) -- (0.377,1.135) -- (0.361,1.127) -- (0.345,1.119) -- (0.329,1.111) -- (0.313,1.102) -- (0.298,1.094) -- (0.282,1.086) -- (0.267,1.077) -- (0.252,1.069) -- (0.237,1.060) -- (0.222,1.052) -- (0.208,1.043) -- (0.193,1.035) -- (0.179,1.026) -- (0.165,1.018) -- (0.151,1.009) -- (0.136,1.001) -- (0.123,0.992) -- (0.109,0.984) -- (0.095,0.976) -- (0.081,0.968) -- (0.068,0.959) -- (0.054,0.951) -- (0.040,0.943) -- (0.027,0.935) -- (0.013,0.928) -- (-0.000,0.920) -- (-0.013,0.912) -- (-0.027,0.905) -- (-0.040,0.897) -- (-0.054,0.890) -- (-0.067,0.882) -- (-0.081,0.875) -- (-0.094,0.868) -- (-0.108,0.861) -- (-0.122,0.854) -- (-0.135,0.848) -- (-0.149,0.841) -- (-0.163,0.834) -- (-0.177,0.828) -- (-0.191,0.822) -- (-0.205,0.815) -- (-0.219,0.809) -- (-0.234,0.803) -- (-0.248,0.797) -- (-0.263,0.792) -- (-0.278,0.786) -- (-0.293,0.780) -- (-0.308,0.775) -- (-0.323,0.769) -- (-0.338,0.764) -- (-0.354,0.759) -- (-0.370,0.754) -- (-0.386,0.749) -- (-0.402,0.744);
\draw[fillback] (-0.402,0.744) -- (-0.418,0.739) -- (-0.434,0.734) -- (-0.451,0.729) -- (-0.468,0.724) -- (-0.485,0.720) -- (-0.502,0.715) -- (-0.520,0.710) -- (-0.537,0.706) -- (-0.555,0.701) -- (-0.573,0.697) -- (-0.592,0.692) -- (-0.610,0.688) -- (-0.629,0.683) -- (-0.648,0.679) -- (-0.667,0.674) -- (-0.686,0.670) -- (-0.706,0.665) -- (-0.726,0.661) -- (-0.746,0.656) -- (-0.766,0.652) -- (-0.786,0.647) -- (-0.807,0.642) -- (-0.828,0.637) -- (-0.849,0.632) -- (-0.870,0.628) -- (-0.891,0.622) -- (-0.913,0.617) -- (-0.935,0.612) -- (-0.957,0.607) -- (-0.979,0.601) -- (-1.001,0.596) -- (-1.023,0.590) -- (-1.046,0.584) -- (-1.069,0.578) -- (-1.091,0.572) -- (-1.114,0.566) -- (-1.137,0.560) -- (-1.161,0.553) -- (-1.184,0.547) -- (-1.207,0.540) -- (-1.231,0.533) -- (-1.254,0.525) -- (-1.278,0.518) -- (-1.301,0.510) -- (-1.325,0.503) -- (-1.349,0.494) -- (-1.372,0.486) -- (-1.396,0.478) -- (-1.420,0.469) -- (-1.443,0.460) -- (-1.467,0.451) -- (-1.490,0.442) -- (-1.514,0.432) -- (-1.537,0.423) -- (-1.561,0.413) -- (-1.584,0.402) -- (-1.607,0.392) -- (-1.630,0.381) -- (-1.653,0.370) -- (-1.676,0.359) -- (-1.699,0.347) -- (-1.721,0.336) -- (-1.743,0.324) -- (-1.766,0.311) -- (-1.787,0.299) -- (-1.809,0.286) -- (-1.830,0.273) -- (-1.852,0.260) -- (-1.872,0.247) -- (-1.893,0.233) -- (-1.913,0.219) -- (-1.933,0.205) -- (-1.953,0.191) -- (-1.972,0.176) -- (-1.991,0.161) -- (-2.010,0.146) -- (-2.028,0.131) -- (-2.046,0.115) -- (-2.064,0.100) -- (-2.081,0.084) -- (-2.098,0.068) -- (-2.114,0.051) -- (-2.130,0.035) -- (-2.145,0.018) -- (-2.160,0.001) -- (-2.174,-0.016) -- (-2.188,-0.033) -- (-2.202,-0.051) -- (-2.214,-0.068) -- (-2.227,-0.086) -- (-2.239,-0.104) -- (-2.250,-0.122) -- (-2.261,-0.140) -- (-2.271,-0.158) -- (-2.281,-0.176) -- (-2.290,-0.195) -- (-2.298,-0.213) -- (-2.306,-0.232) -- (-2.313,-0.251) -- (-2.320,-0.269) -- (-2.326,-0.288) -- (-2.331,-0.307) -- (-2.336,-0.326) -- (-2.341,-0.345) -- (-2.344,-0.364) -- (-2.347,-0.383) -- (-2.349,-0.402) -- (-2.351,-0.421) -- (-2.352,-0.440) -- (-2.352,-0.459) -- (-2.352,-0.478) -- (-2.351,-0.497) -- (-2.350,-0.515) -- (-2.347,-0.534) -- (-2.345,-0.553) -- (-2.341,-0.571) -- (-2.337,-0.590) -- (-2.332,-0.608) -- (-2.326,-0.626) -- (-2.320,-0.645);
\draw[fill] (-2.320,-0.645) -- (-2.313,-0.663) -- (-2.306,-0.680) -- (-2.298,-0.698) -- (-2.289,-0.716) -- (-2.280,-0.733) -- (-2.270,-0.750) -- (-2.259,-0.767) -- (-2.248,-0.784) -- (-2.236,-0.800) -- (-2.224,-0.817) -- (-2.211,-0.833) -- (-2.197,-0.849) -- (-2.183,-0.864) -- (-2.168,-0.879) -- (-2.152,-0.895) -- (-2.136,-0.909) -- (-2.120,-0.924) -- (-2.103,-0.938) -- (-2.085,-0.952) -- (-2.067,-0.966) -- (-2.048,-0.979) -- (-2.029,-0.992) -- (-2.009,-1.005) -- (-1.989,-1.017) -- (-1.969,-1.029) -- (-1.948,-1.041) -- (-1.926,-1.052) -- (-1.904,-1.063) -- (-1.882,-1.074) -- (-1.859,-1.084) -- (-1.836,-1.094) -- (-1.812,-1.103) -- (-1.788,-1.112) -- (-1.764,-1.121) -- (-1.739,-1.130) -- (-1.714,-1.138) -- (-1.689,-1.145) -- (-1.663,-1.153) -- (-1.637,-1.159) -- (-1.611,-1.166) -- (-1.584,-1.172) -- (-1.557,-1.178) -- (-1.531,-1.183) -- (-1.503,-1.188) -- (-1.476,-1.192) -- (-1.448,-1.197) -- (-1.421,-1.200) -- (-1.393,-1.204) -- (-1.365,-1.207) -- (-1.336,-1.209) -- (-1.308,-1.212) -- (-1.280,-1.213) -- (-1.251,-1.215) -- (-1.223,-1.216) -- (-1.194,-1.217) -- (-1.165,-1.217) -- (-1.137,-1.217) -- (-1.108,-1.217) -- (-1.079,-1.216) -- (-1.050,-1.215) -- (-1.022,-1.213) -- (-0.993,-1.212) -- (-0.965,-1.210) -- (-0.936,-1.207) -- (-0.908,-1.205) -- (-0.879,-1.202) -- (-0.851,-1.198) -- (-0.823,-1.195) -- (-0.795,-1.191) -- (-0.767,-1.186) -- (-0.739,-1.182) -- (-0.712,-1.177) -- (-0.684,-1.172) -- (-0.657,-1.167) -- (-0.630,-1.161) -- (-0.603,-1.156) -- (-0.577,-1.150) -- (-0.550,-1.143) -- (-0.524,-1.137) -- (-0.498,-1.130) -- (-0.472,-1.124) -- (-0.447,-1.117) -- (-0.422,-1.109) -- (-0.397,-1.102) -- (-0.372,-1.095) -- (-0.348,-1.087) -- (-0.324,-1.079) -- (-0.300,-1.071) -- (-0.277,-1.063) -- (-0.253,-1.055) -- (-0.231,-1.047) -- (-0.208,-1.039) -- (-0.186,-1.030) -- (-0.164,-1.022) -- (-0.142,-1.013) -- (-0.121,-1.005) -- (-0.100,-0.996) -- (-0.079,-0.988) -- (-0.059,-0.979) -- (-0.039,-0.970) -- (-0.019,-0.962) -- (-0.000,-0.953) -- (0.019,-0.945) -- (0.038,-0.936) -- (0.056,-0.927) -- (0.074,-0.919) -- (0.092,-0.910) -- (0.109,-0.902) -- (0.126,-0.894) -- (0.143,-0.885) -- (0.159,-0.877) -- (0.175,-0.869) -- (0.191,-0.861) -- (0.207,-0.853) -- (0.222,-0.845) -- (0.237,-0.837) -- (0.252,-0.830) -- (0.266,-0.822) -- (0.280,-0.815) -- (0.294,-0.808) -- (0.307,-0.801) -- (0.321,-0.794) -- (0.334,-0.787) -- (0.347,-0.781) -- (0.359,-0.774) -- (0.371,-0.768) -- (0.384,-0.762) -- (0.395,-0.756) -- (0.407,-0.750) -- (0.419,-0.745) -- (0.430,-0.739) -- (0.441,-0.734) -- (0.452,-0.729) -- (0.463,-0.724) -- (0.473,-0.720) -- (0.484,-0.715) -- (0.494,-0.711) -- (0.504,-0.707);
\draw[fillback] (0.504,-0.707) -- (0.514,-0.703) -- (0.524,-0.699) -- (0.534,-0.696) -- (0.544,-0.692) -- (0.554,-0.689) -- (0.563,-0.686) -- (0.573,-0.683) -- (0.582,-0.680) -- (0.592,-0.678) -- (0.601,-0.675) -- (0.610,-0.673) -- (0.620,-0.671) -- (0.629,-0.669) -- (0.638,-0.668) -- (0.648,-0.666) -- (0.657,-0.664) -- (0.666,-0.663) -- (0.676,-0.662) -- (0.685,-0.661) -- (0.694,-0.660) -- (0.704,-0.659) -- (0.713,-0.658) -- (0.723,-0.657) -- (0.733,-0.657) -- (0.742,-0.656) -- (0.752,-0.656) -- (0.762,-0.655) -- (0.772,-0.655) -- (0.782,-0.654) -- (0.792,-0.654) -- (0.802,-0.654) -- (0.813,-0.654) -- (0.823,-0.653) -- (0.834,-0.653) -- (0.845,-0.653) -- (0.855,-0.652) -- (0.866,-0.652) -- (0.878,-0.652) -- (0.889,-0.651) -- (0.900,-0.651) -- (0.912,-0.650) -- (0.923,-0.649) -- (0.935,-0.649) -- (0.947,-0.648) -- (0.959,-0.647) -- (0.972,-0.646) -- (0.984,-0.645) -- (0.996,-0.643) -- (1.009,-0.642) -- (1.022,-0.640) -- (1.035,-0.638) -- (1.048,-0.636) -- (1.061,-0.634) -- (1.074,-0.631) -- (1.088,-0.629) -- (1.101,-0.626) -- (1.115,-0.622) -- (1.129,-0.619) -- (1.143,-0.615) -- (1.157,-0.611) -- (1.171,-0.607) -- (1.185,-0.603) -- (1.199,-0.598) -- (1.214,-0.593) -- (1.228,-0.588) -- (1.243,-0.582) -- (1.257,-0.576) -- (1.272,-0.569) -- (1.286,-0.563) -- (1.301,-0.556) -- (1.316,-0.548) -- (1.330,-0.540) -- (1.345,-0.532) -- (1.360,-0.524) -- (1.375,-0.515) -- (1.389,-0.506) -- (1.404,-0.496) -- (1.419,-0.486) -- (1.433,-0.475) -- (1.448,-0.465) -- (1.462,-0.453) -- (1.476,-0.442) -- (1.491,-0.430) -- (1.505,-0.417) -- (1.519,-0.404) -- (1.533,-0.391) -- (1.546,-0.377) -- (1.560,-0.363) -- (1.573,-0.348) -- (1.587,-0.333) -- (1.600,-0.317) -- (1.613,-0.301) -- (1.625,-0.285) -- (1.638,-0.268) -- (1.650,-0.251) -- (1.662,-0.233) -- (1.674,-0.215) -- (1.685,-0.197) -- (1.696,-0.178) -- (1.707,-0.159) -- (1.718,-0.139) -- (1.728,-0.119) -- (1.738,-0.098) -- (1.747,-0.077) -- (1.757,-0.056) -- (1.765,-0.034) -- (1.774,-0.012) -- (1.782,0.010) -- (1.790,0.033) -- (1.797,0.056) -- (1.804,0.080) -- (1.810,0.104) -- (1.816,0.128) -- (1.822,0.152) -- (1.827,0.177) -- (1.831,0.202) -- (1.836,0.228) -- (1.839,0.254) -- (1.842,0.280) -- (1.845,0.306) -- (1.847,0.333) -- (1.849,0.360) -- (1.850,0.387) -- (1.851,0.414) -- (1.851,0.441) -- (1.850,0.469) -- (1.849,0.497) -- (1.847,0.525) -- (1.845,0.553) -- (1.843,0.582) -- (1.839,0.610) -- (1.836,0.639) -- (1.831,0.668) -- (1.826,0.696) -- (1.821,0.725) -- (1.814,0.754) -- (1.808,0.783) -- (1.800,0.813) -- (1.792,0.842) -- (1.784,0.871) -- (1.775,0.900) -- (1.765,0.929) -- (1.755,0.958) -- (1.744,0.987) -- (1.733,1.016) -- (1.721,1.045) -- (1.708,1.074) -- (1.695,1.103) -- (1.681,1.131) -- (1.667,1.160) -- (1.652,1.188) -- (1.636,1.216) -- (1.620,1.244) -- (1.604,1.272) -- (1.587,1.299) -- (1.569,1.327) -- (1.551,1.354) -- (1.532,1.381) -- (1.513,1.407) -- (1.493,1.434) -- (1.473,1.460) -- (1.452,1.485) -- (1.431,1.511) -- (1.409,1.536) -- (1.387,1.560) -- (1.365,1.585) -- (1.342,1.609) -- (1.318,1.632) -- (1.294,1.656) -- (1.270,1.678) -- (1.245,1.701) -- (1.220,1.723) -- (1.195,1.744) -- (1.169,1.765) -- (1.143,1.786) -- (1.116,1.806) -- (1.089,1.826) -- (1.062,1.845) -- (1.035,1.863) -- (1.007,1.882) -- (0.979,1.899) -- (0.951,1.916) -- (0.922,1.933) -- (0.893,1.949) -- (0.864,1.964) -- (0.835,1.979) -- (0.806,1.994) -- (0.776,2.007) -- (0.747,2.021) -- (0.717,2.033) -- (0.687,2.045) -- (0.657,2.057) -- (0.627,2.068) -- (0.597,2.078) -- (0.566,2.088) -- (0.536,2.097) -- (0.506,2.105) -- (0.475,2.113) -- (0.445,2.120) -- (0.415,2.127) -- (0.384,2.133) -- (0.354,2.138) -- (0.324,2.143) -- (0.294,2.147) -- (0.264,2.151) -- (0.234,2.154) -- (0.204,2.156);
\draw[fill] (0.204,2.156) -- (0.174,2.158) -- (0.145,2.159) -- (0.116,2.160) -- (0.086,2.160) -- (0.057,2.159) -- (0.029,2.158) -- (-0.000,2.156) -- (-0.028,2.154) -- (-0.056,2.151) -- (-0.084,2.148) -- (-0.112,2.144) -- (-0.139,2.139) -- (-0.166,2.134) -- (-0.193,2.128) -- (-0.220,2.122) -- (-0.246,2.115) -- (-0.272,2.108) -- (-0.297,2.100) -- (-0.322,2.091) -- (-0.347,2.083) -- (-0.371,2.073) -- (-0.395,2.063) -- (-0.419,2.053) -- (-0.442,2.042) -- (-0.465,2.031) -- (-0.488,2.019) -- (-0.510,2.007) -- (-0.531,1.994) -- (-0.553,1.981) -- (-0.573,1.968) -- (-0.594,1.954) -- (-0.614,1.940) -- (-0.633,1.925) -- (-0.652,1.910) -- (-0.671,1.895) -- (-0.689,1.879) -- (-0.707,1.863) -- (-0.724,1.846) -- (-0.741,1.830) -- (-0.757,1.813) -- (-0.773,1.795) -- (-0.789,1.778) -- (-0.804,1.760) -- (-0.818,1.741) -- (-0.832,1.723) -- (-0.846,1.704) -- (-0.859,1.685) -- (-0.872,1.666) -- (-0.884,1.647) -- (-0.896,1.627) -- (-0.908,1.607) -- (-0.919,1.587) -- (-0.929,1.567) -- (-0.939,1.547) -- (-0.949,1.526) -- (-0.958,1.506) -- (-0.967,1.485) -- (-0.976,1.464) -- (-0.984,1.443) -- (-0.992,1.422) -- (-0.999,1.401) -- (-1.006,1.380) -- (-1.012,1.358) -- (-1.019,1.337) -- (-1.024,1.315) -- (-1.030,1.294) -- (-1.035,1.272) -- (-1.040,1.251) -- (-1.044,1.229) -- (-1.049,1.207) -- (-1.052,1.186) -- (-1.056,1.164) -- (-1.059,1.142) -- (-1.062,1.121) -- (-1.065,1.099) -- (-1.068,1.077) -- (-1.070,1.056) -- (-1.072,1.034) -- (-1.074,1.013) -- (-1.075,0.991) -- (-1.076,0.969) -- (-1.077,0.948) -- (-1.078,0.927) -- (-1.079,0.905) -- (-1.080,0.884) -- (-1.080,0.863) -- (-1.080,0.842) -- (-1.080,0.821) -- (-1.080,0.800) -- (-1.080,0.779) -- (-1.080,0.758) -- (-1.080,0.737) -- (-1.079,0.716) -- (-1.079,0.696) -- (-1.078,0.675) -- (-1.078,0.655) -- (-1.077,0.634) -- (-1.076,0.614) -- (-1.075,0.594) -- (-1.074,0.574) -- (-1.074,0.554) -- (-1.073,0.534) -- (-1.072,0.514) -- (-1.071,0.494) -- (-1.070,0.474) -- (-1.069,0.454) -- (-1.069,0.435) -- (-1.068,0.415) -- (-1.067,0.396) -- (-1.067,0.377) -- (-1.066,0.357) -- (-1.066,0.338) -- (-1.065,0.319) -- (-1.065,0.300) -- (-1.064,0.281) -- (-1.064,0.262) -- (-1.064,0.243) -- (-1.064,0.224) -- (-1.064,0.205) -- (-1.064,0.186) -- (-1.064,0.167) -- (-1.065,0.148) -- (-1.065,0.130) -- (-1.066,0.111) -- (-1.066,0.092) -- (-1.067,0.073) -- (-1.068,0.055) -- (-1.069,0.036) -- (-1.070,0.017);
\draw[fillback] (-1.070,0.017) -- (-1.071,-0.001) -- (-1.073,-0.020) -- (-1.074,-0.039) -- (-1.076,-0.058) -- (-1.077,-0.076) -- (-1.079,-0.095) -- (-1.081,-0.114) -- (-1.083,-0.133) -- (-1.085,-0.152) -- (-1.087,-0.171) -- (-1.089,-0.190) -- (-1.092,-0.209) -- (-1.094,-0.228) -- (-1.097,-0.247) -- (-1.099,-0.266) -- (-1.102,-0.285) -- (-1.104,-0.304) -- (-1.107,-0.324) -- (-1.110,-0.343) -- (-1.113,-0.363) -- (-1.116,-0.382) -- (-1.119,-0.402) -- (-1.122,-0.422) -- (-1.125,-0.442) -- (-1.128,-0.461) -- (-1.131,-0.481) -- (-1.134,-0.501) -- (-1.137,-0.522) -- (-1.140,-0.542) -- (-1.143,-0.562) -- (-1.145,-0.583) -- (-1.148,-0.603) -- (-1.151,-0.624) -- (-1.154,-0.644) -- (-1.156,-0.665) -- (-1.159,-0.686) -- (-1.161,-0.707) -- (-1.164,-0.728) -- (-1.166,-0.749) -- (-1.168,-0.770) -- (-1.170,-0.791) -- (-1.172,-0.813) -- (-1.173,-0.834) -- (-1.175,-0.855) -- (-1.176,-0.877) -- (-1.177,-0.899) -- (-1.178,-0.920) -- (-1.179,-0.942) -- (-1.179,-0.964) -- (-1.179,-0.986) -- (-1.179,-1.007) -- (-1.179,-1.029) -- (-1.178,-1.051) -- (-1.178,-1.073) -- (-1.177,-1.095) -- (-1.175,-1.117) -- (-1.174,-1.139) -- (-1.172,-1.161) -- (-1.169,-1.183) -- (-1.167,-1.205) -- (-1.164,-1.227) -- (-1.160,-1.249) -- (-1.157,-1.271) -- (-1.153,-1.293) -- (-1.148,-1.315) -- (-1.143,-1.336) -- (-1.138,-1.358) -- (-1.133,-1.380) -- (-1.127,-1.401) -- (-1.120,-1.423) -- (-1.113,-1.444) -- (-1.106,-1.465) -- (-1.098,-1.486) -- (-1.090,-1.507) -- (-1.082,-1.528) -- (-1.073,-1.549) -- (-1.063,-1.569) -- (-1.054,-1.589) -- (-1.043,-1.610) -- (-1.032,-1.629) -- (-1.021,-1.649) -- (-1.009,-1.669) -- (-0.997,-1.688) -- (-0.985,-1.707) -- (-0.971,-1.726) -- (-0.958,-1.745) -- (-0.944,-1.763) -- (-0.929,-1.781) -- (-0.914,-1.799) -- (-0.899,-1.816) -- (-0.883,-1.833) -- (-0.866,-1.850) -- (-0.849,-1.867) -- (-0.832,-1.883) -- (-0.814,-1.898) -- (-0.795,-1.914) -- (-0.777,-1.929) -- (-0.757,-1.944) -- (-0.738,-1.958) -- (-0.717,-1.972) -- (-0.697,-1.985) -- (-0.676,-1.998) -- (-0.654,-2.011) -- (-0.632,-2.023) -- (-0.610,-2.035) -- (-0.587,-2.046) -- (-0.564,-2.057) -- (-0.541,-2.067) -- (-0.517,-2.077) -- (-0.492,-2.086) -- (-0.468,-2.095) -- (-0.442,-2.104) -- (-0.417,-2.111) -- (-0.391,-2.119) -- (-0.365,-2.126) -- (-0.339,-2.132) -- (-0.312,-2.137) -- (-0.285,-2.143) -- (-0.257,-2.147) -- (-0.230,-2.151) -- (-0.202,-2.155) -- (-0.174,-2.157) -- (-0.145,-2.160) -- (-0.117,-2.162) -- (-0.088,-2.163);
\draw[fill] (-0.088,-2.163) -- (-0.059,-2.163) -- (-0.029,-2.163) -- (-0.000,-2.163) -- (0.030,-2.161) -- (0.059,-2.160) -- (0.089,-2.157) -- (0.119,-2.154) -- (0.149,-2.150) -- (0.180,-2.146) -- (0.210,-2.141) -- (0.240,-2.136) -- (0.271,-2.130) -- (0.301,-2.123) -- (0.331,-2.116) -- (0.362,-2.108) -- (0.392,-2.100) -- (0.423,-2.091) -- (0.453,-2.081) -- (0.483,-2.071) -- (0.514,-2.060) -- (0.544,-2.049) -- (0.574,-2.037) -- (0.604,-2.025) -- (0.633,-2.012) -- (0.663,-1.998) -- (0.693,-1.984) -- (0.722,-1.969) -- (0.751,-1.954) -- (0.780,-1.938) -- (0.808,-1.922) -- (0.837,-1.905) -- (0.865,-1.888) -- (0.893,-1.870) -- (0.920,-1.851) -- (0.948,-1.833) -- (0.975,-1.813) -- (1.001,-1.794) -- (1.028,-1.773) -- (1.054,-1.753) -- (1.079,-1.732) -- (1.105,-1.710) -- (1.129,-1.688) -- (1.154,-1.666) -- (1.178,-1.643) -- (1.202,-1.620) -- (1.225,-1.597) -- (1.248,-1.573) -- (1.270,-1.549) -- (1.292,-1.525) -- (1.313,-1.500) -- (1.334,-1.475) -- (1.355,-1.449) -- (1.375,-1.424) -- (1.394,-1.398) -- (1.413,-1.372) -- (1.431,-1.345) -- (1.449,-1.319) -- (1.467,-1.292) -- (1.483,-1.265) -- (1.500,-1.238) -- (1.516,-1.211) -- (1.531,-1.183) -- (1.545,-1.156) -- (1.559,-1.128) -- (1.573,-1.100) -- (1.586,-1.072) -- (1.598,-1.045) -- (1.610,-1.017) -- (1.622,-0.988) -- (1.632,-0.960) -- (1.642,-0.932) -- (1.652,-0.904) -- (1.661,-0.876) -- (1.669,-0.848) -- (1.677,-0.820) -- (1.685,-0.792) -- (1.691,-0.764) -- (1.698,-0.736) -- (1.703,-0.709) -- (1.708,-0.681) -- (1.713,-0.654) -- (1.717,-0.626) -- (1.720,-0.599) -- (1.723,-0.572) -- (1.725,-0.545) -- (1.727,-0.518) -- (1.728,-0.492) -- (1.729,-0.466) -- (1.729,-0.440) -- (1.729,-0.414) -- (1.728,-0.388) -- (1.727,-0.363) -- (1.725,-0.338) -- (1.723,-0.313) -- (1.720,-0.289) -- (1.717,-0.264) -- (1.714,-0.240) -- (1.710,-0.217) -- (1.705,-0.193) -- (1.700,-0.171) -- (1.695,-0.148) -- (1.689,-0.126) -- (1.683,-0.104) -- (1.677,-0.082) -- (1.670,-0.061) -- (1.663,-0.040) -- (1.655,-0.019) -- (1.647,0.001) -- (1.639,0.020) -- (1.631,0.040) -- (1.622,0.059) -- (1.613,0.077) -- (1.603,0.096) -- (1.594,0.113) -- (1.584,0.131) -- (1.574,0.148) -- (1.563,0.164) -- (1.553,0.180) -- (1.542,0.196) -- (1.531,0.212) -- (1.519,0.226) -- (1.508,0.241) -- (1.496,0.255) -- (1.485,0.269) -- (1.473,0.282) -- (1.461,0.295) -- (1.449,0.308) -- (1.437,0.320) -- (1.424,0.332) -- (1.412,0.343) -- (1.399,0.354) -- (1.387,0.364) -- (1.374,0.375) -- (1.362,0.384) -- (1.349,0.394) -- (1.336,0.403) -- (1.323,0.412) -- (1.311,0.420) -- (1.298,0.428) -- (1.285,0.436) -- (1.273,0.443) -- (1.260,0.450) -- (1.247,0.457) -- (1.235,0.463) -- (1.222,0.469) -- (1.210,0.475) -- (1.197,0.481) -- (1.185,0.486) -- (1.173,0.491) -- (1.161,0.496) -- (1.148,0.500) -- (1.137,0.504) -- (1.125,0.508) -- (1.113,0.512) -- (1.101,0.516) -- (1.090,0.519) -- (1.078,0.522) -- (1.067,0.525) -- (1.056,0.528) -- (1.045,0.531) -- (1.034,0.533) -- (1.023,0.535) -- (1.012,0.538) -- (1.002,0.540) -- (0.991,0.542) -- (0.981,0.543) -- (0.971,0.545) -- (0.961,0.547) -- (0.951,0.549) -- (0.941,0.550) -- (0.931,0.552) -- (0.922,0.553) -- (0.912,0.554) -- (0.903,0.556) -- (0.894,0.557) -- (0.885,0.559) -- (0.876,0.560) -- (0.867,0.561) -- (0.858,0.563) -- (0.850,0.564) -- (0.841,0.566) -- (0.833,0.567) -- (0.825,0.569) -- (0.816,0.571) -- (0.808,0.572) -- (0.800,0.574) -- (0.792,0.576) -- (0.784,0.578) -- (0.776,0.581) -- (0.768,0.583) -- (0.760,0.585) -- (0.753,0.588) -- (0.745,0.590) -- (0.737,0.593) -- (0.729,0.596) -- (0.721,0.599) -- (0.713,0.603) -- (0.706,0.606) -- (0.698,0.610) -- (0.690,0.613) -- (0.682,0.617) -- (0.674,0.622) -- (0.666,0.626) -- (0.658,0.630) -- (0.649,0.635) -- (0.641,0.640);
\draw[fillback] (0.641,0.640) -- (0.633,0.645) -- (0.624,0.650) -- (0.616,0.656) -- (0.607,0.662) -- (0.598,0.668) -- (0.589,0.674) -- (0.580,0.680) -- (0.570,0.686) -- (0.561,0.693) -- (0.551,0.700) -- (0.541,0.707) -- (0.531,0.714) -- (0.521,0.722) -- (0.510,0.730) -- (0.499,0.737) -- (0.488,0.745) -- (0.477,0.754) -- (0.466,0.762) -- (0.454,0.771) -- (0.442,0.779) -- (0.430,0.788) -- (0.417,0.797) -- (0.404,0.807) -- (0.391,0.816) -- (0.378,0.825) -- (0.364,0.835) -- (0.350,0.845) -- (0.336,0.855) -- (0.321,0.865) -- (0.306,0.875) -- (0.291,0.885) -- (0.275,0.896) -- (0.259,0.906) -- (0.243,0.917) -- (0.226,0.927) -- (0.209,0.938) -- (0.192,0.949) -- (0.175,0.959) -- (0.157,0.970) -- (0.138,0.981) -- (0.119,0.992) -- (0.100,1.003) -- (0.081,1.013) -- (0.061,1.024) -- (0.041,1.035) -- (0.021,1.046) -- (0.000,1.057) -- (-0.021,1.068) -- (-0.043,1.078) -- (-0.064,1.089) -- (-0.087,1.099) -- (-0.109,1.110) -- (-0.132,1.120) -- (-0.155,1.130) -- (-0.178,1.140) -- (-0.202,1.150) -- (-0.226,1.160) -- (-0.251,1.170) -- (-0.275,1.180) -- (-0.300,1.189) -- (-0.326,1.198) -- (-0.351,1.207) -- (-0.377,1.216) -- (-0.403,1.225) -- (-0.429,1.233) -- (-0.456,1.241) -- (-0.483,1.249) -- (-0.510,1.257) -- (-0.537,1.264) -- (-0.565,1.271) -- (-0.592,1.278) -- (-0.620,1.285) -- (-0.648,1.291) -- (-0.676,1.297) -- (-0.705,1.303) -- (-0.733,1.309) -- (-0.762,1.314) -- (-0.791,1.318) -- (-0.819,1.323) -- (-0.848,1.327) -- (-0.877,1.331) -- (-0.906,1.334) -- (-0.935,1.337) -- (-0.965,1.340) -- (-0.994,1.342) -- (-1.023,1.344) -- (-1.052,1.346) -- (-1.081,1.347) -- (-1.110,1.348) -- (-1.140,1.348) -- (-1.169,1.348) -- (-1.198,1.348) -- (-1.226,1.347) -- (-1.255,1.346) -- (-1.284,1.344) -- (-1.312,1.342) -- (-1.341,1.340) -- (-1.369,1.337) -- (-1.397,1.334) -- (-1.425,1.330) -- (-1.452,1.326) -- (-1.480,1.321) -- (-1.507,1.316) -- (-1.534,1.311) -- (-1.560,1.305) -- (-1.587,1.299) -- (-1.613,1.293) -- (-1.638,1.286) -- (-1.664,1.278) -- (-1.689,1.270) -- (-1.713,1.262) -- (-1.738,1.254) -- (-1.762,1.245) -- (-1.785,1.235) -- (-1.809,1.225) -- (-1.831,1.215) -- (-1.854,1.205) -- (-1.875,1.194) -- (-1.897,1.183) -- (-1.918,1.171) -- (-1.938,1.159) -- (-1.958,1.147) -- (-1.978,1.134) -- (-1.997,1.121) -- (-2.015,1.108) -- (-2.033,1.094) -- (-2.050,1.080) -- (-2.067,1.066) -- (-2.084,1.051) -- (-2.099,1.036) -- (-2.114,1.021) -- (-2.129,1.006) -- (-2.143,0.990) -- (-2.156,0.974) -- (-2.169,0.958) -- (-2.181,0.941) -- (-2.193,0.925) -- (-2.204,0.908) -- (-2.214,0.891) -- (-2.224,0.873) -- (-2.233,0.856) -- (-2.241,0.838);
\draw[fill] (-2.241,0.838) -- (-2.249,0.820) -- (-2.256,0.802) -- (-2.263,0.784) -- (-2.269,0.766) -- (-2.274,0.747) -- (-2.278,0.729) -- (-2.282,0.710) -- (-2.285,0.691) -- (-2.288,0.672) -- (-2.290,0.653) -- (-2.291,0.634) -- (-2.292,0.615) -- (-2.292,0.596) -- (-2.291,0.577) -- (-2.290,0.557) -- (-2.288,0.538) -- (-2.286,0.519) -- (-2.283,0.499) -- (-2.279,0.480) -- (-2.275,0.461) -- (-2.270,0.442) -- (-2.264,0.422) -- (-2.258,0.403) -- (-2.251,0.384) -- (-2.244,0.365) -- (-2.236,0.346) -- (-2.228,0.327) -- (-2.219,0.308) -- (-2.209,0.289) -- (-2.199,0.271) -- (-2.188,0.252) -- (-2.177,0.234) -- (-2.165,0.215) -- (-2.153,0.197) -- (-2.140,0.179) -- (-2.127,0.161) -- (-2.114,0.144) -- (-2.100,0.126) -- (-2.085,0.109) -- (-2.070,0.091) -- (-2.055,0.074) -- (-2.039,0.057) -- (-2.023,0.041) -- (-2.006,0.024) -- (-1.989,0.008) -- (-1.972,-0.008) -- (-1.954,-0.024) -- (-1.936,-0.040) -- (-1.917,-0.055) -- (-1.899,-0.071) -- (-1.880,-0.086) -- (-1.860,-0.101) -- (-1.841,-0.115) -- (-1.821,-0.130) -- (-1.801,-0.144) -- (-1.781,-0.158) -- (-1.760,-0.172) -- (-1.739,-0.185) -- (-1.718,-0.199) -- (-1.697,-0.212) -- (-1.676,-0.225) -- (-1.654,-0.237) -- (-1.633,-0.250) -- (-1.611,-0.262) -- (-1.589,-0.274) -- (-1.567,-0.286) -- (-1.545,-0.297) -- (-1.523,-0.309) -- (-1.501,-0.320) -- (-1.479,-0.331) -- (-1.456,-0.341) -- (-1.434,-0.352) -- (-1.412,-0.362) -- (-1.389,-0.372) -- (-1.367,-0.382) -- (-1.345,-0.392) -- (-1.322,-0.401) -- (-1.300,-0.411) -- (-1.278,-0.420) -- (-1.255,-0.429) -- (-1.233,-0.438) -- (-1.211,-0.446) -- (-1.189,-0.455) -- (-1.167,-0.463) -- (-1.146,-0.471) -- (-1.124,-0.480) -- (-1.102,-0.488) -- (-1.081,-0.495) -- (-1.060,-0.503) -- (-1.038,-0.511) -- (-1.017,-0.518) -- (-0.997,-0.525) -- (-0.976,-0.533) -- (-0.955,-0.540) -- (-0.935,-0.547) -- (-0.915,-0.554) -- (-0.895,-0.561) -- (-0.875,-0.568) -- (-0.855,-0.575) -- (-0.835,-0.581) -- (-0.816,-0.588) -- (-0.797,-0.595) -- (-0.778,-0.601) -- (-0.759,-0.608) -- (-0.741,-0.615) -- (-0.722,-0.621) -- (-0.704,-0.628) -- (-0.686,-0.634) -- (-0.669,-0.641) -- (-0.651,-0.648) -- (-0.634,-0.654) -- (-0.616,-0.661) -- (-0.599,-0.667) -- (-0.583,-0.674) -- (-0.566,-0.681) -- (-0.550,-0.688) -- (-0.533,-0.694) -- (-0.517,-0.701);
\draw[fillback] (-0.517,-0.701) -- (-0.502,-0.708) -- (-0.486,-0.715) -- (-0.470,-0.722) -- (-0.455,-0.729) -- (-0.440,-0.736) -- (-0.425,-0.744) -- (-0.410,-0.751) -- (-0.395,-0.759) -- (-0.381,-0.766) -- (-0.366,-0.774) -- (-0.352,-0.781) -- (-0.338,-0.789) -- (-0.324,-0.797) -- (-0.310,-0.805) -- (-0.296,-0.813) -- (-0.282,-0.821) -- (-0.268,-0.829) -- (-0.255,-0.838) -- (-0.241,-0.846) -- (-0.228,-0.855) -- (-0.214,-0.864) -- (-0.201,-0.872) -- (-0.187,-0.881) -- (-0.174,-0.890) -- (-0.161,-0.899) -- (-0.147,-0.908) -- (-0.134,-0.918) -- (-0.121,-0.927) -- (-0.108,-0.936) -- (-0.094,-0.946) -- (-0.081,-0.955) -- (-0.068,-0.965) -- (-0.054,-0.975) -- (-0.041,-0.985) -- (-0.027,-0.995) -- (-0.014,-1.004) -- (-0.000,-1.014) -- (0.014,-1.024) -- (0.028,-1.034) -- (0.042,-1.045) -- (0.056,-1.055) -- (0.070,-1.065) -- (0.084,-1.075) -- (0.099,-1.085) -- (0.113,-1.095) -- (0.128,-1.106) -- (0.143,-1.116) -- (0.158,-1.126) -- (0.173,-1.136) -- (0.188,-1.146) -- (0.204,-1.156) -- (0.220,-1.166) -- (0.236,-1.176) -- (0.252,-1.186) -- (0.268,-1.196) -- (0.285,-1.205) -- (0.301,-1.215) -- (0.318,-1.224) -- (0.335,-1.234) -- (0.353,-1.243) -- (0.370,-1.252) -- (0.388,-1.261) -- (0.406,-1.270) -- (0.425,-1.278) -- (0.443,-1.287) -- (0.462,-1.295) -- (0.481,-1.303) -- (0.500,-1.311) -- (0.520,-1.318) -- (0.539,-1.326) -- (0.559,-1.333) -- (0.579,-1.340) -- (0.600,-1.346) -- (0.620,-1.353) -- (0.641,-1.359) -- (0.662,-1.364) -- (0.684,-1.370) -- (0.705,-1.375) -- (0.727,-1.380) -- (0.749,-1.384) -- (0.771,-1.388) -- (0.794,-1.392) -- (0.816,-1.395) -- (0.839,-1.398) -- (0.862,-1.401) -- (0.885,-1.403) -- (0.909,-1.405) -- (0.932,-1.406) -- (0.956,-1.407) -- (0.980,-1.407) -- (1.004,-1.407) -- (1.028,-1.407) -- (1.053,-1.406) -- (1.077,-1.405) -- (1.102,-1.403) -- (1.126,-1.401) -- (1.151,-1.398) -- (1.176,-1.395) -- (1.201,-1.391) -- (1.226,-1.387) -- (1.251,-1.382) -- (1.276,-1.377) -- (1.301,-1.371) -- (1.326,-1.364) -- (1.352,-1.357) -- (1.377,-1.350) -- (1.402,-1.342) -- (1.427,-1.333) -- (1.452,-1.324) -- (1.477,-1.315) -- (1.502,-1.305) -- (1.527,-1.294) -- (1.552,-1.283) -- (1.577,-1.271) -- (1.601,-1.258) -- (1.626,-1.245) -- (1.650,-1.232) -- (1.675,-1.218) -- (1.699,-1.203) -- (1.722,-1.188) -- (1.746,-1.172) -- (1.770,-1.156) -- (1.793,-1.139) -- (1.816,-1.122) -- (1.839,-1.104) -- (1.861,-1.085) -- (1.883,-1.066) -- (1.905,-1.047) -- (1.927,-1.027) -- (1.948,-1.006) -- (1.969,-0.985) -- (1.990,-0.963) -- (2.010,-0.941) -- (2.030,-0.919) -- (2.049,-0.895) -- (2.068,-0.872) -- (2.087,-0.848) -- (2.105,-0.823) -- (2.123,-0.798) -- (2.141,-0.773) -- (2.157,-0.747) -- (2.174,-0.721) -- (2.190,-0.694) -- (2.205,-0.667) -- (2.220,-0.639) -- (2.235,-0.612) -- (2.249,-0.583) -- (2.262,-0.555) -- (2.275,-0.526) -- (2.287,-0.496) -- (2.299,-0.467) -- (2.310,-0.437) -- (2.321,-0.407) -- (2.331,-0.376) -- (2.340,-0.345) -- (2.349,-0.314) -- (2.357,-0.283) -- (2.365,-0.252) -- (2.372,-0.220) -- (2.378,-0.188) -- (2.384,-0.156) -- (2.389,-0.124) -- (2.393,-0.091) -- (2.397,-0.059) -- (2.400,-0.026) -- (2.403,0.007) -- (2.404,0.040) -- (2.406,0.073) -- (2.406,0.106) -- (2.406,0.139) -- (2.405,0.172) -- (2.404,0.205) -- (2.401,0.238) -- (2.399,0.271) -- (2.395,0.304) -- (2.391,0.337) -- (2.386,0.370) -- (2.381,0.403) -- (2.375,0.435) -- (2.368,0.468) -- (2.361,0.500) -- (2.353,0.533) -- (2.344,0.565) -- (2.335,0.597) -- (2.325,0.628) -- (2.314,0.660) -- (2.303,0.691) -- (2.291,0.722) -- (2.279,0.753) -- (2.266,0.784);
\draw[fill] (2.266,0.784) -- (2.252,0.814) -- (2.238,0.844) -- (2.223,0.874) -- (2.208,0.903) -- (2.192,0.932) -- (2.176,0.961) -- (2.159,0.989) -- (2.142,1.017) -- (2.124,1.044) -- (2.105,1.072) -- (2.086,1.098) -- (2.067,1.125) -- (2.047,1.150) -- (2.026,1.176) -- (2.006,1.201) -- (1.984,1.225) -- (1.963,1.249) -- (1.940,1.273) -- (1.918,1.296) -- (1.895,1.318) -- (1.872,1.340) -- (1.848,1.362) -- (1.824,1.382) -- (1.800,1.403) -- (1.775,1.423) -- (1.750,1.442) -- (1.725,1.461) -- (1.700,1.479) -- (1.674,1.497) -- (1.648,1.514) -- (1.622,1.530) -- (1.595,1.546) -- (1.569,1.561) -- (1.542,1.576) -- (1.515,1.590) -- (1.488,1.604) -- (1.461,1.617) -- (1.433,1.629) -- (1.406,1.641) -- (1.378,1.652) -- (1.350,1.662) -- (1.322,1.672) -- (1.295,1.682) -- (1.267,1.690) -- (1.239,1.699) -- (1.211,1.706) -- (1.183,1.713) -- (1.155,1.719) -- (1.127,1.725) -- (1.099,1.730) -- (1.071,1.735) -- (1.043,1.739) -- (1.016,1.743) -- (0.988,1.746) -- (0.961,1.748) -- (0.933,1.750) -- (0.906,1.751) -- (0.879,1.752) -- (0.852,1.752) -- (0.825,1.751) -- (0.798,1.751) -- (0.771,1.749) -- (0.745,1.747) -- (0.719,1.745) -- (0.693,1.742) -- (0.667,1.738) -- (0.642,1.734) -- (0.616,1.730) -- (0.591,1.725) -- (0.566,1.720) -- (0.542,1.714) -- (0.517,1.708) -- (0.493,1.701) -- (0.469,1.694) -- (0.446,1.687) -- (0.422,1.679) -- (0.399,1.671) -- (0.377,1.662) -- (0.354,1.653) -- (0.332,1.644) -- (0.310,1.634) -- (0.288,1.624) -- (0.267,1.614) -- (0.246,1.604) -- (0.225,1.593) -- (0.205,1.581) -- (0.185,1.570) -- (0.165,1.558) -- (0.145,1.546) -- (0.126,1.534) -- (0.107,1.521) -- (0.089,1.509) -- (0.070,1.496) -- (0.052,1.482) -- (0.035,1.469) -- (0.017,1.456) -- (0.000,1.442) -- (-0.017,1.428) -- (-0.033,1.414) -- (-0.050,1.400) -- (-0.066,1.385) -- (-0.082,1.371) -- (-0.097,1.356) -- (-0.112,1.342) -- (-0.127,1.327) -- (-0.142,1.312) -- (-0.157,1.297) -- (-0.171,1.282) -- (-0.185,1.267) -- (-0.199,1.252) -- (-0.213,1.237) -- (-0.226,1.222) -- (-0.239,1.207) -- (-0.252,1.191) -- (-0.265,1.176) -- (-0.278,1.161) -- (-0.290,1.146) -- (-0.303,1.130) -- (-0.315,1.115) -- (-0.327,1.100) -- (-0.339,1.085) -- (-0.351,1.070) -- (-0.363,1.055) -- (-0.374,1.040) -- (-0.386,1.025) -- (-0.397,1.010) -- (-0.408,0.995) -- (-0.419,0.980) -- (-0.431,0.965) -- (-0.442,0.950) -- (-0.453,0.936) -- (-0.464,0.921) -- (-0.475,0.906) -- (-0.486,0.892) -- (-0.497,0.878) -- (-0.508,0.863) -- (-0.519,0.849) -- (-0.529,0.835) -- (-0.540,0.821) -- (-0.551,0.807) -- (-0.562,0.793) -- (-0.574,0.779) -- (-0.585,0.765) -- (-0.596,0.752) -- (-0.607,0.738) -- (-0.618,0.725) -- (-0.630,0.711) -- (-0.641,0.698) -- (-0.653,0.685) -- (-0.665,0.671) -- (-0.676,0.658) -- (-0.688,0.645) -- (-0.700,0.632) -- (-0.712,0.619) -- (-0.724,0.606) -- (-0.737,0.593) -- (-0.749,0.581) -- (-0.762,0.568) -- (-0.775,0.555) -- (-0.788,0.543) -- (-0.801,0.530) -- (-0.814,0.517) -- (-0.827,0.505) -- (-0.841,0.492);
\draw[fillback] (-0.841,0.492) -- (-0.854,0.480) -- (-0.868,0.467) -- (-0.882,0.455) -- (-0.896,0.442) -- (-0.910,0.430) -- (-0.925,0.417) -- (-0.939,0.404) -- (-0.954,0.392) -- (-0.969,0.379) -- (-0.984,0.367) -- (-0.999,0.354) -- (-1.014,0.341) -- (-1.030,0.329) -- (-1.045,0.316) -- (-1.061,0.303) -- (-1.077,0.290) -- (-1.093,0.277) -- (-1.109,0.264) -- (-1.125,0.251) -- (-1.142,0.238) -- (-1.158,0.225) -- (-1.175,0.211) -- (-1.191,0.198) -- (-1.208,0.184) -- (-1.225,0.170) -- (-1.242,0.157) -- (-1.259,0.143) -- (-1.276,0.129) -- (-1.293,0.115) -- (-1.310,0.100) -- (-1.327,0.086) -- (-1.344,0.071) -- (-1.361,0.057) -- (-1.379,0.042) -- (-1.396,0.027) -- (-1.413,0.012) -- (-1.430,-0.003) -- (-1.447,-0.019) -- (-1.464,-0.034) -- (-1.481,-0.050) -- (-1.498,-0.066) -- (-1.515,-0.082) -- (-1.532,-0.098) -- (-1.548,-0.114) -- (-1.565,-0.131) -- (-1.581,-0.148) -- (-1.598,-0.164) -- (-1.614,-0.181) -- (-1.630,-0.199) -- (-1.645,-0.216) -- (-1.661,-0.233) -- (-1.676,-0.251) -- (-1.692,-0.269) -- (-1.707,-0.287) -- (-1.721,-0.305) -- (-1.736,-0.323) -- (-1.750,-0.342) -- (-1.764,-0.360) -- (-1.778,-0.379) -- (-1.791,-0.398) -- (-1.804,-0.417) -- (-1.817,-0.436) -- (-1.829,-0.455) -- (-1.841,-0.475) -- (-1.853,-0.494) -- (-1.865,-0.514) -- (-1.875,-0.534) -- (-1.886,-0.553) -- (-1.896,-0.573) -- (-1.906,-0.593) -- (-1.915,-0.614) -- (-1.924,-0.634) -- (-1.933,-0.654) -- (-1.941,-0.674) -- (-1.948,-0.695) -- (-1.955,-0.715) -- (-1.962,-0.736) -- (-1.968,-0.756) -- (-1.974,-0.777) -- (-1.979,-0.798) -- (-1.983,-0.818) -- (-1.987,-0.839) -- (-1.991,-0.860) -- (-1.994,-0.880) -- (-1.996,-0.901) -- (-1.998,-0.921) -- (-1.999,-0.942) -- (-2.000,-0.963) -- (-2.000,-0.983) -- (-2.000,-1.003) -- (-1.999,-1.024) -- (-1.997,-1.044) -- (-1.995,-1.064) -- (-1.992,-1.084) -- (-1.989,-1.104) -- (-1.985,-1.124) -- (-1.980,-1.143) -- (-1.975,-1.163) -- (-1.969,-1.182) -- (-1.963,-1.201) -- (-1.956,-1.220) -- (-1.948,-1.239) -- (-1.940,-1.258) -- (-1.931,-1.276) -- (-1.921,-1.294) -- (-1.911,-1.312) -- (-1.901,-1.330) -- (-1.889,-1.347) -- (-1.878,-1.365) -- (-1.865,-1.382) -- (-1.852,-1.398) -- (-1.838,-1.415) -- (-1.824,-1.431) -- (-1.809,-1.446) -- (-1.794,-1.462) -- (-1.778,-1.477);
\draw[fill] (-1.778,-1.477) -- (-1.762,-1.492) -- (-1.745,-1.506) -- (-1.727,-1.520) -- (-1.709,-1.534) -- (-1.690,-1.547) -- (-1.671,-1.560) -- (-1.651,-1.573) -- (-1.631,-1.585) -- (-1.610,-1.597) -- (-1.589,-1.608) -- (-1.568,-1.619) -- (-1.545,-1.630) -- (-1.523,-1.640) -- (-1.500,-1.650) -- (-1.476,-1.659) -- (-1.452,-1.668) -- (-1.428,-1.676) -- (-1.403,-1.684) -- (-1.378,-1.691) -- (-1.353,-1.698) -- (-1.327,-1.705) -- (-1.301,-1.711) -- (-1.275,-1.717) -- (-1.248,-1.722) -- (-1.221,-1.726) -- (-1.193,-1.730) -- (-1.165,-1.734) -- (-1.138,-1.737) -- (-1.109,-1.740) -- (-1.081,-1.742) -- (-1.052,-1.744) -- (-1.023,-1.745) -- (-0.994,-1.745) -- (-0.965,-1.746) -- (-0.936,-1.745) -- (-0.906,-1.745) -- (-0.877,-1.743) -- (-0.847,-1.742) -- (-0.817,-1.739) -- (-0.787,-1.737) -- (-0.757,-1.733) -- (-0.727,-1.730) -- (-0.697,-1.726) -- (-0.667,-1.721) -- (-0.636,-1.716) -- (-0.606,-1.710) -- (-0.576,-1.704) -- (-0.546,-1.698) -- (-0.516,-1.691) -- (-0.486,-1.684) -- (-0.456,-1.676) -- (-0.426,-1.668) -- (-0.397,-1.659) -- (-0.367,-1.650) -- (-0.338,-1.640) -- (-0.309,-1.631) -- (-0.279,-1.620) -- (-0.250,-1.610) -- (-0.222,-1.599) -- (-0.193,-1.587) -- (-0.165,-1.576) -- (-0.137,-1.564) -- (-0.109,-1.551) -- (-0.081,-1.538) -- (-0.054,-1.525) -- (-0.027,-1.512) -- (0.000,-1.498) -- (0.027,-1.484) -- (0.053,-1.470) -- (0.079,-1.456) -- (0.104,-1.441) -- (0.130,-1.426) -- (0.154,-1.410) -- (0.179,-1.395) -- (0.203,-1.379) -- (0.227,-1.363) -- (0.251,-1.347) -- (0.274,-1.331) -- (0.297,-1.315) -- (0.319,-1.298) -- (0.341,-1.281) -- (0.362,-1.264) -- (0.384,-1.247) -- (0.404,-1.230) -- (0.425,-1.213) -- (0.445,-1.196) -- (0.464,-1.178) -- (0.483,-1.161) -- (0.502,-1.144) -- (0.520,-1.126) -- (0.538,-1.108) -- (0.556,-1.091) -- (0.573,-1.073) -- (0.589,-1.056) -- (0.605,-1.038) -- (0.621,-1.021) -- (0.637,-1.003) -- (0.651,-0.986) -- (0.666,-0.968) -- (0.680,-0.951) -- (0.694,-0.934) -- (0.707,-0.917) -- (0.720,-0.900) -- (0.732,-0.883) -- (0.744,-0.866) -- (0.756,-0.849) -- (0.767,-0.833) -- (0.778,-0.816) -- (0.788,-0.800) -- (0.799,-0.784) -- (0.808,-0.768) -- (0.818,-0.752) -- (0.827,-0.737) -- (0.835,-0.722) -- (0.844,-0.706) -- (0.852,-0.691) -- (0.859,-0.677) -- (0.867,-0.662) -- (0.874,-0.648) -- (0.880,-0.634) -- (0.887,-0.620) -- (0.893,-0.606) -- (0.899,-0.593) -- (0.904,-0.580) -- (0.910,-0.567) -- (0.915,-0.554) -- (0.920,-0.542) -- (0.924,-0.530) -- (0.929,-0.518) -- (0.933,-0.506) -- (0.937,-0.495) -- (0.941,-0.484) -- (0.944,-0.473) -- (0.948,-0.462) -- (0.951,-0.452) -- (0.954,-0.442) -- (0.957,-0.432) -- (0.960,-0.422) -- (0.962,-0.413) -- (0.965,-0.404) -- (0.967,-0.395) -- (0.970,-0.387) -- (0.972,-0.378) -- (0.974,-0.370) -- (0.976,-0.362) -- (0.978,-0.355) -- (0.980,-0.347) -- (0.982,-0.340) -- (0.984,-0.333) -- (0.986,-0.327) -- (0.988,-0.320) -- (0.990,-0.314) -- (0.992,-0.308) -- (0.994,-0.302) -- (0.996,-0.296) -- (0.997,-0.291) -- (0.999,-0.285) -- (1.001,-0.280) -- (1.003,-0.275) -- (1.006,-0.270) -- (1.008,-0.266) -- (1.010,-0.261);
\draw[fillback] (1.010,-0.261) -- (1.012,-0.256) -- (1.015,-0.252) -- (1.017,-0.248) -- (1.019,-0.244) -- (1.022,-0.240) -- (1.025,-0.236) -- (1.028,-0.232) -- (1.030,-0.228) -- (1.033,-0.224) -- (1.036,-0.220) -- (1.040,-0.216) -- (1.043,-0.213) -- (1.046,-0.209) -- (1.050,-0.205) -- (1.054,-0.202) -- (1.057,-0.198) -- (1.061,-0.194) -- (1.065,-0.190) -- (1.069,-0.186) -- (1.073,-0.182) -- (1.078,-0.178) -- (1.082,-0.174) -- (1.087,-0.170) -- (1.091,-0.166) -- (1.096,-0.161) -- (1.101,-0.157) -- (1.106,-0.152) -- (1.111,-0.148) -- (1.116,-0.143) -- (1.121,-0.137) -- (1.127,-0.132) -- (1.132,-0.127) -- (1.138,-0.121) -- (1.143,-0.115) -- (1.149,-0.109) -- (1.154,-0.103) -- (1.160,-0.097) -- (1.166,-0.090) -- (1.172,-0.083) -- (1.177,-0.076) -- (1.183,-0.068) -- (1.189,-0.061) -- (1.195,-0.053) -- (1.201,-0.044) -- (1.207,-0.036) -- (1.212,-0.027) -- (1.218,-0.018) -- (1.224,-0.009) -- (1.230,0.001) -- (1.235,0.011) -- (1.241,0.021) -- (1.246,0.032) -- (1.252,0.043) -- (1.257,0.054) -- (1.262,0.066) -- (1.267,0.077) -- (1.272,0.090) -- (1.277,0.102) -- (1.282,0.115) -- (1.286,0.128) -- (1.291,0.142) -- (1.295,0.156) -- (1.299,0.170) -- (1.303,0.185) -- (1.306,0.200) -- (1.309,0.215) -- (1.313,0.230) -- (1.315,0.246) -- (1.318,0.262) -- (1.320,0.279) -- (1.322,0.296) -- (1.324,0.313) -- (1.326,0.331) -- (1.327,0.349) -- (1.328,0.367) -- (1.328,0.385) -- (1.328,0.404) -- (1.328,0.423) -- (1.328,0.442) -- (1.327,0.462) -- (1.325,0.482) -- (1.324,0.502) -- (1.322,0.522) -- (1.319,0.543) -- (1.316,0.564) -- (1.313,0.585) -- (1.310,0.606) -- (1.305,0.628) -- (1.301,0.650) -- (1.296,0.672) -- (1.290,0.694) -- (1.284,0.717) -- (1.278,0.739) -- (1.271,0.762) -- (1.264,0.785) -- (1.256,0.808) -- (1.248,0.831) -- (1.239,0.854) -- (1.230,0.878) -- (1.220,0.901) -- (1.210,0.925) -- (1.199,0.948) -- (1.188,0.972) -- (1.176,0.996) -- (1.164,1.020) -- (1.151,1.044) -- (1.138,1.067) -- (1.124,1.091) -- (1.110,1.115) -- (1.095,1.139) -- (1.080,1.163) -- (1.064,1.186) -- (1.048,1.210) -- (1.031,1.233) -- (1.014,1.257) -- (0.996,1.280) -- (0.978,1.303) -- (0.959,1.326) -- (0.940,1.349) -- (0.920,1.372) -- (0.900,1.395) -- (0.880,1.417) -- (0.859,1.439) -- (0.837,1.461) -- (0.815,1.483) -- (0.793,1.505) -- (0.770,1.526) -- (0.747,1.547) -- (0.723,1.568) -- (0.699,1.588) -- (0.675,1.608) -- (0.650,1.628) -- (0.625,1.648) -- (0.599,1.667) -- (0.573,1.686) -- (0.547,1.704) -- (0.520,1.722) -- (0.494,1.740) -- (0.466,1.757) -- (0.439,1.774) -- (0.411,1.791) -- (0.383,1.807) -- (0.355,1.822) -- (0.326,1.838) -- (0.297,1.852) -- (0.268,1.867) -- (0.239,1.881) -- (0.210,1.894) -- (0.180,1.907) -- (0.151,1.919) -- (0.121,1.931) -- (0.091,1.943) -- (0.061,1.953) -- (0.030,1.964) -- (0.000,1.974) -- (-0.030,1.983) -- (-0.061,1.992) -- (-0.091,2.000) -- (-0.122,2.008) -- (-0.152,2.015) -- (-0.183,2.022) -- (-0.213,2.028) -- (-0.244,2.033) -- (-0.274,2.038) -- (-0.305,2.043) -- (-0.335,2.046) -- (-0.365,2.050) -- (-0.395,2.053) -- (-0.425,2.055) -- (-0.455,2.056) -- (-0.484,2.057) -- (-0.514,2.058) -- (-0.543,2.058) -- (-0.572,2.057) -- (-0.601,2.056) -- (-0.630,2.054) -- (-0.658,2.052) -- (-0.686,2.049) -- (-0.714,2.046) -- (-0.742,2.042) -- (-0.769,2.037) -- (-0.796,2.032) -- (-0.822,2.027) -- (-0.849,2.021) -- (-0.875,2.014) -- (-0.900,2.007) -- (-0.925,2.000) -- (-0.950,1.991) -- (-0.975,1.983) -- (-0.999,1.974) -- (-1.022,1.964);
\draw[fill] (-1.022,1.964) -- (-1.046,1.954) -- (-1.068,1.944) -- (-1.091,1.933) -- (-1.113,1.921) -- (-1.134,1.909) -- (-1.155,1.897) -- (-1.175,1.884) -- (-1.195,1.871) -- (-1.215,1.857) -- (-1.234,1.843) -- (-1.252,1.829) -- (-1.270,1.814) -- (-1.288,1.799) -- (-1.305,1.783) -- (-1.321,1.768) -- (-1.337,1.751) -- (-1.352,1.735) -- (-1.367,1.718) -- (-1.382,1.701) -- (-1.395,1.683) -- (-1.409,1.665) -- (-1.421,1.647) -- (-1.433,1.629) -- (-1.445,1.610) -- (-1.456,1.591) -- (-1.466,1.572) -- (-1.476,1.553) -- (-1.486,1.533) -- (-1.495,1.514) -- (-1.503,1.494) -- (-1.511,1.473) -- (-1.518,1.453) -- (-1.525,1.433) -- (-1.531,1.412) -- (-1.537,1.391) -- (-1.542,1.370) -- (-1.546,1.349) -- (-1.551,1.328) -- (-1.554,1.306) -- (-1.557,1.285) -- (-1.560,1.263) -- (-1.562,1.242) -- (-1.564,1.220) -- (-1.565,1.198) -- (-1.566,1.177) -- (-1.566,1.155) -- (-1.566,1.133) -- (-1.565,1.111) -- (-1.564,1.089) -- (-1.562,1.067) -- (-1.560,1.045) -- (-1.558,1.023) -- (-1.555,1.001) -- (-1.552,0.979) -- (-1.549,0.957) -- (-1.545,0.936) -- (-1.541,0.914) -- (-1.536,0.892) -- (-1.531,0.870) -- (-1.526,0.849) -- (-1.520,0.827) -- (-1.514,0.805) -- (-1.508,0.784) -- (-1.502,0.763) -- (-1.495,0.741) -- (-1.488,0.720) -- (-1.481,0.699) -- (-1.473,0.678) -- (-1.466,0.657) -- (-1.458,0.636) -- (-1.450,0.615) -- (-1.441,0.595) -- (-1.433,0.574) -- (-1.424,0.554) -- (-1.415,0.533) -- (-1.406,0.513) -- (-1.397,0.493) -- (-1.388,0.473) -- (-1.378,0.453) -- (-1.369,0.433) -- (-1.359,0.414) -- (-1.350,0.394) -- (-1.340,0.375) -- (-1.330,0.355) -- (-1.320,0.336) -- (-1.310,0.317) -- (-1.300,0.298) -- (-1.290,0.279) -- (-1.280,0.261) -- (-1.270,0.242) -- (-1.260,0.223) -- (-1.250,0.205) -- (-1.240,0.187) -- (-1.230,0.168) -- (-1.221,0.150) -- (-1.211,0.132) -- (-1.201,0.114) -- (-1.191,0.096) -- (-1.181,0.078) -- (-1.172,0.061) -- (-1.162,0.043) -- (-1.153,0.025) -- (-1.143,0.008) -- (-1.134,-0.010) -- (-1.125,-0.027) -- (-1.116,-0.044) -- (-1.107,-0.062) -- (-1.098,-0.079) -- (-1.089,-0.096) -- (-1.080,-0.113) -- (-1.072,-0.130) -- (-1.063,-0.148) -- (-1.055,-0.165) -- (-1.047,-0.182) -- (-1.038,-0.199) -- (-1.030,-0.216) -- (-1.023,-0.233) -- (-1.015,-0.250);
\draw[fillback] (-1.015,-0.250) -- (-1.007,-0.267) -- (-1.000,-0.284) -- (-0.992,-0.301) -- (-0.985,-0.318) -- (-0.978,-0.336) -- (-0.971,-0.353) -- (-0.964,-0.370) -- (-0.957,-0.387) -- (-0.951,-0.404) -- (-0.944,-0.422) -- (-0.938,-0.439) -- (-0.931,-0.456) -- (-0.925,-0.474) -- (-0.919,-0.491) -- (-0.913,-0.509) -- (-0.907,-0.526) -- (-0.901,-0.544) -- (-0.895,-0.562) -- (-0.889,-0.579) -- (-0.883,-0.597) -- (-0.878,-0.615) -- (-0.872,-0.633) -- (-0.866,-0.651) -- (-0.861,-0.669) -- (-0.855,-0.687) -- (-0.849,-0.706) -- (-0.844,-0.724) -- (-0.838,-0.743) -- (-0.833,-0.761) -- (-0.827,-0.780) -- (-0.821,-0.798) -- (-0.816,-0.817) -- (-0.810,-0.836) -- (-0.804,-0.855) -- (-0.798,-0.874) -- (-0.792,-0.893) -- (-0.786,-0.912) -- (-0.780,-0.931) -- (-0.774,-0.950) -- (-0.768,-0.969) -- (-0.761,-0.988) -- (-0.755,-1.008) -- (-0.748,-1.027) -- (-0.741,-1.047) -- (-0.734,-1.066) -- (-0.727,-1.085) -- (-0.720,-1.105) -- (-0.712,-1.125) -- (-0.705,-1.144) -- (-0.697,-1.164) -- (-0.689,-1.183) -- (-0.680,-1.203) -- (-0.672,-1.222) -- (-0.663,-1.242) -- (-0.654,-1.261) -- (-0.645,-1.281) -- (-0.635,-1.300) -- (-0.625,-1.320) -- (-0.615,-1.339) -- (-0.605,-1.358) -- (-0.594,-1.378) -- (-0.583,-1.397) -- (-0.572,-1.416) -- (-0.560,-1.435) -- (-0.548,-1.454) -- (-0.536,-1.473) -- (-0.523,-1.491) -- (-0.510,-1.510) -- (-0.497,-1.528) -- (-0.483,-1.546) -- (-0.470,-1.564) -- (-0.455,-1.582) -- (-0.441,-1.600) -- (-0.425,-1.618) -- (-0.410,-1.635) -- (-0.394,-1.652) -- (-0.378,-1.669) -- (-0.362,-1.685) -- (-0.345,-1.702) -- (-0.327,-1.718) -- (-0.310,-1.734) -- (-0.292,-1.750) -- (-0.273,-1.765) -- (-0.255,-1.780) -- (-0.235,-1.795) -- (-0.216,-1.809) -- (-0.196,-1.823) -- (-0.176,-1.837) -- (-0.155,-1.850) -- (-0.134,-1.863) -- (-0.112,-1.876) -- (-0.091,-1.888) -- (-0.069,-1.900) -- (-0.046,-1.911) -- (-0.023,-1.922) -- (0.000,-1.933) -- (0.024,-1.943) -- (0.047,-1.953) -- (0.072,-1.962) -- (0.096,-1.971) -- (0.121,-1.979) -- (0.146,-1.987) -- (0.172,-1.994) -- (0.197,-2.001) -- (0.223,-2.008) -- (0.250,-2.013) -- (0.276,-2.019) -- (0.303,-2.024) -- (0.330,-2.028) -- (0.358,-2.032) -- (0.385,-2.035) -- (0.413,-2.037) -- (0.441,-2.039) -- (0.469,-2.041) -- (0.498,-2.042) -- (0.526,-2.042) -- (0.555,-2.042) -- (0.584,-2.041) -- (0.613,-2.039) -- (0.642,-2.037) -- (0.671,-2.035) -- (0.701,-2.031) -- (0.730,-2.027) -- (0.760,-2.023) -- (0.789,-2.018) -- (0.819,-2.012) -- (0.849,-2.006) -- (0.878,-1.999) -- (0.908,-1.991) -- (0.938,-1.983) -- (0.967,-1.974) -- (0.997,-1.965) -- (1.026,-1.955) -- (1.056,-1.944) -- (1.085,-1.933) -- (1.114,-1.921) -- (1.143,-1.909) -- (1.172,-1.895) -- (1.201,-1.882) -- (1.230,-1.867) -- (1.259,-1.853) -- (1.287,-1.837);
\draw[fill] (1.287,-1.837) -- (1.315,-1.821) -- (1.343,-1.804) -- (1.371,-1.787) -- (1.398,-1.769) -- (1.425,-1.751) -- (1.452,-1.732) -- (1.479,-1.713) -- (1.505,-1.693) -- (1.531,-1.672) -- (1.556,-1.651) -- (1.582,-1.630) -- (1.607,-1.608) -- (1.631,-1.585) -- (1.655,-1.562) -- (1.679,-1.539) -- (1.702,-1.515) -- (1.725,-1.491) -- (1.748,-1.466) -- (1.770,-1.440) -- (1.791,-1.415) -- (1.812,-1.389) -- (1.833,-1.362) -- (1.853,-1.335) -- (1.872,-1.308) -- (1.891,-1.281) -- (1.910,-1.253) -- (1.928,-1.225) -- (1.945,-1.196) -- (1.962,-1.167) -- (1.979,-1.138) -- (1.994,-1.109) -- (2.009,-1.079) -- (2.024,-1.049) -- (2.038,-1.019) -- (2.052,-0.989) -- (2.064,-0.958) -- (2.077,-0.927) -- (2.088,-0.896) -- (2.099,-0.865) -- (2.110,-0.834) -- (2.119,-0.803) -- (2.128,-0.772) -- (2.137,-0.740) -- (2.145,-0.709) -- (2.152,-0.677) -- (2.159,-0.645) -- (2.165,-0.614) -- (2.170,-0.582) -- (2.175,-0.550) -- (2.179,-0.519) -- (2.182,-0.487) -- (2.185,-0.455) -- (2.187,-0.424) -- (2.189,-0.392) -- (2.190,-0.361) -- (2.190,-0.330) -- (2.190,-0.299) -- (2.189,-0.268) -- (2.187,-0.237) -- (2.185,-0.206) -- (2.182,-0.176) -- (2.179,-0.146) -- (2.175,-0.116) -- (2.171,-0.086) -- (2.166,-0.056) -- (2.160,-0.027) -- (2.154,0.002) -- (2.147,0.031) -- (2.140,0.060) -- (2.132,0.088) -- (2.123,0.116) -- (2.114,0.143) -- (2.105,0.171) -- (2.095,0.197) -- (2.085,0.224) -- (2.074,0.250) -- (2.062,0.276) -- (2.051,0.301) -- (2.038,0.326) -- (2.025,0.351) -- (2.012,0.375) -- (1.999,0.399) -- (1.985,0.422) -- (1.970,0.445) -- (1.956,0.468) -- (1.940,0.490) -- (1.925,0.511) -- (1.909,0.533) -- (1.893,0.553) -- (1.876,0.573) -- (1.860,0.593) -- (1.842,0.612) -- (1.825,0.631) -- (1.807,0.650) -- (1.789,0.667) -- (1.771,0.685) -- (1.753,0.702) -- (1.734,0.718) -- (1.715,0.734) -- (1.696,0.749) -- (1.677,0.764) -- (1.658,0.778) -- (1.638,0.792) -- (1.619,0.805) -- (1.599,0.818) -- (1.579,0.831) -- (1.559,0.843) -- (1.539,0.854) -- (1.519,0.865) -- (1.498,0.875) -- (1.478,0.885) -- (1.458,0.895) -- (1.437,0.904) -- (1.417,0.912) -- (1.397,0.920) -- (1.376,0.928) -- (1.356,0.935) -- (1.335,0.942) -- (1.315,0.948) -- (1.295,0.954) -- (1.274,0.960) -- (1.254,0.965) -- (1.234,0.969) -- (1.214,0.973) -- (1.194,0.977) -- (1.174,0.981) -- (1.154,0.984) -- (1.134,0.986) -- (1.115,0.988) -- (1.095,0.990) -- (1.076,0.992) -- (1.057,0.993) -- (1.038,0.994) -- (1.019,0.995) -- (1.000,0.995) -- (0.981,0.995) -- (0.963,0.995) -- (0.944,0.994) -- (0.926,0.993) -- (0.908,0.992) -- (0.890,0.991) -- (0.873,0.989) -- (0.855,0.987) -- (0.838,0.985) -- (0.821,0.983) -- (0.804,0.981) -- (0.787,0.978) -- (0.771,0.975) -- (0.754,0.972) -- (0.738,0.969) -- (0.722,0.966) -- (0.706,0.962) -- (0.690,0.959) -- (0.675,0.955) -- (0.659,0.951) -- (0.644,0.948) -- (0.629,0.944) -- (0.614,0.940) -- (0.600,0.936) -- (0.585,0.931) -- (0.571,0.927) -- (0.556,0.923) -- (0.542,0.919) -- (0.528,0.915) -- (0.514,0.910) -- (0.501,0.906) -- (0.487,0.902) -- (0.474,0.898) -- (0.460,0.893) -- (0.447,0.889) -- (0.434,0.885) -- (0.421,0.881) -- (0.408,0.877) -- (0.395,0.873) -- (0.382,0.869) -- (0.369,0.865) -- (0.356,0.861) -- (0.343,0.858) -- (0.331,0.854) -- (0.318,0.851) -- (0.305,0.847) -- (0.293,0.844) -- (0.280,0.841) -- (0.267,0.838) -- (0.255,0.835) -- (0.242,0.832) -- (0.229,0.830) -- (0.217,0.827) -- (0.204,0.825) -- (0.191,0.822) -- (0.178,0.820) -- (0.165,0.818) -- (0.152,0.816) -- (0.139,0.814) -- (0.125,0.813) -- (0.112,0.811) -- (0.099,0.810) -- (0.085,0.809) -- (0.071,0.808) -- (0.057,0.807) -- (0.043,0.806) -- (0.029,0.805) -- (0.015,0.805) -- (0.000,0.804) -- (-0.015,0.804) -- (-0.030,0.804);
\draw[fillback] (-0.030,0.804) -- (-0.045,0.804) -- (-0.060,0.804) -- (-0.076,0.804) -- (-0.091,0.804) -- (-0.107,0.805) -- (-0.123,0.805) -- (-0.140,0.806) -- (-0.156,0.806) -- (-0.173,0.807) -- (-0.190,0.808) -- (-0.207,0.809) -- (-0.225,0.810) -- (-0.243,0.811) -- (-0.261,0.812) -- (-0.279,0.813) -- (-0.298,0.814) -- (-0.316,0.816) -- (-0.335,0.817) -- (-0.355,0.818) -- (-0.374,0.820) -- (-0.394,0.821) -- (-0.414,0.822) -- (-0.435,0.823) -- (-0.455,0.825) -- (-0.476,0.826) -- (-0.497,0.827) -- (-0.519,0.828) -- (-0.540,0.830) -- (-0.562,0.831) -- (-0.584,0.832) -- (-0.607,0.833) -- (-0.630,0.834) -- (-0.652,0.834) -- (-0.675,0.835) -- (-0.699,0.836) -- (-0.722,0.836) -- (-0.746,0.837) -- (-0.770,0.837) -- (-0.794,0.837) -- (-0.819,0.837) -- (-0.843,0.837) -- (-0.868,0.836) -- (-0.893,0.836) -- (-0.918,0.835) -- (-0.944,0.834) -- (-0.969,0.833) -- (-0.995,0.832) -- (-1.020,0.831) -- (-1.046,0.829) -- (-1.072,0.827) -- (-1.098,0.825) -- (-1.124,0.823) -- (-1.150,0.820) -- (-1.177,0.817) -- (-1.203,0.814) -- (-1.229,0.811) -- (-1.256,0.807) -- (-1.282,0.804) -- (-1.309,0.800) -- (-1.335,0.795) -- (-1.362,0.791) -- (-1.388,0.786) -- (-1.415,0.781) -- (-1.441,0.775) -- (-1.467,0.769) -- (-1.493,0.763) -- (-1.519,0.757) -- (-1.545,0.750) -- (-1.571,0.743) -- (-1.597,0.736) -- (-1.623,0.729) -- (-1.648,0.721) -- (-1.673,0.713) -- (-1.698,0.704) -- (-1.723,0.695) -- (-1.748,0.686) -- (-1.772,0.677) -- (-1.797,0.667) -- (-1.820,0.657) -- (-1.844,0.647) -- (-1.867,0.636) -- (-1.890,0.625) -- (-1.913,0.614) -- (-1.936,0.603) -- (-1.958,0.591) -- (-1.979,0.579) -- (-2.001,0.566) -- (-2.021,0.554) -- (-2.042,0.541) -- (-2.062,0.527) -- (-2.082,0.514) -- (-2.101,0.500) -- (-2.120,0.486) -- (-2.138,0.472) -- (-2.156,0.457) -- (-2.174,0.442) -- (-2.191,0.427) -- (-2.207,0.412) -- (-2.223,0.397) -- (-2.238,0.381) -- (-2.253,0.365) -- (-2.267,0.349) -- (-2.281,0.332) -- (-2.294,0.316) -- (-2.307,0.299) -- (-2.319,0.282) -- (-2.330,0.265) -- (-2.341,0.248) -- (-2.352,0.230) -- (-2.361,0.213) -- (-2.370,0.195) -- (-2.379,0.177) -- (-2.386,0.159) -- (-2.393,0.141) -- (-2.400,0.123) -- (-2.406,0.105) -- (-2.411,0.087) -- (-2.416,0.068) -- (-2.419,0.050) -- (-2.423,0.032) -- (-2.425,0.013) -- (-2.427,-0.006) -- (-2.428,-0.024) -- (-2.429,-0.043) -- (-2.429,-0.061) -- (-2.428,-0.080);
\draw[fill] (-2.428,-0.080) -- (-2.427,-0.098) -- (-2.425,-0.117) -- (-2.422,-0.135) -- (-2.419,-0.154) -- (-2.415,-0.172) -- (-2.410,-0.190) -- (-2.405,-0.209) -- (-2.399,-0.227) -- (-2.392,-0.245) -- (-2.385,-0.263) -- (-2.377,-0.280) -- (-2.368,-0.298) -- (-2.359,-0.316) -- (-2.349,-0.333) -- (-2.339,-0.350) -- (-2.328,-0.367) -- (-2.316,-0.384) -- (-2.304,-0.401) -- (-2.291,-0.417) -- (-2.278,-0.434) -- (-2.264,-0.450) -- (-2.249,-0.466) -- (-2.234,-0.481) -- (-2.218,-0.497) -- (-2.202,-0.512) -- (-2.186,-0.527) -- (-2.168,-0.542) -- (-2.151,-0.556) -- (-2.133,-0.571) -- (-2.114,-0.584) -- (-2.095,-0.598) -- (-2.075,-0.612) -- (-2.055,-0.625) -- (-2.035,-0.638) -- (-2.014,-0.650) -- (-1.993,-0.662) -- (-1.971,-0.674) -- (-1.949,-0.686) -- (-1.927,-0.697) -- (-1.904,-0.708) -- (-1.881,-0.719) -- (-1.857,-0.729) -- (-1.834,-0.739) -- (-1.810,-0.749) -- (-1.785,-0.759) -- (-1.761,-0.768) -- (-1.736,-0.776) -- (-1.711,-0.785) -- (-1.686,-0.793) -- (-1.660,-0.801) -- (-1.634,-0.808) -- (-1.608,-0.816) -- (-1.582,-0.822) -- (-1.556,-0.829) -- (-1.530,-0.835) -- (-1.503,-0.841) -- (-1.477,-0.847) -- (-1.450,-0.852) -- (-1.423,-0.857) -- (-1.396,-0.862) -- (-1.370,-0.866) -- (-1.343,-0.870) -- (-1.316,-0.874) -- (-1.289,-0.877) -- (-1.262,-0.880) -- (-1.235,-0.883) -- (-1.208,-0.886) -- (-1.181,-0.888) -- (-1.154,-0.890) -- (-1.127,-0.892) -- (-1.100,-0.894) -- (-1.074,-0.895) -- (-1.047,-0.896) -- (-1.020,-0.897) -- (-0.994,-0.898) -- (-0.968,-0.898) -- (-0.942,-0.898) -- (-0.916,-0.898) -- (-0.890,-0.898) -- (-0.864,-0.897) -- (-0.839,-0.897) -- (-0.813,-0.896) -- (-0.788,-0.895) -- (-0.763,-0.894) -- (-0.739,-0.892) -- (-0.714,-0.891) -- (-0.690,-0.889) -- (-0.666,-0.888) -- (-0.642,-0.886) -- (-0.618,-0.884) -- (-0.595,-0.882) -- (-0.572,-0.880) -- (-0.549,-0.877) -- (-0.526,-0.875) -- (-0.504,-0.873) -- (-0.482,-0.870) -- (-0.460,-0.868) -- (-0.438,-0.865) -- (-0.417,-0.862) -- (-0.396,-0.860) -- (-0.375,-0.857) -- (-0.355,-0.854) -- (-0.334,-0.852) -- (-0.314,-0.849) -- (-0.295,-0.846) -- (-0.275,-0.843) -- (-0.256,-0.841) -- (-0.237,-0.838) -- (-0.219,-0.836) -- (-0.200,-0.833) -- (-0.182,-0.831) -- (-0.164,-0.828) -- (-0.147,-0.826) -- (-0.130,-0.823) -- (-0.113,-0.821) -- (-0.096,-0.819) -- (-0.079,-0.817) -- (-0.063,-0.815) -- (-0.047,-0.813) -- (-0.031,-0.811) -- (-0.015,-0.810) -- (-0.000,-0.808) -- (0.015,-0.807) -- (0.030,-0.805) -- (0.045,-0.804) -- (0.059,-0.803) -- (0.074,-0.802) -- (0.088,-0.801);
\draw[fillback] (0.088,-0.801) -- (0.102,-0.801) -- (0.116,-0.800) -- (0.130,-0.800) -- (0.143,-0.800) -- (0.156,-0.799) -- (0.170,-0.799) -- (0.183,-0.800) -- (0.196,-0.800) -- (0.209,-0.800) -- (0.222,-0.801) -- (0.234,-0.802) -- (0.247,-0.803) -- (0.259,-0.804) -- (0.272,-0.805) -- (0.284,-0.806) -- (0.297,-0.808) -- (0.309,-0.809) -- (0.321,-0.811) -- (0.333,-0.813) -- (0.345,-0.815) -- (0.358,-0.817) -- (0.370,-0.819) -- (0.382,-0.822) -- (0.394,-0.824) -- (0.407,-0.827) -- (0.419,-0.829) -- (0.431,-0.832) -- (0.443,-0.835) -- (0.456,-0.838) -- (0.468,-0.841) -- (0.481,-0.844) -- (0.493,-0.847) -- (0.506,-0.850) -- (0.519,-0.853) -- (0.532,-0.856) -- (0.545,-0.860) -- (0.558,-0.863) -- (0.571,-0.866) -- (0.584,-0.870) -- (0.598,-0.873) -- (0.612,-0.876) -- (0.625,-0.879) -- (0.639,-0.883) -- (0.653,-0.886) -- (0.667,-0.889) -- (0.682,-0.892) -- (0.696,-0.895) -- (0.711,-0.898) -- (0.725,-0.901) -- (0.740,-0.904) -- (0.756,-0.906) -- (0.771,-0.909) -- (0.786,-0.911) -- (0.802,-0.913) -- (0.818,-0.916) -- (0.834,-0.917) -- (0.850,-0.919) -- (0.866,-0.921) -- (0.883,-0.922) -- (0.899,-0.923) -- (0.916,-0.924) -- (0.933,-0.925) -- (0.950,-0.926) -- (0.967,-0.926) -- (0.985,-0.926) -- (1.002,-0.926) -- (1.020,-0.925) -- (1.038,-0.924) -- (1.056,-0.923) -- (1.074,-0.922) -- (1.093,-0.920) -- (1.111,-0.918) -- (1.130,-0.916) -- (1.148,-0.913) -- (1.167,-0.910) -- (1.186,-0.907) -- (1.205,-0.903) -- (1.224,-0.899) -- (1.243,-0.894) -- (1.262,-0.889) -- (1.281,-0.884) -- (1.301,-0.878) -- (1.320,-0.872) -- (1.339,-0.865) -- (1.359,-0.858) -- (1.378,-0.851) -- (1.398,-0.843) -- (1.417,-0.835) -- (1.437,-0.826) -- (1.456,-0.817) -- (1.475,-0.807) -- (1.495,-0.797) -- (1.514,-0.786) -- (1.533,-0.775) -- (1.552,-0.764) -- (1.571,-0.752) -- (1.590,-0.739) -- (1.609,-0.726) -- (1.627,-0.713) -- (1.646,-0.699) -- (1.664,-0.684) -- (1.683,-0.670) -- (1.700,-0.654) -- (1.718,-0.638) -- (1.736,-0.622) -- (1.753,-0.605) -- (1.770,-0.588) -- (1.787,-0.570) -- (1.804,-0.552) -- (1.820,-0.533) -- (1.836,-0.514) -- (1.852,-0.494) -- (1.868,-0.474) -- (1.883,-0.453) -- (1.898,-0.432) -- (1.912,-0.411) -- (1.926,-0.389) -- (1.940,-0.367) -- (1.954,-0.344) -- (1.967,-0.321) -- (1.979,-0.297) -- (1.991,-0.273) -- (2.003,-0.248) -- (2.014,-0.224) -- (2.025,-0.198) -- (2.036,-0.173) -- (2.046,-0.147) -- (2.055,-0.121) -- (2.064,-0.094) -- (2.072,-0.067) -- (2.080,-0.040) -- (2.088,-0.012) -- (2.095,0.016) -- (2.101,0.044) -- (2.107,0.073) -- (2.112,0.102) -- (2.117,0.131) -- (2.121,0.160) -- (2.125,0.189) -- (2.128,0.219) -- (2.130,0.249) -- (2.132,0.279) -- (2.133,0.310) -- (2.134,0.340) -- (2.134,0.371) -- (2.133,0.401) -- (2.132,0.432) -- (2.130,0.463) -- (2.128,0.494) -- (2.125,0.526) -- (2.121,0.557) -- (2.117,0.588) -- (2.112,0.619) -- (2.106,0.651) -- (2.100,0.682) -- (2.093,0.713) -- (2.086,0.744) -- (2.078,0.776) -- (2.069,0.807) -- (2.060,0.838) -- (2.050,0.869) -- (2.039,0.900) -- (2.028,0.930) -- (2.017,0.961) -- (2.004,0.992) -- (1.991,1.022) -- (1.978,1.052) -- (1.964,1.082) -- (1.949,1.112) -- (1.933,1.141) -- (1.918,1.170) -- (1.901,1.199) -- (1.884,1.228) -- (1.867,1.256) -- (1.848,1.284) -- (1.830,1.312) -- (1.811,1.340) -- (1.791,1.367) -- (1.771,1.393) -- (1.750,1.420) -- (1.729,1.446) -- (1.707,1.471) -- (1.685,1.497) -- (1.663,1.521) -- (1.640,1.546) -- (1.616,1.570) -- (1.593,1.593) -- (1.568,1.616) -- (1.544,1.639) -- (1.519,1.661) -- (1.493,1.682) -- (1.468,1.703) -- (1.441,1.724) -- (1.415,1.744) -- (1.388,1.763) -- (1.361,1.782) -- (1.334,1.801) -- (1.307,1.818) -- (1.279,1.836) -- (1.251,1.852) -- (1.222,1.869) -- (1.194,1.884) -- (1.165,1.899) -- (1.136,1.913) -- (1.107,1.927) -- (1.078,1.940) -- (1.049,1.953);
\draw[fill] (1.049,1.953) -- (1.019,1.965) -- (0.990,1.976) -- (0.960,1.987) -- (0.931,1.997) -- (0.901,2.007) -- (0.871,2.016) -- (0.841,2.024) -- (0.811,2.032) -- (0.782,2.039) -- (0.752,2.045) -- (0.722,2.051) -- (0.692,2.056) -- (0.662,2.061) -- (0.633,2.065) -- (0.603,2.068) -- (0.574,2.071) -- (0.544,2.073) -- (0.515,2.075) -- (0.486,2.076) -- (0.457,2.076) -- (0.429,2.076) -- (0.400,2.075) -- (0.372,2.074) -- (0.343,2.072) -- (0.315,2.069) -- (0.288,2.066) -- (0.260,2.063) -- (0.233,2.058) -- (0.206,2.054) -- (0.179,2.048) -- (0.153,2.043) -- (0.126,2.036) -- (0.101,2.029) -- (0.075,2.022) -- (0.050,2.014) -- (0.025,2.006) -- (0.000,1.997) -- (-0.024,1.988) -- (-0.048,1.978) -- (-0.072,1.968) -- (-0.095,1.957) -- (-0.118,1.946) -- (-0.141,1.934) -- (-0.163,1.922) -- (-0.185,1.910) -- (-0.206,1.897) -- (-0.227,1.884) -- (-0.248,1.871) -- (-0.268,1.857) -- (-0.288,1.843) -- (-0.308,1.828) -- (-0.327,1.813) -- (-0.346,1.798) -- (-0.364,1.782) -- (-0.382,1.766) -- (-0.399,1.750) -- (-0.417,1.734) -- (-0.433,1.717) -- (-0.450,1.700) -- (-0.466,1.683) -- (-0.481,1.666) -- (-0.497,1.648) -- (-0.511,1.630) -- (-0.526,1.612) -- (-0.540,1.594) -- (-0.554,1.576) -- (-0.567,1.557) -- (-0.580,1.538) -- (-0.592,1.519) -- (-0.605,1.500) -- (-0.617,1.481) -- (-0.628,1.462) -- (-0.639,1.442) -- (-0.650,1.423) -- (-0.661,1.403) -- (-0.671,1.384) -- (-0.681,1.364) -- (-0.691,1.344) -- (-0.700,1.324) -- (-0.709,1.304) -- (-0.718,1.284) -- (-0.726,1.264) -- (-0.735,1.244) -- (-0.742,1.224) -- (-0.750,1.204) -- (-0.758,1.184) -- (-0.765,1.164) -- (-0.772,1.143) -- (-0.779,1.123) -- (-0.785,1.103) -- (-0.792,1.083) -- (-0.798,1.063) -- (-0.804,1.043) -- (-0.810,1.023) -- (-0.816,1.003) -- (-0.821,0.983) -- (-0.827,0.964) -- (-0.832,0.944) -- (-0.837,0.924) -- (-0.842,0.904) -- (-0.847,0.885) -- (-0.852,0.865) -- (-0.857,0.846) -- (-0.861,0.826) -- (-0.866,0.807) -- (-0.871,0.788) -- (-0.875,0.769) -- (-0.880,0.750) -- (-0.884,0.730) -- (-0.889,0.711) -- (-0.893,0.693) -- (-0.897,0.674) -- (-0.902,0.655) -- (-0.906,0.636) -- (-0.911,0.618) -- (-0.915,0.599) -- (-0.920,0.581) -- (-0.924,0.562) -- (-0.929,0.544) -- (-0.934,0.526) -- (-0.939,0.508) -- (-0.943,0.490) -- (-0.948,0.472) -- (-0.953,0.454) -- (-0.958,0.436) -- (-0.963,0.418) -- (-0.969,0.400) -- (-0.974,0.382) -- (-0.979,0.364) -- (-0.985,0.347) -- (-0.990,0.329) -- (-0.996,0.311) -- (-1.002,0.294) -- (-1.008,0.276) -- (-1.014,0.258) -- (-1.020,0.241) -- (-1.027,0.223) -- (-1.033,0.206);
\draw[fillback] (-1.033,0.206) -- (-1.040,0.188) -- (-1.046,0.170) -- (-1.053,0.153) -- (-1.060,0.135) -- (-1.067,0.118) -- (-1.074,0.100) -- (-1.081,0.082) -- (-1.089,0.065) -- (-1.096,0.047) -- (-1.104,0.029) -- (-1.111,0.012) -- (-1.119,-0.006) -- (-1.127,-0.024) -- (-1.135,-0.042) -- (-1.143,-0.060) -- (-1.151,-0.078) -- (-1.160,-0.096) -- (-1.168,-0.114) -- (-1.176,-0.133) -- (-1.185,-0.151) -- (-1.193,-0.169) -- (-1.202,-0.188) -- (-1.210,-0.206) -- (-1.219,-0.225) -- (-1.228,-0.244) -- (-1.236,-0.263) -- (-1.245,-0.282) -- (-1.254,-0.301) -- (-1.262,-0.320) -- (-1.271,-0.339) -- (-1.280,-0.358) -- (-1.288,-0.378) -- (-1.297,-0.397) -- (-1.305,-0.417) -- (-1.314,-0.437) -- (-1.322,-0.457) -- (-1.330,-0.477) -- (-1.339,-0.497) -- (-1.347,-0.517) -- (-1.355,-0.537) -- (-1.363,-0.558) -- (-1.370,-0.578) -- (-1.378,-0.599) -- (-1.385,-0.620) -- (-1.393,-0.641) -- (-1.400,-0.661) -- (-1.407,-0.682) -- (-1.413,-0.704) -- (-1.420,-0.725) -- (-1.426,-0.746) -- (-1.432,-0.768) -- (-1.438,-0.789) -- (-1.443,-0.811) -- (-1.449,-0.832) -- (-1.454,-0.854) -- (-1.458,-0.876) -- (-1.463,-0.897) -- (-1.467,-0.919) -- (-1.471,-0.941) -- (-1.474,-0.963) -- (-1.477,-0.985) -- (-1.480,-1.007) -- (-1.482,-1.029) -- (-1.484,-1.051) -- (-1.486,-1.073) -- (-1.487,-1.095) -- (-1.488,-1.117) -- (-1.488,-1.139) -- (-1.488,-1.161) -- (-1.488,-1.183) -- (-1.487,-1.205) -- (-1.486,-1.227) -- (-1.484,-1.249) -- (-1.482,-1.271) -- (-1.479,-1.293) -- (-1.476,-1.314) -- (-1.473,-1.336) -- (-1.468,-1.357) -- (-1.464,-1.379) -- (-1.459,-1.400) -- (-1.453,-1.421) -- (-1.447,-1.442) -- (-1.441,-1.463) -- (-1.433,-1.483) -- (-1.426,-1.504) -- (-1.418,-1.524) -- (-1.409,-1.544) -- (-1.400,-1.564) -- (-1.390,-1.584) -- (-1.380,-1.603) -- (-1.369,-1.622) -- (-1.358,-1.641) -- (-1.346,-1.660) -- (-1.334,-1.678) -- (-1.321,-1.697) -- (-1.307,-1.715) -- (-1.293,-1.732) -- (-1.279,-1.750) -- (-1.264,-1.766) -- (-1.248,-1.783) -- (-1.232,-1.799) -- (-1.216,-1.815) -- (-1.198,-1.831) -- (-1.181,-1.846) -- (-1.163,-1.861) -- (-1.144,-1.876) -- (-1.125,-1.890) -- (-1.105,-1.904) -- (-1.085,-1.917) -- (-1.065,-1.930) -- (-1.044,-1.942) -- (-1.022,-1.954) -- (-1.000,-1.966) -- (-0.978,-1.977) -- (-0.955,-1.987) -- (-0.932,-1.997) -- (-0.908,-2.007) -- (-0.884,-2.016) -- (-0.860,-2.025) -- (-0.835,-2.033);
\draw[fill] (-0.835,-2.033) -- (-0.809,-2.041) -- (-0.784,-2.048) -- (-0.758,-2.054) -- (-0.731,-2.061) -- (-0.705,-2.066) -- (-0.678,-2.071) -- (-0.651,-2.076) -- (-0.623,-2.080) -- (-0.595,-2.083) -- (-0.567,-2.086) -- (-0.539,-2.088) -- (-0.510,-2.090) -- (-0.481,-2.091) -- (-0.452,-2.092) -- (-0.423,-2.092) -- (-0.393,-2.091) -- (-0.364,-2.090) -- (-0.334,-2.088) -- (-0.304,-2.086) -- (-0.274,-2.083) -- (-0.244,-2.080) -- (-0.213,-2.076) -- (-0.183,-2.071) -- (-0.153,-2.066) -- (-0.122,-2.061) -- (-0.092,-2.054) -- (-0.061,-2.047) -- (-0.031,-2.040) -- (-0.000,-2.032) -- (0.030,-2.024) -- (0.061,-2.015) -- (0.091,-2.005) -- (0.122,-1.995) -- (0.152,-1.985) -- (0.182,-1.973) -- (0.212,-1.962) -- (0.242,-1.950) -- (0.272,-1.937) -- (0.301,-1.924) -- (0.331,-1.910) -- (0.360,-1.896) -- (0.389,-1.881) -- (0.418,-1.866) -- (0.446,-1.850) -- (0.475,-1.834) -- (0.503,-1.818) -- (0.530,-1.801) -- (0.558,-1.783) -- (0.585,-1.766) -- (0.612,-1.747) -- (0.639,-1.729) -- (0.665,-1.710) -- (0.691,-1.691) -- (0.716,-1.671) -- (0.741,-1.651) -- (0.766,-1.630) -- (0.791,-1.610) -- (0.815,-1.589) -- (0.838,-1.567) -- (0.861,-1.546) -- (0.884,-1.524) -- (0.906,-1.502) -- (0.928,-1.479) -- (0.950,-1.456) -- (0.971,-1.433) -- (0.991,-1.410) -- (1.011,-1.387) -- (1.031,-1.363) -- (1.050,-1.340) -- (1.068,-1.316) -- (1.086,-1.292) -- (1.104,-1.268) -- (1.121,-1.243) -- (1.138,-1.219) -- (1.154,-1.195) -- (1.169,-1.170) -- (1.184,-1.145) -- (1.199,-1.121) -- (1.213,-1.096) -- (1.226,-1.071) -- (1.239,-1.047) -- (1.252,-1.022) -- (1.264,-0.997) -- (1.275,-0.972) -- (1.286,-0.948) -- (1.297,-0.923) -- (1.307,-0.899) -- (1.316,-0.874) -- (1.325,-0.850) -- (1.333,-0.825) -- (1.341,-0.801) -- (1.348,-0.777) -- (1.355,-0.753) -- (1.362,-0.730) -- (1.368,-0.706) -- (1.373,-0.682) -- (1.378,-0.659) -- (1.382,-0.636) -- (1.386,-0.613) -- (1.390,-0.591) -- (1.393,-0.568) -- (1.396,-0.546) -- (1.398,-0.524) -- (1.400,-0.502) -- (1.401,-0.481) -- (1.402,-0.460) -- (1.403,-0.439) -- (1.403,-0.418) -- (1.403,-0.398) -- (1.402,-0.378) -- (1.401,-0.358) -- (1.400,-0.338) -- (1.398,-0.319) -- (1.396,-0.300) -- (1.394,-0.282) -- (1.391,-0.264) -- (1.389,-0.246) -- (1.385,-0.228) -- (1.382,-0.211) -- (1.378,-0.194) -- (1.374,-0.178) -- (1.370,-0.161) -- (1.365,-0.146) -- (1.360,-0.130) -- (1.355,-0.115) -- (1.350,-0.100) -- (1.345,-0.086) -- (1.339,-0.072) -- (1.333,-0.058) -- (1.327,-0.044) -- (1.321,-0.031) -- (1.315,-0.019) -- (1.309,-0.006) -- (1.302,0.006) -- (1.296,0.017) -- (1.289,0.029) -- (1.282,0.040) -- (1.275,0.050) -- (1.268,0.061) -- (1.261,0.071) -- (1.254,0.081) -- (1.247,0.090) -- (1.240,0.099) -- (1.233,0.108) -- (1.226,0.116) -- (1.218,0.125) -- (1.211,0.132) -- (1.204,0.140) -- (1.197,0.148) -- (1.190,0.155) -- (1.183,0.162) -- (1.176,0.168) -- (1.169,0.175) -- (1.162,0.181);

  \end{scope}
  \node[lab, anchor=north] at (9.65,-4.95) {never closes, never a knot};
  \begin{axis}[snd, at={(12.85cm,-4.95cm)}, xmin=0, xmax=5.8, ymax=1.75, xtick={0,1,...,5}, xlabel={frequency ÷ f₁ \quad (f₂ = 1.618 f₁)}]
    \grid
    \addplot[ycomb, line width=1.4pt, ink] table[x=freq, y=amp] {tkco-open.dat};
    \node[lab, anchor=north east, align=right] at (axis cs:5.8,1.78) {most lines fall between};
  \end{axis}
\end{scope}
\node[font=\scriptsize, anchor=north west, text=ink2, text width=19.6cm] at (0,-12.4) {A path that winds p times one way and q times the other is a (p, q) torus knot, and the winding numbers do not change when the speed does. Whether that invariance is audible is an open listening test. Lines are m·f₁ + n·f₂ for |m| + |n| ≤ 5, heights falling with |m| + |n|.};

% ---------- title block ----------
\draw[ink, line width=0.8pt] (-0.3,-13.6) rectangle (19.9,-14.5);
\draw[ink, line width=0.3pt] (9.6,-13.6) -- (9.6,-14.5) (14.6,-13.6) -- (14.6,-14.5);
\node[font=\scriptsize\bfseries, text=ink, anchor=west] at (-0.1,-14.05) {DSP IN LOOPING DIMENSIONS · TORUS KNOT, CLOSED AND OPEN};
\node[font=\scriptsize, text=ink, anchor=west] at (9.8,-14.05) {INK computed};
\node[font=\scriptsize, text=ink2, anchor=west] at (14.8,-14.05) {FIG · 2026-11-09 · PAPER};

\begin{scope}[on background layer]
  \fill[paper] ($(current bounding box.south west)+(-0.4,-0.4)$) rectangle ($(current bounding box.north east)+(0.4,0.4)$);
\end{scope}
\end{tikzpicture}
\end{document}
